\documentclass[11pt]{article}
\usepackage{amsfonts,mathrsfs,amssymb,amsthm,mathptm}
\usepackage[hidelinks]{hyperref}
%\usepackage{hyperref}
\usepackage{amsmath,amscd}
\usepackage{pslatex}
\usepackage{color}
\usepackage[all]{xy}
\usepackage{fancyhdr}
\usepackage{tikz-cd}
\usepackage[title]{appendix}
\usepackage{comment}
\usepackage{needspace}

\oddsidemargin -10pt \voffset=-2.5cm \addtolength{\textheight}{4cm}
\addtolength{\textwidth}{3.5cm}
\setlength{\emergencystretch}{3em}

\begin{document}
%\sloppy
\newtheorem{Theo}{Theorem}[section]
\newtheorem{Prop}[Theo]{Proposition}
\newtheorem{Lem}[Theo]{Lemma}
\newtheorem{Koro}[Theo]{Corollary}
\theoremstyle{definition}
\newtheorem{Rem}[Theo]{Remark}

\renewcommand{\labelenumi}{\textup{(\arabic{enumi})}}
\renewcommand{\theenumi}{\arabic{enumi}}

\newcommand{\add}{{\rm add}}
\newcommand{\con}{{\rm con}}
\newcommand{\gd}{{\rm gldim}}
\newcommand{\sd}{{\rm stdim}}
\newcommand{\sr}{{\rm sr}}
\newcommand{\dm}{{\rm domdim}}
\newcommand{\cdm}{{\rm codomdim}}
\newcommand{\tdim}{{\rm dim}}
\newcommand{\E}{{\rm E}}
\newcommand{\Mor}{{\rm Morph}}
\newcommand{\End}{{\rm End}}
\newcommand{\ind}{{\rm ind}}
\newcommand{\rsd}{{\rm resdim}}
\newcommand{\rd} {{\rm rd}}
\newcommand{\ol}{\overline}
\newcommand{\overpr}{$\hfill\square$}
\newcommand{\rad}{{\rm rad}}
\newcommand{\soc}{{\rm soc}}
\renewcommand{\top}{{\rm top}}
\newcommand{\pd}{{\rm pdim}}
\newcommand{\id}{{\rm idim}}
\newcommand{\fld}{{\rm fdim}}
\newcommand{\Fac}{{\rm Fac}}
\newcommand{\Gen}{{\rm Gen}}
\newcommand{\fd} {{\rm findim}}
\newcommand{\Fd} {{\rm Findim}}
\newcommand{\Pf}[1]{{\mathscr P}^{<\infty}(#1)}
\newcommand{\DTr}{{\rm DTr}}
\newcommand{\cpx}[1]{#1^{\bullet}}
\newcommand{\D}[1]{{\mathscr D}(#1)}
\newcommand{\Dz}[1]{{\mathscr D}^+(#1)}
\newcommand{\Df}[1]{{\mathscr D}^-(#1)}
\newcommand{\Db}[1]{{\mathscr D}^b(#1)}
\newcommand{\C}[1]{{\mathscr C}(#1)}
\newcommand{\Cz}[1]{{\mathscr C}^+(#1)}
\newcommand{\Cf}[1]{{\mathscr C}^-(#1)}
\newcommand{\Cb}[1]{{\mathscr C}^b(#1)}
\newcommand{\Dc}[1]{{\mathscr D}^c(#1)}
\newcommand{\K}[1]{{\mathscr K}(#1)}
\newcommand{\Kz}[1]{{\mathscr K}^+(#1)}
\newcommand{\Kf}[1]{{\mathscr  K}^-(#1)}
\newcommand{\Kb}[1]{{\mathscr K}^b(#1)}
\newcommand{\DF}[1]{{\mathscr D}_F(#1)}
\newcommand{\Kac}[1]{{\mathscr K}_{\rm ac}(#1)}
\newcommand{\Keac}[1]{{\mathscr K}_{\mbox{\rm e-ac}}(#1)}
\newcommand{\modcat}{\ensuremath{\mbox{{\rm -mod}}}}
\newcommand{\Modcat}{\ensuremath{\mbox{{\rm -Mod}}}}
\newcommand{\Spec}{{\rm Spec}}
\newcommand{\stmc}[1]{#1\mbox{{\rm -{\underline{mod}}}}}
\newcommand{\Stmc}[1]{#1\mbox{{\rm -{\underline{Mod}}}}}
\newcommand{\prj}[1]{#1\mbox{{\rm -proj}}}
\newcommand{\inj}[1]{#1\mbox{{\rm -inj}}}
\newcommand{\Prj}[1]{#1\mbox{{\rm -Proj}}}
\newcommand{\Inj}[1]{#1\mbox{{\rm -Inj}}}
\newcommand{\PI}[1]{#1\mbox{{\rm -Prinj}}}
\newcommand{\GP}[1]{#1\mbox{{\rm -GProj}}}
\newcommand{\GI}[1]{#1\mbox{{\rm -GInj}}}
\newcommand{\gp}[1]{#1\mbox{{\rm -Gproj}}}
\newcommand{\gi}[1]{#1\mbox{{\rm -Ginj}}}
\newcommand{\opp}{^{\rm op}}
\newcommand{\otimesL}{\otimes^{\rm\mathbb L}}
\newcommand{\rHom}{{\rm\mathbb R}{\rm Hom}\,}
\newcommand{\pdim}{\pd}
\newcommand{\Hom}{{\rm Hom}}
\newcommand{\Coker}{{\rm Coker}}
\newcommand{ \Ker  }{{\rm Ker}}
\newcommand{ \Cone }{{\rm Con}}
\newcommand{ \Img  }{{\rm Im}}
\newcommand{\Ext}{{\rm Ext}}
\newcommand{\StHom}{{\rm \underline{Hom}}}
\newcommand{\StEnd}{{\rm \underline{End}}}
\newcommand{\KK}{I\!\!K}
\newcommand{\gm}{{\rm _{\Gamma_M}}}
\newcommand{\gmr}{{\rm _{\Gamma_M^R}}}
\def\vez{\varepsilon}\def\bz{\bigoplus}  \def\sz {\oplus}
\def\epa{\xrightarrow} \def\inja{\hookrightarrow}
\newcommand{\lra}{\longrightarrow}
\newcommand{\llra}{\longleftarrow}
\newcommand{\lraf}[1]{\stackrel{#1}{\lra}}
\newcommand{\llaf}[1]{\stackrel{#1}{\llra}}
\newcommand{\ra}{\rightarrow}
\newcommand{\dk}{{\rm dim_{_{k}}}}
\newcommand{\holim}{{\rm Holim}}
\newcommand{\hocolim}{{\rm Hocolim}}
\newcommand{\colim}{{\rm colim\, }}
\newcommand{\limt}{{\rm lim\, }}
\newcommand{\Add}{{\rm Add }}
\newcommand{\Prod}{{\rm Prod }}
\newcommand{\Tor}{{\rm Tor}}
\newcommand{\Cogen}{{\rm Cogen}}
\newcommand{\Tria}{{\rm Tria}}
\newcommand{\Loc}{{\rm Loc}}
\newcommand{\Coloc}{{\rm Coloc}}
\newcommand{\tria}{{\rm tria}}
\newcommand{\Con}{{\rm Con}}
\newcommand{\Thick}{{\rm Thick}}
\newcommand{\thick}{{\rm thick}}
\newcommand{\Sum}{{\rm Sum}}



{\Large \bf
\begin{center}
From Frobenius seeds to self-orthogonal modules
over symmetric algebras
\end{center}}

\medskip\centerline{\textbf{Hongxing Chen}}
\hypersetup{pdfauthor={Hongxing Chen}}

\renewcommand{\thefootnote}{\alph{footnote}}
\setcounter{footnote}{-1} \footnote{2020 Mathematics Subject
Classification: Primary 16G10, 16E35, 16D50; Secondary 18G80, 16E10.}
\renewcommand{\thefootnote}{\alph{footnote}}
\setcounter{footnote}{-1} \footnote{Keywords: Frobenius algebra; Self-orthogonal module;
Symmetric algebra; Tachikawa's second conjecture; Trivial extension.}

\begin{abstract}
We construct finite-dimensional symmetric algebras admitting nonprojective self-orthogonal modules from Frobenius algebras equipped with a module satisfying explicit stable and homological conditions. The resulting modules have complexity three, and their stable endomorphism algebras are isomorphic to the algebra of dual numbers over the ground field. The construction combines a variable input with a fixed auxiliary factor carrying the required twisting, and preserves the input through specified quotient and corner constructions. The proof combines derived induction, an obstruction to perfect factorization, and a cone construction over a Gorenstein triangular matrix algebra. For a symmetric algebra and a self-orthogonal generator, we describe the derived category of the opposite endomorphism algebra through a recollement with a two-term $2$-symmetric dg algebra. For our counterexamples, it is formal and quasi-isomorphic to an explicit four-dimensional dg algebra. Applications to cyclic coverings of a quantum exterior algebra yield connected examples with explicitly determined numbers of simple modules. The construction works over every field containing an element of infinite multiplicative order.
\end{abstract}


{\footnotesize\tableofcontents}
\clearpage

\section{Introduction}

We first review recent developments concerning Tachikawa's second
conjecture and present our main result. We then explain its relation to
earlier constructions and outline the proof strategy and the
organization of the paper.

\subsection{Background and main result}

Tachikawa's second conjecture (TC2) asserts that every finitely generated
self-orthogonal module over a finite-dimensional self-injective algebra
is projective. Here a module is called \emph{self-orthogonal} if all its
positive self-extension groups vanish. Enomoto and Marczinzik \cite{EM}
proved that, for Artin algebras over a fixed commutative artinian ring,
(TC2) is globally equivalent to
Tachikawa's first conjecture and the Auslander--Reiten, generalized
Nakayama, and Nakayama conjectures. Since the finitistic dimension
conjecture also implies (TC2), a counterexample to (TC2) disproves all
these conjectures (see also Remark~\ref{audit:section}).

Recent progress also includes algebra constructions and positive results
for specific classes. Chen, Fang and Xi \cite{CFX} introduced
mirror-reflective algebras and developed iterative constructions linking
nonprojective self-orthogonal modules over symmetric algebras to
nontrivial stratifying ideals in other symmetric algebras and to derived
recollements. Chen and Xi \cite{CXSO} developed a homological theory of
self-orthogonal generators using recollements of relative stable
categories and descriptions of compact objects, thereby establishing the
Nakayama conjecture for Gorenstein--Morita algebras. They also gave
projectivity criteria involving virtually Gorenstein endomorphism
algebras and the Nakayama functor \cite{CXVG}. Li \cite{LiVG} proved that virtually Gorenstein
Artin algebras are weakly Gorenstein and satisfy both the Nakayama
conjecture and Tachikawa's first conjecture. The conjecture (TC2) holds
for self-injective Igusa--Todorov algebras, including those with radical
cube zero \cite{ZZ}, and quantum complete intersections \cite{Xia}.

OpenAI \cite{O} constructed a finite-dimensional symmetric algebra
with a nonprojective self-orthogonal module over a rational function
field of characteristic two. Building on the ten-dimensional
starting algebra and character map used in that construction,
Tang \cite{F} subsequently extended the construction to every field
containing an element of infinite multiplicative order. He adjusted
the multiplication signs and used an involution to twist the right
action on the target dual bimodule. We work under the same field
hypothesis.

The main result reads as follows; see Theorem~\ref{general:construction}
for its full statement and Subsection~\ref{general:complexity-section}
for the definition of complexity.
\begin{quote}
\textbf{Main theorem.}
\emph{Let $k$ be a field containing an element of infinite
multiplicative order, $H$ a finite-dimensional Frobenius $k$-algebra, and $M_0$ a finite-dimensional indecomposable nonprojective left $H$-module satisfying}
\[
 \tau_HM_0\simeq M_0,\qquad \underline{\End}_H(M_0)\simeq k,\qquad
 \Ext_H^i(M_0,M_0)=0\quad(i\ge2),
\]
\emph{where $\tau_H$ denotes the Auslander--Reiten translation
for $H$-modules and $\underline{\End}_H(M_0)$ denotes the stable
endomorphism algebra of $M_0$.}
\emph{Set $C=\End_H(H\oplus M_0)^{\mathrm{op}}$ and
$S=C/CeC$, where $e\in C$ is the idempotent corresponding to the direct summand $H$. Then the construction
produces a finite-dimensional symmetric $k$-algebra $A$, a finite-dimensional
nonprojective self-orthogonal left $A$-module $M$ of complexity three, and a specified
ideal $I$ of $A$ such that}
\[
 A/I\simeq C,\qquad M/IM\simeq S,\qquad
 \underline{\End}_A(M)\simeq k[\varepsilon]/(\varepsilon^2).
\]
\emph{If $H$ is connected and split basic with $r$ simple modules,
then $A$ is connected and split basic with $4(r+1)$ simple modules.}
\end{quote}

Our construction retains the methods of derived induction, bimodule
lifting, triangular cones and transfer through trivial extensions
used in \cite{O,F}, while allowing the Frobenius algebra $H$ and
the module $M_0$ to vary. We take Tang's ten-dimensional seed as a
fixed auxiliary factor, expressed here as an endomorphism algebra.
Its involution acts only on this fixed factor, so no compatible
automorphism of the given Frobenius algebra is required.
The specified quotient and corner identifications in
Theorem~\ref{general:construction}\textup{(\ref{general:construction-recovery})} recover the input
algebra and module.

We also change the intermediate symmetric algebra: we take the
trivial extension after tensoring the seeds, whereas \cite{O,F}
use a tensor product of two trivial extensions. The distinguished
module then has a Yoneda algebra with one generator in degree five
(Proposition~\ref{ind:one-generator}), in place of two generators
in degree three. Its stable Hom groups are computed from a single
induction triangle, replacing the two-variable Koszul calculation
in \cite{O,F}. We construct bimodule maps that give nonzero
morphisms in the stable category after tensoring with the
distinguished module. The key step is to show that a twisted
character map does not factor through a perfect complex of
bimodules (Proposition~\ref{lift:trace-obstruction}). We prove this
using a fixed corner algebra isomorphic to $k[t]/(t^2)$.

We also compare the derived categories of a symmetric algebra $A$
and $\Gamma=\End_A(A\oplus M)^{\mathrm{op}}$ for a self-orthogonal
$A$-module $M$. For the idempotent $e\in\Gamma$ corresponding to
the direct summand $A$ of $A\oplus M$, the kernel of the functor $e(-)$ is triangle equivalent to the derived category of a two-term
$2$-symmetric dg algebra
(Proposition~\ref{dg:symmetric-application}). For the pairs
constructed here, this dg algebra is formal and admits an explicit
four-dimensional model (Corollary~\ref{dg:four-dimensional}).

The complexity assertion is proved in Proposition~\ref{general:complexity}
and shows that (TC2) fails even for modules of finite complexity. For comparison,
every finite-dimensional self-orthogonal module of complexity at most
one over a finite-dimensional self-injective algebra is projective
(Remark~\ref{general:bounded-selforthogonal}).

The field hypothesis and the complexity calculation lead to the
following questions.
\begin{enumerate}
\item
Does (TC2) hold for finite-dimensional self-injective algebras over
finite fields?
\item
Let $R$ be a finite-dimensional self-injective algebra over a field.
Is every finite-dimensional self-orthogonal $R$-module of complexity
at most two projective?
\end{enumerate}

Appendix~\ref{family:section} uses cyclic coverings of a quantum exterior algebra to construct Frobenius algebras and modules satisfying the hypotheses of the main theorem for $n\ge2$. The original algebra and module in Subsection~\ref{seed:concrete-pair} give the case $n=1$. Applying the main theorem yields, over each such field, connected symmetric algebras with exactly $4(n+1)$ simple modules for every $n\ge1$.

We provide detailed proofs of the construction and its key homological steps.
These arguments can be checked without consulting \cite{O,F}.

\subsection{Proof strategy and organization}

We explain the main steps in the proof. For the self-orthogonality
assertion, Theorem~\ref{frame:transfer} reduces the task to constructing
a finite-dimensional algebra $\Lambda$ and a
nonprojective $\Lambda$-module $Z$ with a complete resolution whose
self-Hom cohomology vanishes outside degrees $-1$ and $0$. The cone
criterion of Theorem~\ref{frame:cone} provides a way to construct such a
pair. Proposition~\ref{frame:factory} expresses sufficient conditions
for this criterion in terms of stable morphism spaces and two bimodule
maps with the same nonzero stable evaluation. These are general results
established in Section~\ref{prelim:section}.

We first construct the symmetric algebra over which these stable
morphisms will be computed. Proposition~\ref{general:seed} associates
to $(H,M_0)$ the algebra $C$ and a self-orthogonal simple module $S$.
Fix $q\in k^\times$ of infinite multiplicative order. We combine
$(C,S)$ with the corresponding fixed signed seed $(C_*,S_*,\gamma_*)$
of Notation~\ref{seed:auxiliary-notation}.
Set
\[
 B=C_*\otimes_k C,\qquad X=S_*\otimes_k S,\qquad E=B\ltimes DB,
\]
where $D=\Hom_k(-,k)$. We regard $X$ as an $E$-module via the
canonical projection $E\twoheadrightarrow B$.
By Lemma~\ref{ind:tensor-input}, the image of $X$ under the derived
Nakayama functor $DB\otimesL_B(-)$ is isomorphic to $X[4]$.
Thus the short exact sequence
$0\to DB\to E\to B\to0$ yields the induction triangle
in the bounded derived category of $E$-modules
\[
 X[4]\longrightarrow E\otimesL_BX
 \longrightarrow X\xrightarrow{\varpi}X[5].
\]
Derived adjunction then gives $\Ext_E^*(X,X)\simeq k[\varpi]$,
with $\varpi$ of degree five (Proposition~\ref{ind:one-generator}).
Stable duality determines the negative-degree groups as well
(Proposition~\ref{ind:Tate-and-middle}). These calculations supply the
stable self-Hom conditions required by Proposition~\ref{frame:factory}.

The next task is to realize the derived data by a finite-dimensional
bimodule, so that we can form an ordinary finite-dimensional algebra.
Proposition~\ref{ind:Y-evaluation} constructs a finite-dimensional
$E$-bimodule $Y$ from $E\otimesL_BE$. It is projective as both a left
and a right $E$-module. Its evaluation at $X$ satisfies the remaining
support condition: the spaces
$\underline{\Hom}_E(X,(Y\otimes_E X)[i])$ are one-dimensional for
$i=-3,0$ and zero otherwise.

We next construct two bimodule maps into $Y$ with the same nonzero
stable evaluation at $X$. The twisted regular bimodules
$U_1,U_q$ have two complementary roles. Their tensor functors
$U_1\otimes_E(-)$ and $U_q\otimes_E(-)$ induce automorphisms of
$\underline{\Hom}_E(X,X[i])$ whose difference is invertible for
$i\ne0$; Proposition~\ref{ind:weights} proves this using the infinite
multiplicative order of $q$. To construct ordinary bimodule maps from
these sources to $Y$, we first establish the perfect-factorization
obstruction of Proposition~\ref{lift:trace-obstruction}.
Using this obstruction and stable duality, we construct $E$-bimodule
maps $g_1:U_1\to Y$ and $g_q:U_q\to Y$.
We then use the stable tensor--Hom adjunction of
Lemma~\ref{lift:stable-adjunction} to prove that tensoring these maps
with $X$ gives nonzero morphisms in the stable category
(Proposition~\ref{lift:nonzero-lifts}). Since the degree-zero evaluation
space is one-dimensional, we can rescale them so that
$g_1\otimes_E1_X$ and $g_q\otimes_E1_X$ give the same nonzero
stable morphism $X\to Y\otimes_E X$, under the source identifications
$U_1\otimes_E X\simeq X\simeq U_q\otimes_E X$.

Choose an $E^e$-projective cover $\pi_Q:Q\twoheadrightarrow Y$,
and define $F$ by
\[
 0\longrightarrow F\longrightarrow U_1\oplus U_q\oplus Q
 \xrightarrow{(g_1,-g_q,\pi_Q)}Y\longrightarrow0.
\]
This sequence splits on each $E$-side. All hypotheses of
Proposition~\ref{frame:factory} are now
satisfied. It constructs the stable morphism $v:X\to F\otimes_E X$
required by the cone criterion. For the triangular matrix algebra
$\Lambda=\left(\begin{smallmatrix}E&0\\F&E\end{smallmatrix}\right)$,
Theorem~\ref{frame:cone} therefore produces $Z$ and its complete
resolution with the required self-Hom vanishing. Applying
Theorem~\ref{frame:transfer} to $A=\Lambda\ltimes D\Lambda$ and
$M=A\otimes_\Lambda Z$ gives the desired nonprojective self-orthogonal
module over a symmetric algebra. The recovery relations and the
assertions about connectedness and simple modules are verified in
the final construction in Theorem~\ref{general:construction}.

Section~\ref{prelim:section} develops the homological tools and abstract
construction criteria. Section~\ref{seed:section} presents the general
Frobenius seeds and the fixed signed auxiliary seed.
Section~\ref{ind:section} treats derived induction and the
finite-dimensional bimodule model, while
Section~\ref{lift:construction-section} constructs the twists and the
required bimodule maps. Section~\ref{app:section} gives the final
construction and discusses its consequences for related homological
conjectures. Subsection~\ref{general:complexity-section} computes the
complexity of the constructed module. Subsection~\ref{general:dg-models}
constructs two-term dg models, establishes the corresponding derived
recollements, and proves formality for the constructed examples.
Appendix~\ref{family:section} provides the cyclic-covering examples.

\section{Homological tools and abstract construction} \label{prelim:section}

We first fix the homological conventions. The main results are the
transfer theorem (Theorem~\ref{frame:transfer}) and the cone criterion
(Theorem~\ref{frame:cone}).

\subsection{Conventions and homological preliminaries}\label{prelim:notation}


Throughout, $k$ is a field, and tensor products without a subscript
are taken over $k$. Ordinary algebras are associative unital finite-dimensional
$k$-algebras. Modules and bimodules, including graded modules, are
finite-dimensional over $k$ unless otherwise stated, and modules
are unital left modules unless a side is specified. Graded
Yoneda algebras need not have finite total dimension. For an $R$-bimodule $L$,
we write $ {}_RL$ and $L_R$ for its left and right restrictions.
Composition is written from right to left.

The opposite and enveloping algebras are denoted by $R^{\rm op}$ and
$R^e=R\otimes_kR^{\rm op}$, respectively. An $R$-bimodule $Z$ is regarded
as a left $R^e$-module by $(a\otimes b^{\rm op})z=azb$ for $a,b\in R$ and $z\in Z$.
It is \emph{separately projective} if its two restrictions are projective; this condition does
not require projectivity over $R^e$.

We use the duality $D=\Hom_k(-,k)$. The dual of an $R$-bimodule $Z$ has
the action $(a\lambda b)(z)=\lambda(bza)$, where $\lambda\in DZ$. Thus duality interchanges the two sides. We denote the category of left modules by $R\text{-}\mathrm{mod}$, and its full
subcategories of projective and injective modules by
$R\text{-}\mathrm{proj}$ and $R\text{-}\mathrm{inj}$. The notation
$\add(L)$ denotes the direct summands of finite direct sums of $L$.
The Jacobson radical of $R$ is denoted by $\rad(R)$.
An algebra $R$ is \emph{connected} if it is not isomorphic to a direct
product of two nonzero algebras.
It is \emph{split basic} if $R/\rad(R)\simeq k^m$ for some
$m\ge1$.  A \emph{character} is a unital $k$-algebra homomorphism $\chi:R\to k$; the corresponding
one-dimensional module is denoted by $k_\chi$.

Projective and injective dimensions are denoted by $\pd_R(U)$ and
$\id_R(U)$ on the left, and by a subscript $R^{\rm op}$ on the right.
For $U\ne0$, these are the least lengths of a projective resolution
and an injective coresolution, respectively, with value $\infty$
when no finite resolution or coresolution exists.
The ordinary extension groups $\Ext_R^a(U,V)$ for $a\in\mathbb{Z}$ are computed using a projective resolution of $U$ or an injective coresolution of $V$; we set them equal to zero for $a<0$.


An algebra $R$ is \emph{self-injective} if the regular module ${_R}R$ (or equivalently, $R_R$)
is injective. It is \emph{Frobenius} if $ {}_RR\simeq {}_RDR$,
equivalently, if there is a linear form $\lambda:R\to k$ (called a \emph{Frobenius functional}) such that
$(a,b)\mapsto\lambda(ab)$ is nondegenerate. It is \emph{symmetric}
if $R\simeq DR$ as $R$-bimodules. Equivalently, there is a linear form
$t:R\to k$ satisfying $t(ab)=t(ba)$ for which
$(a,b)\mapsto t(ab)$ is nondegenerate. Indeed, such a form defines
a bimodule isomorphism
$$R\lra DR, \;\; a\mapsto[b\mapsto t(ba)],\quad a, b\in R.$$
Conversely, a bimodule isomorphism $R\to DR$ determines $t$ by evaluation at $1$.


For a minimal injective coresolution $0\to {}_RR\to I^0\to I^1\to I^2\to\cdots$, we define the \emph{dominant dimension} of $R$ to be
\[
\operatorname{domdim}(R)=\sup\{d\in\mathbb{Z}_{\ge0}\mid I^j\text{ is projective for every }0\le j<d\}.
\]
This value is infinite if every term $I^j$ is projective.
The \emph{Nakayama conjecture} predicts that
if $\operatorname{domdim}(R)=\infty$, then $R$ is self-injective.

\paragraph{Trivial extensions.}
For an algebra $R$, its trivial extension $R\ltimes DR$ is
the vector space $R\oplus DR$ with multiplication
$(r,\lambda)(s,\mu)=(rs,r\mu+\lambda s)$ and unit $(1,0)$, where $r,s \in R$ and $\lambda,\mu\in DR$.
It is symmetric, with symmetrizing form $t(r,\lambda)=\lambda(1)$,
and has dimension $2\dim_k(R)$. The ideal $DR=0\oplus DR$ is
square-zero, and the standard inclusion $R\hookrightarrow R\ltimes DR$
and projection $R\ltimes DR\twoheadrightarrow R$ are algebra homomorphisms.

\paragraph{Restriction and twists.}
For an algebra homomorphism $\varphi:R\to T$, restriction is the
functor
\[
 \operatorname{Res}_\varphi:T\text{-}\mathrm{mod}\longrightarrow R\text{-}\mathrm{mod},\quad U\mapsto {_\varphi}U:\;\;r\cdot u=\varphi(r)u, \quad r\in R, u\in U.
\]
It leaves the underlying vector spaces and maps unchanged and
extends to complexes termwise, with the same differentials.
For a $k$-algebra automorphism $\sigma$ of $R$, the left twist
$ {}_\sigma U=\operatorname{Res}_\sigma(U)$ has action
$r\cdot_\sigma u=\sigma(r)u$.
For a map $f:U\to V$, the notation
$ {}_\sigma f:{}_\sigma U\to{}_\sigma V$ denotes the same
underlying map. Likewise, $ {}_\sigma P^\bullet$ is obtained by
twisting each term of $P^\bullet$ without changing its differential.
If $\sigma^2=1$, the double twist is identified with the original
module or complex.

For bimodules, $R_\sigma$ denotes the right twist, with its usual
left action and right action $a\cdot b=a\sigma(b)$ for $a,b\in R$.
More generally, $L_\sigma$ and $ {}_\sigma L$ twist only the right
and left actions, respectively. There is a natural isomorphism
\[
 R_\sigma\otimes_RU\xrightarrow{\sim}
 {}_{\sigma^{-1}}U=\operatorname{Res}_{\sigma^{-1}}(U),
 \qquad r\otimes u\longmapsto\sigma^{-1}(r)u,\quad r\in R, u\in U.
\]

\paragraph{Complexes and derived objects.}
Complexes over an algebra $R$ are cohomologically indexed and are written as
$P^\bullet=(P^j, d_P^j)_{j\in\mathbb{Z}}$. Projective models and complete resolutions may be unbounded, but
their terms are finite-dimensional modules. Full Hom complexes
use products and may have infinite-dimensional terms.
For $j,a\in\mathbb Z$, our conventions are
\begin{equation*}
 \begin{gathered}
 d_P^j:P^j\longrightarrow P^{j+1},\qquad
 c^j(P^\bullet)=\operatorname{coker}(d_P^{j-1}:P^{j-1}\to P^j),\\
 (P^\bullet[a])^j=P^{j+a},\qquad d_{P[a]}^j=(-1)^a d_P^{j+a}.
 \end{gathered}
\end{equation*}
Thus $c^j(P^\bullet)$ is a cokernel module, whereas
$H^j(P^\bullet)=\ker(d_P^j)/\operatorname{im}(d_P^{j-1})$ is the cohomology.
A complex is \emph{bounded above} if its terms vanish in sufficiently
large degrees, and has \emph{bounded} cohomology if its cohomology
vanishes outside a finite interval. Projective resolutions are
placed in degrees at most zero, and modules are concentrated in
degree zero. We also write $P_n=P^{-n}$ and
$\partial_n=d_P^{-n}:P_n\to P_{n-1}$. If $H^j(P^\bullet)=0$, the
differential induces an embedding
$c^j(P^\bullet)\hookrightarrow P^{j+1}$ with image $\operatorname{im}(d_P^j)$.

For arbitrary complexes, we use the (product) Hom complex
$\Hom_R^\bullet(P^\bullet,Q^\bullet)$, defined by
\begin{equation*}
 \Hom_R^a(P^\bullet,Q^\bullet)
 =\prod_{j\in\mathbb Z}\Hom_R(P^j,Q^{j+a}),\qquad
 (\partial f)^j=d_Q^{j+a}f^j-(-1)^a f^{j+1}d_P^j,
\end{equation*}
where $f^j:P^j\to Q^{j+a}$. Degree-zero cycles are cochain maps,
and degree-zero boundaries are nullhomotopic maps. Thus a homotopy
between cochain maps $f,g$ consists of maps $h^j:P^j\to Q^{j-1}$
such that $f^j-g^j=d_Q^{j-1}h^j+h^{j+1}d_P^j$.
For homogeneous maps $f$ of degree $a$ and $g$ of degree $b$,
graded composition is $(gf)^j=g^{j+a}f^j$.
These conventions apply also to unbounded complete resolutions.
For a complex $T^\bullet$ of $k$-vector spaces, the dual complex is
$DT^\bullet=\Hom_k^\bullet(T^\bullet,k)$, with terms
$(DT^\bullet)^n=D(T^{-n})$ and differential
$d_{DT}^n=(-1)^{n+1}D(d_T^{-n-1})$. The signed cone complex $\operatorname{Cone}(f)^\bullet$ of a
cochain map $f:P^\bullet\to Q^\bullet$ is
\begin{equation}\label{prelim:signed-cone}
 \operatorname{Cone}(f)^j=Q^j\oplus P^{j+1},\qquad
 d_{\operatorname{Cone}(f)}^j=
 \begin{pmatrix}d_Q^j&f^{j+1}\\0&-d_P^{j+1}\end{pmatrix}.
\end{equation}
It gives the triangle
$P^\bullet\xrightarrow{f}Q^\bullet\to\operatorname{Cone}(f)^\bullet\to P^\bullet[1]$;
the last arrow after one rotation is $-f[1]$.

For an additive category $\mathcal A$, we write $\mathscr K(\mathcal A)$
for the homotopy category of complexes in $\mathcal A$, and
$\mathscr K^b(\mathcal A)$ for its bounded subcategory.
We write $\mathscr D^b(R)=\mathscr D^b(R\text{-}\mathrm{mod})$ for
the bounded derived category of modules. The canonical functor
\[
 \mathscr K^b(R\text{-}\mathrm{proj})\longrightarrow\mathscr D^b(R)
\]
is fully faithful. Its essential image consists precisely of the objects
isomorphic in $\mathscr D^b(R)$ to bounded complexes of projective
$R$-modules. These are the \emph{perfect objects} of $\mathscr D^b(R)$,
and we identify $\mathscr K^b(R\text{-}\mathrm{proj})$ with this full subcategory. A bounded-above complex of
projective modules with bounded cohomology is a projective model
for an object of $\mathscr D^b(R)$. Such a model computes derived
morphisms and quasi-isomorphisms between projective models are homotopy
equivalences.  Derived tensor products are denoted by
$L\otimes_R^{\mathbf L}V$ for $L\in \mathscr D^b(R\opp)$ and $V\in\mathscr D^b(R)$, and
computed using suitable projective resolutions.

For a triangulated category $\mathscr T$ with suspension $[1]$, we write
$\Ext_{\mathscr T}^a(X,Y)=\Hom_{\mathscr T}(X,Y[a])$ for $X, Y\in\mathscr{T}$ and $a\in\mathbb Z$.
For $f\in\Ext_{\mathscr T}^a(X,Y)$ and
$g\in\Ext_{\mathscr T}^b(Y,Z)$, their product is
$gf:=g[a]\circ f\in\Ext_{\mathscr T}^{a+b}(X,Z)$.
For modules, $\Ext_{\mathscr D^b(R)}^a(U,V)\simeq\Ext_R^a(U,V)$
for all $a\in\mathbb Z$.


\subsection{Stable categories and derived models}\label{prelim:stable}

For $R$-modules $U,V$, let $\mathcal P_R(U,V)\subseteq\Hom_R(U,V)$ be the
subspace of maps factoring through projective modules. We set
$\underline{\Hom}_R(U,V)=\Hom_R(U,V)/\mathcal P_R(U,V)$ and
$\underline{\End}_R(U)=\underline{\Hom}_R(U,U)$, and write
$\underline f$ for the class of an ordinary map $f$ in $\underline{\Hom}_R(U,V)$.
The stable category $R\text{-}\underline{\mathrm{mod}}$ has the same
objects as $R\text{-}\mathrm{mod}$ and these quotient morphism spaces.
When $R$ is self-injective, it is triangulated, with inverse
suspension $[-1]=\Omega$ and suspension $[1]=\Omega^{-1}$.
These are represented by the kernel of a projective
cover and the cokernel of an injective envelope, respectively.

\paragraph{Nakayama functors and Auslander--Reiten translation.}
\label{prelim:AR-duality}

For an algebra $R$, the Nakayama functor
$\nu_R=D\Hom_R(-,R)\simeq DR\otimes_R(-)$ is right exact and
restricts to an equivalence from projective to injective modules.
For a projective module $Q$ and a module $V$, there is a natural isomorphism
$D\Hom_R(Q,V)\simeq\Hom_R(V,\nu_RQ)$;
see \cite[Section~6.4 and Lemma~6.4.5]{KrauseHTR}.

To define the Auslander--Reiten translation, take a minimal projective
presentation $P_1\xrightarrow{d_1}P_0\to U\to0$ and define the transpose by
\[
 \operatorname{Tr}_R(U)=\operatorname{coker}\bigl(\Hom_R(P_0,R)
       \xrightarrow{\Hom_R(d_1,R)}\Hom_R(P_1,R)\bigr).
\]
The transpose is a right $R$-module, and its dual
$\tau_RU=D\operatorname{Tr}_R(U)$ is a left $R$-module.

If $R$ is self-injective, $\nu_R$ is an autoequivalence and
$\tau_RU\simeq\Omega_R^2\nu_RU$ as modules. On
$R\text{-}\underline{\mathrm{mod}}$, this gives a natural isomorphism
of triangle functors $\tau_R\simeq\Omega_R^2\nu_R$;
see \cite[Proposition~IV.3.7(a)]{ARS}.
In the same setting, the \emph{Auslander--Reiten formula} gives natural isomorphisms
\[
 D\Ext_R^1(U,V)\simeq\underline{\Hom}_R(V,\tau_RU)
 \qquad(U,V\in R\text{-}\mathrm{mod});
\]
see \cite[Proposition~6.3.5]{KrauseHTR}.

For a Frobenius algebra, we choose a Frobenius functional $\lambda:R\to k$ and define its Nakayama automorphism $\eta$ by $\lambda(ab)=\lambda(b\eta(a))$ for $a, b\in R$. Then $DR\simeq R_\eta$ as bimodules
and $\nu_R\simeq\operatorname{Res}_{\eta^{-1}}$. In particular, $\nu_R$ is naturally isomorphic to the identity when $R$ is symmetric.

\phantomsection
For a self-injective algebra $R$, we use the following convention for
connecting morphisms:

Given a short exact sequence of $R$-modules
$0\to U\xrightarrow{f}V\xrightarrow{g}W\to0$, choose an embedding
$i:U\to J$ into a projective (also injective) module.
For the map $(f,-i):U\to V\oplus J$, $u\mapsto(f(u),-i(u))$, we use the quotient model
\[
 \operatorname{coker}(f,-i)
 =(V\oplus J)/\{(f(u),-i(u))\mid u\in U\}.
\]
The map $\operatorname{coker}(f,-i)\to W$, $[v,j]\mapsto g(v)$,
is a stable isomorphism. The map
$\operatorname{coker}(f,-i)\to J/i(U)=U[1]$,
$[v,j]\mapsto j+i(U)$, transported along this stable isomorphism,
gives the connecting morphism.
In particular, the connecting morphism of
$0\to U\to J\to U[1]\to0$ is $-1_{U[1]}$. We use this convention throughout.

A complex $P^\bullet$ of projective $R$-modules is \emph{totally acyclic}
if both $P^\bullet$ and $\Hom_R^\bullet(P^\bullet,R)$ are exact. It is called a
\emph{complete resolution} of $U$ if, in addition,
$c^0(P^\bullet)\simeq U$. A module admitting a complete resolution
is called \emph{Gorenstein-projective}. Every module over a self-injective algebra
has a complete resolution: splicing a projective resolution and an injective coresolution gives an exact complex of projective modules.

The next lemma shows that, for every $R$-module $V$, the groups
$H^a\Hom_R^\bullet(P^\bullet,V)$ for $a\in\mathbb Z$ are independent
of the chosen complete resolution $P^\bullet$ of a fixed module $U$
and agree with $\Ext_R^a(U,V)$ for $a\ge1$.


\begin{Lem}\label{frame:comparison}
Let $P^\bullet$ be a totally acyclic complex over $R$, and let $Q^\bullet$ be an
exact complex of projective $R$-modules. For every $a\in\mathbb Z$, there are isomorphisms
\begin{equation*}
 H^a\Hom_R^\bullet(P^\bullet,Q^\bullet)
 \simeq\underline{\Hom}_R(c^{-a}(P^\bullet),c^0(Q^\bullet))
 \simeq H^a\Hom_R^\bullet(P^\bullet,c^0(Q^\bullet)).
\end{equation*}
More generally, for every $R$-module $V$ and every $a\in\mathbb Z$,
there is a natural isomorphism
\[
 H^a\Hom_R^\bullet(P^\bullet,V)
 \simeq\underline{\Hom}_R(c^{-a}(P^\bullet),V).
\]
For $a\ge1$, these groups are naturally isomorphic to
$\Ext_R^a(c^0(P^\bullet),V)$.
\end{Lem}
\begin{proof}
We prove the first isomorphism by dualizing and applying Krause's
comparison lemma for injective complexes. Since the duality $D$
interchanges projective left modules and injective
right modules, $DQ^\bullet$ and $DP^\bullet$ are exact
complexes of injective right modules.
For every injective right module $J$, termwise transposition
gives an isomorphism of complexes
$\Hom_{R^{\rm op}}^\bullet(J,DP^\bullet)\simeq\Hom_R^\bullet(P^\bullet,DJ).$
The complex on the right is exact because $DJ$ is projective and $P^\bullet$ is totally acyclic.
Since $DP^\bullet$ is exact and $J$ is injective, $\Hom_{R^{\rm op}}^\bullet(DP^\bullet,J)$ is exact.
Thus $DP^\bullet$ is totally acyclic in the injective sense used in \cite[Section~7]{KrauseSD}.
We apply \cite[Lemma~7.1]{KrauseSD} to the abelian category $R\opp\modcat$, with $X=DQ^\bullet$ and $Y=DP^\bullet$. Its hypotheses hold since $X$ is exact
and $Y$ is totally acyclic. Using the canonical isomorphisms $\ker(d_{DT^\bullet}^0)\simeq D(c^0(T^\bullet))$ for $T^\bullet=P^\bullet, Q^\bullet$, we have
\[
 \Hom_{\mathscr K(R\text{-}\mathrm{proj})}(P^\bullet,Q^\bullet)
 \simeq\Hom_{\mathscr K(R^{\rm op}\text{-}\mathrm{inj})}(DQ^\bullet,DP^\bullet)
 \simeq\underline{\Hom}_R(c^0(P^\bullet),c^0(Q^\bullet)).
\]
For the last isomorphism, $D$ identifies the factorizations through
injective right modules in Krause's comparison with the factorizations through projective left modules. Applying this degree-zero comparison to $P^\bullet[-a]$ proves the first
isomorphism for every $a\in\mathbb Z$.

For the second isomorphism, we identify cycles and boundaries directly.
For every $R$-module $V$, degree-$a$ cycles of $\Hom_R^\bullet(P^\bullet,V)$ are maps
$c^{-a}(P^\bullet)\to V$. Its boundaries are those extending along $c^{-a}(P^\bullet)\hookrightarrow P^{1-a}$. Total acyclicity implies that every map from $c^{-a}(P^\bullet)$ to a projective module extends along this embedding. Thus the boundaries are precisely the projective factorizations.
This proves the isomorphism for arbitrary $V$; taking $V=c^0(Q^\bullet)$
gives the second isomorphism in the first display of the lemma.

The comparisons are induced by cokernel maps and respect shifted
composition. Finally, the nonpositive part of $P^\bullet$ is a deleted projective resolution of
$c^0(P^\bullet)$, so its complete Hom cohomology agrees with
ordinary Ext in degrees $a\ge1$.
\end{proof}

\begin{Rem}
For self-injective $R$ and $R$-modules $U,V$, we choose a complete
resolution $P^\bullet$ of $U$. The comparison lemma gives
\[
 H^a\Hom_R^\bullet(P^\bullet,V)\simeq\underline{\Hom}_R(U,V[a])
 \simeq\Ext_{R\text{-}\underline{\mathrm{mod}}}^a(U,V),
\]
using $c^{-a}(P^\bullet)\simeq\Omega^aU=U[-a]$ in the stable category.
Negative stable morphism spaces may be nonzero, although
negative ordinary Ext groups vanish by convention.
\end{Rem}


For symmetric algebras, the following duality recovers negative stable
morphism groups from nonnegative ones. Its composition formula will
also let us compare multiplication in the two ranges.

\begin{Lem}\label{tool:duality}
Let $R$ be symmetric. For all $R$-modules $U,V$ and all
$a\in\mathbb Z$, there are natural isomorphisms
\[
 \vartheta^a_{U,V}:\underline{\Hom}_R(V,U[-a-1])
 \xrightarrow{\sim}D\underline{\Hom}_R(U,V[a]).
\]
For each $U$, define
\[
 \lambda_U:\underline{\Hom}_R(U,U[-1])\longrightarrow k,
 \qquad \lambda_U(g)=\vartheta^0_{U,U}(g)(1_U).
\]
These isomorphisms satisfy
\[
 \vartheta^a_{U,V}(g)(f)
   =\lambda_U\bigl(g[a]\circ f\bigr)
\]
for $f\in\underline{\Hom}_R(U,V[a])$ and
$g\in\underline{\Hom}_R(V,U[-a-1])$.
\end{Lem}
\begin{proof}
We obtain the degree-zero duality from the Serre functor and then
shift to the other degrees. Since $R$ is symmetric, its Nakayama
functor is naturally isomorphic to the identity. Thus $\Omega_R=[-1]$ is a Serre functor
on the stable module category by \cite[Proposition~6.4.2]{KrauseHTR}, equivalently by the
Auslander--Reiten formula (for example, see \cite[Proposition~6.3.5]{KrauseHTR}).
Write the resulting natural isomorphism as
\[
 \vartheta^0_{U,V}:\underline{\Hom}_R(V,U[-1])
 \xrightarrow{\sim}D\underline{\Hom}_R(U,V),
\]
and define $\lambda_U$ as in the statement. Naturality gives
$\vartheta^0_{U,V}(g)(f)=\lambda_U(gf)$, the Serre trace identity
of \cite[Lemma~6.3.1]{KrauseHTR}.
For general $a$, use the natural shift identification and set
$\vartheta^a_{U,V}(g)=\vartheta^0_{U,V[a]}\bigl(g[a]\bigr)$.
The degree-zero identity then gives the stated composition formula.
\end{proof}

\paragraph{Stable images via the singularity category.}
Let $R$ be self-injective, and let
$q_R:\mathscr D^b(R)\to\mathscr D_{\rm sg}(R)
=\mathscr D^b(R)/\mathscr K^b(R\text{-}\mathrm{proj})$
be the quotient functor to the \emph{singularity category}.
The inclusion of modules as stalk complexes induces a triangle equivalence
(see \cite[Theorem~6.2.5]{KrauseHTR}):
\[
 \Gamma_R:R\text{-}\underline{\mathrm{mod}}
 \xrightarrow{\sim}\mathscr D_{\rm sg}(R).
\]
Choose its triangle-functor structure compatible with the convention for
connecting morphisms fixed in Subsection~\ref{prelim:stable}.
Fix a triangle quasi-inverse $\Gamma_R^{-1}$ and set
\[
 \Phi_R=\Gamma_R^{-1}\circ q_R:\;
 \mathscr D^b(R)\lra R\text{-}\underline{\mathrm{mod}}.
\]
For a bounded derived object $W^\bullet$, write
$\overline W=\Phi_R(W^\bullet)$ for its stable image, represented by an
$R$-module. The functor $\Phi_R$ preserves triangles and annihilates
perfect objects. Under the natural identifications $\overline M\simeq M$,
an ordinary module map $u:M\to N$ has stable image $\underline u$.

If $P^\bullet$ is a bounded-above complex of projective modules with
bounded cohomology and $n$ is a positive integer such that
$H^j(P^\bullet)=0$ for all $j\le-n$, then
\[
 \overline P\simeq(c^{-n}(P^\bullet))[n].
\]
Indeed, the kernel of the projection onto the brutal truncation in degrees
at most $-n$ is a bounded complex of projective modules. The truncated
complex is quasi-isomorphic to $(c^{-n}(P^\bullet))[n]$.

For a cochain map $f:P^\bullet\to Q^\bullet$ between such models,
the corresponding stable morphism is $\underline{c^{-n}(f)}[n]$
at a common sufficiently low cutoff. These identifications are natural
in cochain maps and compatible with changes of projective model and
cutoff, as well as with shift comparisons.

\paragraph{Perfect factorizations.}

The next lemma lets us prove nonvanishing in the stable category by
excluding perfect factorizations.

\begin{Lem}\label{lift:zero-perfect}
Let $R$ be self-injective. Suppose that $K^\bullet$ is a bounded-above
complex of projective $R$-modules with bounded cohomology and
$U$ is an $R$-module. If a derived morphism
$\alpha:K^\bullet\to U$ satisfies $\Phi_R(\alpha)=0$, then it
factors in $\mathscr D^b(R)$ through a bounded complex of projective
modules.
\end{Lem}
\begin{proof}
By definition $\Phi_R(\alpha)=\Gamma_R^{-1}q_R(\alpha)$.
Since $\Gamma_R^{-1}$ is an equivalence, $q_R(\alpha)=0$.
By the calculus of fractions, there is a morphism
$s:(K')^\bullet\to K^\bullet$ with perfect cone and $\alpha s=0$.
Complete $s$ to a triangle
\[
 (K')^\bullet\xrightarrow{s}K^\bullet\xrightarrow{t}C^\bullet
 \longrightarrow(K')^\bullet[1].
\]
Exactness of $\Hom_{\mathscr D^b(R)}(-,U)$ gives a morphism
$\beta:C^\bullet\to U$ with $\alpha=\beta t$.
Since $C^\bullet$ is perfect, it is isomorphic to a bounded
complex of projective modules, giving the required factorization;
see also \cite[Remark~3.2.3(2)]{KrauseHTR}.
\end{proof}

\subsection{Transfer to a symmetric algebra}

Part \textup{(1)} of the following theorem is the transfer principle of
\cite[Proposition~3.2]{O}. We include a proof with cohomological
indexing and explicit signs, valid in arbitrary characteristic.

\Needspace{10\baselineskip}
\begin{Theo}\label{frame:transfer}
Let $R$ be an algebra and let $P^\bullet$ be a totally acyclic complex of projective
$R$-modules such that $Z=c^0(P^\bullet)$ is nonprojective.
Set $A=R\ltimes DR$ and $M=A\otimes_RZ$. Then the following hold.
\begin{enumerate}
\item If $H^a\big(\Hom_R^\bullet(P^\bullet,Z)\big)=0$ for $a\ge1$ or $a\le-2$,
then $A$ is symmetric and ${_A}M$ is nonprojective with
$\Ext_A^i(M,M)=0$ for every $i>0$.
\item If $H^a\big(\Hom_R^\bullet(P^\bullet,Z)\big)\simeq k$ for
$a=-1,0$, then
\mbox{$\underline{\End}_A(M)\simeq k[\varepsilon]/(\varepsilon^2)$}
as $k$-algebras.
\end{enumerate}
\end{Theo}
\begin{proof}
Both assertions use the induced complex $A\otimes_RP^\bullet$
and a duality for $I^\bullet=DR\otimes_RP^\bullet$.
We first establish this duality and the exactness of these complexes.


For every projective $R$-module $L$, the natural Nakayama isomorphism
and projective Hom duality give
\begin{equation*}
 DR\otimes_RL\simeq D\Hom_R(L,R),\qquad
 \Hom_R(Z,DR\otimes_RL)\simeq D\Hom_R(L,Z).
\end{equation*}
Multiplying these termwise identifications by $(-1)^n$ in degree $n$ gives
isomorphisms of complexes
\[
 I^\bullet\simeq D\Hom_R^\bullet(P^\bullet,R),\qquad
 \Hom_R^\bullet(Z,I^\bullet)\simeq D\Hom_R^\bullet(P^\bullet,Z).
\]
Since $\cpx{P}$ is totally acyclic, the complex $\Hom_R^\bullet(P^\bullet, R)$ is exact.
Thus $I^\bullet$ is exact. Its terms are injective, since the
Nakayama functor sends projective modules to injective modules.

We can now use $A\otimes_RP^\bullet$ to resolve ${_A}M$. Its restriction
to $R$ is $P^\bullet\oplus I^\bullet$, so it is exact. Its terms
are projective over $A$, since the induction functor $A\otimes_R-$ preserves projective modules.
Right exactness gives
$c^0(A\otimes_RP^\bullet)=A\otimes_Rc^0(P^\bullet)=M$ and
$c^0(I^\bullet)=DR\otimes_RZ$.
The trivial extension $A$ is symmetric, hence self-injective.
Thus applying $\Hom_A(-,A)$ to the exact complex
$A\otimes_RP^\bullet$ gives an exact complex. Consequently,
$A\otimes_RP^\bullet$ is totally acyclic and is a complete resolution
of $M$.
Its nonpositive part
augments to a projective resolution of $M$:
$$
\cdots \lra A\otimes_RP^{-i}\lra\cdots\lra A\otimes_RP^{-1}\lra A\otimes_RP^0\lra M\lra 0.
$$
Let $\lambda:R\hookrightarrow A$ be the standard inclusion.

\smallskip\noindent
\textup{(1)} We compute the self-extensions of $M$ by
induction--restriction adjunction.
The adjoint pair
$(A\otimes_R-, \operatorname{Res}_\lambda)$ induces a series of isomorphisms
$$\Ext_A^i(M,M)=\Ext_A^i(A\otimes_RZ, A\otimes_RZ)\simeq \Ext_R^i(Z, A\otimes_RZ)\simeq \Ext_R^i(Z, Z)\oplus \Ext_R^i(Z, DR\otimes_RZ)$$ for every $i\ge 0$.
Note that $\Ext_R^i(Z,Z)\simeq H^i\big(\Hom_R^\bullet(P^\bullet,Z)\big)=0$ for $i>0$ by assumption.
Since $I^\bullet$ is exact with injective terms, its positive part gives
an injective coresolution of ${_R}DR\otimes_RZ$:
$$
0\lra DR\otimes_RZ\lra DR\otimes_RP^1\lra DR\otimes_RP^2\lra\cdots\lra DR\otimes_RP^i\lra\cdots.
$$
Consequently, for $i>0$, we have
\begin{equation*}
 \Ext_R^i(Z,DR\otimes_RZ)\simeq H^{i+1}\big(\Hom_R^\bullet(Z,I^\bullet)\big)
 \simeq H^{i+1}\big(D\Hom_R^\bullet(P^\bullet, Z)\big)\simeq D\big(H^{-i-1}\Hom_R^\bullet(P^\bullet,Z)\big)=0,
\end{equation*}
where the last equality follows from the assumption $H^a\big(\Hom_R^\bullet(P^\bullet,Z)\big)=0$ for $a\le-2$.
Thus $\Ext_A^i(M,M)=0$ for $i>0$.
Base change along the projection $\pi: A\twoheadrightarrow R$ sends projective modules to
projective modules. Since the composite of $\lambda$ with $\pi$ is the identity map of $R$, tensor associativity gives $R\otimes_AM\simeq R\otimes_AA\otimes_RZ\simeq Z$.
Thus projectivity of $M$ would contradict nonprojectivity of $Z$.

\smallskip\noindent
\textup{(2)} To determine the stable endomorphism algebra, we will
identify a one-dimensional square-zero ideal complementary to the
scalar multiples of $\underline{1_M}$. We first use the induced complete
resolution to compute the relevant vector spaces, and then describe
the endomorphisms and their multiplication explicitly.

Lemma~\ref{frame:comparison} gives
\mbox{$\underline{\End}_A(M)\simeq
H^0\big(\Hom_A^\bullet(A\otimes_RP^\bullet,M)\big)$}
for the complete resolution $A\otimes_RP^\bullet$ of $M$.
Induction--restriction adjunction on Hom complexes gives
\[
 H^0\big(\Hom_A^\bullet(A\otimes_RP^\bullet,M)\big)
 \simeq H^0\big(\Hom_R^\bullet(P^\bullet,Z)\big)
 \oplus H^0\big(\Hom_R^\bullet(P^\bullet,DR\otimes_RZ)\big).
\]
Applying Lemma~\ref{frame:comparison} to $P^\bullet$ in each
summand defines a composite isomorphism of vector spaces
\mbox{$\Phi:\underline{\End}_A(M)\xrightarrow{\sim}
\underline{\End}_R(Z)\oplus\underline{\Hom}_R(Z,DR\otimes_RZ)$}.
The summand $\underline{\End}_R(Z)$ is one-dimensional by the
assumption in degree zero. To compute
$\underline{\Hom}_R(Z,DR\otimes_RZ)$, we relate projective
factorizations to the boundaries in $\Hom_R^\bullet(Z,I^\bullet)$.

Let $q:I^0\twoheadrightarrow DR\otimes_RZ=c^0(I^\bullet)$ be the
cokernel map. Exactness of $I^\bullet$ gives an induced injection
$DR\otimes_RZ\hookrightarrow I^1$ whose image is the kernel of
the differential $I^1\to I^2$. It follows that
\[
 H^1\big(\Hom_R^\bullet(Z,I^\bullet)\big)
 \simeq
 \frac{\Hom_R(Z,DR\otimes_RZ)}
 {\{q\circ h\mid h\in\Hom_R(Z,I^0)\}}.
\]
We claim that the denominator consists precisely of maps factoring
through projective $R$-modules. Indeed, if a map factors through a
projective module, its second factor lifts along the epimorphism $q$,
so the map belongs to the denominator. Conversely, every map
$Z\to I^0$ extends along $Z\hookrightarrow P^1$, because $I^0$
is injective. Its composite with $q$ therefore factors through the
projective module $P^1$, proving the claim.
Consequently,
\[
 \underline{\Hom}_R(Z,DR\otimes_RZ)
 \simeq H^1\big(\Hom_R^\bullet(Z,I^\bullet)\big)
 \simeq D\big(H^{-1}\Hom_R^\bullet(P^\bullet,Z)\big)
 \simeq k.
\]
The last isomorphism uses the assumption in degree $-1$.

It remains to identify the multiplication using $\Phi$.
As an $R$-module,
$\operatorname{Res}_\lambda M=Z\oplus(DR\otimes_RZ)$.
For $f\in DR\subset A$, $z\in Z$ and $v\in DR\otimes_RZ$,
the action on $M$ is $f\cdot(z,v)=(0,f\otimes z)$.
Thus every $A$-endomorphism of $M$ has a unique expression
\[
 F_{\alpha,\beta}=
 \begin{pmatrix}
  \alpha&0\\
  \beta&DR\otimes_R\alpha
 \end{pmatrix},
 \qquad
 \alpha\in\End_R(Z),\quad
 \beta\in\Hom_R(Z,DR\otimes_RZ).
\]
Conversely, each such matrix defines an $A$-endomorphism.
For $\gamma\in\End_R(Z)$ and
$\delta\in\Hom_R(Z,DR\otimes_RZ)$, composition is given by
$F_{\alpha,\beta}\circ F_{\gamma,\delta}
=F_{\alpha\circ\gamma,\,
      \beta\circ\gamma+(DR\otimes_R\alpha)\circ\delta}$.
In particular, the endomorphisms $F_{0,\beta}$ form a square-zero
ideal in $\End_A(M)$.
To identify the image of this ideal under $\Phi$, we use the
induction--restriction adjunction
$\End_A(M)\simeq\Hom_R(Z,\operatorname{Res}_\lambda M)$, which sends
$F$ to $z\mapsto F(1_A\otimes z)$.
For $F=F_{\alpha,\beta}$, the resulting map is
$z\mapsto(\alpha(z),\beta(z))$. The cokernel identifications
in Lemma~\ref{frame:comparison}, used to construct $\Phi$, therefore
give $\Phi(\underline{F_{\alpha,\beta}})
=(\underline\alpha,\underline\beta)$.
Thus the image $\mathcal J$ of this square-zero ideal in
$\underline{\End}_A(M)$ is
$\Phi^{-1}(0\oplus\underline{\Hom}_R(Z,DR\otimes_RZ))$ and has
dimension one. The one-dimensional space $\underline{\End}_R(Z)$
is generated by $\underline{1_Z}$, which is nonzero because $Z$ is nonprojective.
As $F_{1_Z,0}=1_M$ and $\Phi(\underline{1_M})=(\underline{1_Z},0)$, we obtain
$\underline{\End}_A(M)=k\,\underline{1_M}\oplus\mathcal J$, with
$\dim_k\mathcal J=1$ and $\mathcal J^2=0$.
Choosing a nonzero $\varepsilon\in\mathcal J$ gives the required
isomorphism of $k$-algebras.
\end{proof}

\subsection{A homological cone criterion}
\label{frame:cone-section}

We construct a nonprojective cokernel of a totally acyclic complex
whose complete self-Hom is concentrated in degrees $-1$ and $0$.
The construction uses an endomorphism sequence for an orthogonal
triangle and two normalized bimodule maps, as in
Proposition~\ref{frame:factory}. Theorem~\ref{frame:transfer} then
gives a nonprojective self-orthogonal module over a symmetric algebra
in Corollary \ref{frame:combined}.

Throughout this subsection, $E$ is {\bf self-injective}. If an
$E$-bimodule $L$ is separately projective, then
$L\otimes_E(-)$ is exact and preserves projective modules.
Since projective and injective $E$-modules coincide, this functor
induces a triangle functor on $E\text{-}\underline{\mathrm{mod}}$,
with natural comparisons
$L\otimes_E X[a]\simeq(L\otimes_E X)[a]$.
These comparisons arise from tensoring projective-injective
cosyzygy sequences with the specified shift convention. They are
natural in the module and in maps of such bimodules.
All shifts and connecting maps use the conventions fixed above.
The inputs for the concrete construction are specified in
Section~\ref{app:section}.

\paragraph{Endomorphisms of an orthogonal triangle.}

Under an orthogonality assumption, the next lemma relates shifted
endomorphisms of a triangle's third object to morphisms involving its
first two objects. We will apply this sequence to the cone triangle.

\begin{Lem}
\label{frame:orthogonal-triangle}
Let $\mathscr T$ be a triangulated category, and let
\[
 B\xrightarrow{u}A\xrightarrow{j}C\xrightarrow{q}B[1]
\]
be a distinguished triangle satisfying
$\Ext_{\mathscr T}^n(A,B)=0$ for every $n\in\mathbb Z$.
Then there exists a long exact sequence of abelian groups
\[
\begin{aligned}
\cdots&\longrightarrow
 \Ext_{\mathscr T}^{a-1}(A,A)\oplus
 \Ext_{\mathscr T}^{a-1}(B,B)
 \xrightarrow{\delta^{a-1}}\Ext_{\mathscr T}^{a-1}(B,A)\\
&\xrightarrow{\kappa_a}\Ext_{\mathscr T}^a(C,C)
 \xrightarrow{\pi_a}
 \Ext_{\mathscr T}^a(A,A)\oplus\Ext_{\mathscr T}^a(B,B)
 \xrightarrow{\delta^a}\Ext_{\mathscr T}^a(B,A)
 \longrightarrow\cdots
\end{aligned}
\]
with the associated maps given by
\[
 \delta^a:\;\;\Ext_{\mathscr T}^a(A,A)\oplus
                   \Ext_{\mathscr T}^a(B,B)
       \longrightarrow\Ext_{\mathscr T}^a(B,A),
 \qquad (g,h)\mapsto gu-u[a]h,
\]
\[
\kappa_a:\;\; \Ext_{\mathscr T}^{a-1}(B,A)
       \longrightarrow\Ext_{\mathscr T}^a(C,C),
 \qquad c\mapsto j[a]c[1]q,
\]
\[
\pi_a:\;\; \Ext_{\mathscr T}^a(C,C)\longrightarrow
\Ext_{\mathscr T}^a(A,A)\oplus\Ext_{\mathscr T}^a(B,B)\qquad
z\mapsto (g_z,h_z),
\]
where the pair $(g_z,h_z)$ is determined by $j[a]g_z=zj$ and $h_z[1]q=(-1)^a q[a]z$.

In particular, for each $a\in\mathbb{Z}$, there is a short exact sequence of abelian groups
\begin{equation*}
 0\longrightarrow\operatorname{coker}(\delta^{a-1})
 \longrightarrow\Ext_{\mathscr T}^a(C,C)
 \longrightarrow\ker(\delta^a)\longrightarrow0.
\end{equation*}
\end{Lem}
%Here $\kappa_a(c)=j[a]c[1]q$ for $c\in \Ext_{\mathscr T}^{a-1}(B,A)$, and $\pi_a(z)=(g,h)$ is
%determined by
%\[
% j[a]g=zj,\qquad h[1]q=(-1)^a q[a]z.
%\]
\begin{proof}
We first show that the map $\pi_a$ is well defined.

Since $\Hom_\mathscr{T}(A,B[a+1])=0$, the long exact Hom
sequence of the shifted triangle gives, for each
$z\in\Ext_{\mathscr T}^a(C,C)$, a morphism
$g_z:A\to A[a]$ such that $j[a]g_z=zj$.
Applying the morphism-of-triangles axiom to the rotations of
the two rows below and to the pair $(g_z,z)$ then gives a
morphism $h_z:B\to B[a]$ making the following diagram commute:
$$
(\ast)\quad
\xymatrix{
B\ar[r]^-{u}\ar[d]^-{h_z}& A\ar[r]^-{j}\ar[d]^-{g_z}&C\ar[r]^-{q}\ar[d]^-{z}& B[1]\ar[d]^-{h_z[1]}\\
B[a]\ar[r]^-{u[a]}& A[a]\ar[r]^-{j[a]}&C[a]\ar[r]^-{(-1)^aq[a]}& B[a+1]
}
$$
If there are other morphisms $g_z': A\to A[a]$ and $h_z':B\to B[a]$ that also make
the above diagram commute, then $g_z-g_z'=u[a]s$ and $h_z[1]-h_z'[1]=-t[1]u[1]$ for some morphisms
$s, t: A\to B[a]$. The assumption $\Hom_\mathscr{T}(A, B[a])=0$ forces $(g_z, h_z)=(g_z', h_z')$. Thus $\pi_a$ is well defined. This also implies $\delta^a\pi_a=0$. For $z=\kappa_a(c)=j[a]c[1]q$, we have $zj=0$ and $q[a]z=0$. The defining identities for $\pi_a$, together with the uniqueness of $(g_z,h_z)$, therefore give $\pi_a\kappa_a=0$. Moreover, since $u[1]q=0=j[a]u[a]$, we have $\kappa_a\delta^{a-1}=0$.

Thus the long sequence in Lemma \ref{frame:orthogonal-triangle} is a complex of abelian groups. To show its exactness, it suffices to show the inclusions
$$\ker(\delta^a)\subseteq \operatorname{im}(\pi_a),\;\; \ker(\pi_a)\subseteq\operatorname{im}(\kappa_a),\;\;
\ker(\kappa_a)\subseteq \operatorname{im}(\delta^{a-1}).$$

Let $(g,h)\in\Ext_{\mathscr T}^a(A,A)\oplus\Ext_{\mathscr T}^a(B,B)$ with $\delta^a(g,h)=gu-u[a]h=0$.
By the morphism-of-triangles axiom, there exists a morphism $z_{g,h}:C\to C[a]$ making the following diagram of triangles commutative:
$$
\xymatrix{
B\ar[r]^-{u}\ar[d]^-{h}& A\ar[r]^-{j}\ar[d]^-{g}&C\ar[r]^-{q}\ar[d]^-{z_{g,h}}& B[1]\ar[d]^-{h[1]}\\
B[a]\ar[r]^-{u[a]}& A[a]\ar[r]^-{j[a]}&C[a]\ar[r]^-{(-1)^aq[a]}& B[a+1].
}
$$
By the preceding discussion, we have $\pi_a(z_{g,h})=(g,h)$. This shows $\ker(\delta^a)\subseteq \operatorname{im}(\pi_a)$.

Let $z\in \Ext_{\mathscr T}^a(C,C)$ with $\pi_a(z)=(g_z, h_z)=(0,0)$. The diagram $(\ast)$ gives
$zj=0=(-1)^aq[a]z$. It follows that there exists a morphism $\ell: B[1]\to C[a]$ with $z=\ell q$.
This leads to $q[a]z=q[a]\ell q=0$. By exactness of the Hom sequence associated with the first row of $(\ast)$, there exists another morphism $m:A[1]\to B[a+1]$ such that $(-1)^aq[a]\ell=-mu[1]$. Then the assumption $\Hom_\mathscr{T}(A[1], B[a+1])=0$ forces $m=0$ and so $(-1)^aq[a]\ell=0$. By exactness of the Hom sequence associated with the second row of $(\ast)$,
there exists a morphism $c: B\to A[a-1]$ such that $\ell=j[a]c[1]$. Thus $z=\ell q=j[a]c[1]q=\kappa_a(c)$.
This shows $\ker(\pi_a)\subseteq\operatorname{im}(\kappa_a)$.

Finally, let $c\in\Ext_{\mathscr T}^{a-1}(B,A)$ with $\kappa_a(c)=j[a]c[1]q=0$. Then there exists a morphism
$t: A[1]\to C[a]$ with $j[a]c[1]=tu[1]$. Since $\Hom_\mathscr{T}(A[1], B[a+1])\simeq \Hom_\mathscr{T}(A, B[a])=0$, we have $(-1)^aq[a]t=0$.
By exactness of the Hom sequence associated with the second row of $(\ast)$, there exists a morphism $x:A\to A[a-1]$  with $t=j[a]x[1]$.
Then $j[a]c[1]=tu[1]=j[a]x[1]u[1]$, or equivalently, $j[a-1](c-xu)=0$.
It follows that there exists a morphism $y:B\to B[a-1]$ with $c-xu=u[a-1]y$, and hence $c=\delta^{a-1}(x,-y)$. This shows $\ker(\kappa_a)\subseteq \operatorname{im}(\delta^{a-1})$.
\end{proof}

\paragraph{Column modules and the cone.}

We now apply this exact sequence in the homotopy category of complete resolutions. The column functors for a triangular matrix algebra will provide the orthogonality required by the lemma.

Let ${_E}X$ be a nonprojective $E$-module, $F$ a separately projective $E$-bimodule and
$v\in\underline{\Hom}_E(X,F\otimes_E X)$. For $a\in\mathbb Z$, we put
$$ H^a=\underline{\Hom}_E(X,X[a]),\qquad V^a=\underline{\Hom}_E(X,(F\otimes_E X)[a]).$$
For $g:X\to X[a]$, we write $F(g)$ for the induced stable map $F\otimes_Eg: F\otimes_E X \to (F\otimes_E X)[a]$, identifying $(F\otimes_E X)[a]$ with $F\otimes_E (X[a])$ via the natural shift comparison for $F\otimes_E(-)$. Define
\begin{equation*}
\Delta^a: H^a\times H^a\longrightarrow V^a,
\qquad\Delta^a(g,h)=F(g)v-v[a]h.
\end{equation*}
Both summands are stable maps from $X$ to $(F\otimes_E X)[a]$.

We choose an embedding $i:X\hookrightarrow J_X$ into an injective module; for example, a natural choice for $i$ is an injective envelope. Since $E$ is self-injective, splicing a projective resolution of $X$ with an injective coresolution of $X$ beginning with $i$ gives a complete resolution $P^\bullet$ such that $$P^1=J_X,\qquad c^0(P^\bullet)=X,\qquad d_P^0=i\epsilon$$
with a surjection $\epsilon:P^0\twoheadrightarrow X$. Let $\widetilde v:X\to F\otimes_E X$ be a homomorphism of $E$-modules, a representative of $v$ in the stable category. The complex $F\otimes_E P^\bullet$ is exact with projective terms, so Lemma~\ref{frame:comparison} gives a cochain representative of
$v$. Its induced cokernel map differs from $\widetilde v$ by a
factorization $X\to T\to F\otimes_E X$ through a projective
module $T$, where the first map extends across $X\hookrightarrow P^1$
because $T$ is injective, and the second lifts through
$F\otimes_E P^0\twoheadrightarrow F\otimes_E X$ by projectivity.
Then the composite $P^1\to F\otimes_E P^0$ defines a nullhomotopic
correction. Thus the representative can be chosen as a cochain lift
$\widetilde u\colon P^\bullet\to F\otimes_E P^\bullet$ inducing the specified ordinary map $\widetilde v$.

We now consider the triangular matrix algebra
$$\Lambda=\begin{pmatrix}E&0\\F&E\end{pmatrix}.$$
By \cite[Chapter III.2]{ARS}, a left $\Lambda$-module is a triple $(L, N, \alpha:F\otimes_E L\to N)$, and a morphism from this triple to another triple $(L', N',\alpha': F\otimes_E L'\to N')$ is a pair $(r: L\to L', s: N\to N')$ satisfying
$s\alpha=\alpha'(1_F\otimes_E r)$.

\Needspace{5\baselineskip}
Define two additive functors between module categories:
\[
\begin{aligned}
D_0\colon E\modcat&\lra\Lambda\modcat,\quad L\mapsto(L,F\otimes_E L,1_{F\otimes_E L}),
\\
D_1\colon E\modcat&\lra\Lambda\modcat,\quad L\mapsto(0,L,0).
\end{aligned}
\]
The lift $\widetilde u$ determines a cochain map
$u:D_1(P^\bullet)\to D_0(P^\bullet)$, with degree-$j$ component
$u^j=(0,\widetilde u^j)$.
With the signed cone convention \eqref{prelim:signed-cone}, define
$$P_Z^\bullet=\operatorname{Cone}(u)^\bullet,\qquad Z=c^0(P_Z^\bullet).$$

\Needspace{5\baselineskip}
The following theorem gives an abstract form of the triangular cone
argument in \cite[Sections~7.2--7.3]{O}, over a general self-injective
algebra. Our proof uses Lemma~\ref{frame:orthogonal-triangle} to
express the required complete self-Hom vanishing in terms of the maps
$\Delta^a$.

\begin{Theo}\label{frame:cone}
The construction above has the following properties.
\begin{enumerate}
\item\label{frame:cone-properties}
The complex $P_Z^\bullet$ is totally acyclic with projective
terms, and $Z$ is nonprojective. Moreover,
\begin{equation*}
 Z=\left(X,
 \frac{F\otimes_E X\oplus J_X}
      {\{(\widetilde v(x),-i(x)):x\in X\}},\alpha\right),
\end{equation*}
where $\alpha$ is induced by inclusion of $F\otimes_E X$.

\item\label{frame:cone-vanishing}
The vanishing condition
\begin{equation*}
 H^a\Hom_\Lambda^\bullet(P_Z^\bullet,Z)=0
       \quad(a\geq1\text{ or }a\leq-2)
\end{equation*}
holds if and only if $\Delta^a$ is an isomorphism whenever $a\geq1$
or $a\leq-3$, $\Delta^0$ is surjective, and $\Delta^{-2}$ is
injective.

\item\label{frame:cone-support}
Under the equivalent vanishing conditions in
\textup{(\ref{frame:cone-vanishing})}, suppose additionally that
$\Delta^{-1}$ is an isomorphism, $\ker(\Delta^0)\cong k$, and
$\operatorname{coker}(\Delta^{-2})\cong k$. Then
\begin{equation*}
 H^a\Hom_\Lambda^\bullet(P_Z^\bullet,Z)\cong
 \begin{cases}k,&a=-1,0,\\0,&a\notin\{-1,0\}\end{cases}.
\end{equation*}
\end{enumerate}
\end{Theo}

\begin{proof}
\textup{(\ref{frame:cone-properties})}
We first prove total acyclicity for the two column complexes;
the cone sequence will then give total acyclicity of $P_Z^\bullet$.
The column functors are exact because $F_E$ is projective.
They are induction functors along the diagonal idempotents. In particular, they are fully faithful.
Moreover, each is left adjoint to an exact component restriction, so it preserves projectives.
The same adjunctions give natural identifications
\begin{equation*}
(\diamondsuit)\quad
\begin{aligned}
  \Hom_\Lambda(D_iL,D_iL')&\simeq\Hom_E(L,L')\quad(i=0,1),\;
  &\Hom_\Lambda(D_0L,D_1L')&=0,\;\\
  \Hom_\Lambda(D_1L,D_0L')&\simeq\Hom_E(L,F\otimes_E L'),
 \end{aligned}
\end{equation*}
where $L, L'\in E\modcat$.
Thus $D_0(P^\bullet)$ and $D_1(P^\bullet)$ are exact complexes
of projective modules. Their Hom complexes into
$\Lambda=D_0(E)\oplus D_1(E)$ are direct sums of
$\Hom_E^\bullet(P^\bullet,E)$,
$\Hom_E^\bullet(P^\bullet,F)$, and zero. These are exact because
$P^\bullet$ is totally acyclic and ${_E}F$ is projective. Both columns are therefore totally acyclic.
The degreewise split sequence
$0\to D_0(P^\bullet)\to P_Z^\bullet \to D_1(P^\bullet)[1]\to 0$
and the sequence obtained by applying
$\Hom_\Lambda^\bullet(-,\Lambda)$ prove total acyclicity of
$P_Z^\bullet$.

We identify $Z$ from the cone differential. Its upper component
is $X$. On the lower component, we first take the cokernel of
$1_F\otimes_Ed_P^{-1}$. Since $\widetilde u$ induces
$\widetilde v$ and $d_P^0=i\epsilon$, the remaining differential is
\[
 \begin{pmatrix}
 (1_F\otimes_E\epsilon)\widetilde u^0\\-d_P^0
 \end{pmatrix}
 =\begin{pmatrix}\widetilde v\\-i\end{pmatrix}\epsilon:\;\;
 P^0\longrightarrow(F\otimes_E X)\oplus J_X.
\]
Since $\epsilon$ is surjective, this map has the same cokernel
as $(\widetilde v,-i)$, namely the pushout of $i$ along
$\widetilde v$. This gives the asserted quotient description
and the short exact sequence
\begin{equation*}
 0\longrightarrow D_1(X)
 \xrightarrow{(\widetilde v,-D_1(i))}
 D_0(X)\oplus D_1(J_X)\longrightarrow Z\longrightarrow0.
\end{equation*}
Here the first component denotes the induced $\Lambda$-module map
$(0,\widetilde v):D_1(X)\to D_0(X)$.
Since pushouts preserve monomorphisms in $E\text{-}\mathrm{mod}$,
the structure map $\alpha$ is injective.
The upper-component functor sends projective $\Lambda$-modules
to projective $E$-modules. Since the upper component of $Z$ is
the nonprojective module $X$, the module $Z$ is nonprojective.

\textup{(\ref{frame:cone-vanishing})}
To express the required vanishing in terms of $\Delta^a$, we apply
Lemma~\ref{frame:orthogonal-triangle} to the cone triangle.
We must first verify orthogonality and identify the difference map.
Let $\mathscr K=\mathscr K_{\mathrm{tac}}(\Lambda\text{-}\mathrm{proj})$
be the homotopy category of totally acyclic complexes of
projective modules. It is triangulated by the degreewise split
cone argument in \textup{(\ref{frame:cone-properties})}.
Put $\mathcal A^\bullet=D_0(P^\bullet)$ and
$\mathcal B^\bullet=D_1(P^\bullet)$. The cone triangle is
\[
 \mathcal B^\bullet\lraf{u}\mathcal A^\bullet
 \longrightarrow P_Z^\bullet\longrightarrow\mathcal B^\bullet[1].
\]
The column-Hom identifications (see $(\diamondsuit)$) and Lemma~\ref{frame:comparison}
give, for every integer $a$,
\begin{equation*}
 \begin{aligned}
 \Ext_{\mathscr K}^a(\mathcal A^\bullet,\mathcal A^\bullet)&\cong H^a,
 &\Ext_{\mathscr K}^a(\mathcal B^\bullet,\mathcal B^\bullet)&\cong H^a,\\
 \Ext_{\mathscr K}^a(\mathcal A^\bullet,\mathcal B^\bullet)&=0,
 &\Ext_{\mathscr K}^a(\mathcal B^\bullet,\mathcal A^\bullet)&\cong V^a.
 \end{aligned}
\end{equation*}
The identifications respect shifts and composition.
In the lower component, composition on $\mathcal A^\bullet$ is
computed by tensoring with $F$, whereas composition on
$\mathcal B^\bullet$ acts on the source of $u$.
Thus the difference map of Lemma~\ref{frame:orthogonal-triangle}
becomes $\Delta^a(g,h)=F(g)v-v[a]h$.
By Lemma~\ref{frame:comparison}, $\Ext_{\mathscr K}^a(P_Z^\bullet,P_Z^\bullet)
\cong H^a\Hom_\Lambda^\bullet(P_Z^\bullet,Z)$.
Lemma~\ref{frame:orthogonal-triangle} therefore yields a short
exact sequence
\[
 0\longrightarrow\operatorname{coker}(\Delta^{a-1})
 \longrightarrow H^a\Hom_\Lambda^\bullet(P_Z^\bullet,Z)
 \longrightarrow\ker(\Delta^a)\longrightarrow0.
\]
Its middle term vanishes if and only if both outer terms vanish.
For every $a\geq1$, this requires injectivity of $\Delta^a$ and
surjectivity of $\Delta^{a-1}$. Together these conditions are
equivalent to bijectivity for all $a\geq1$ and surjectivity in
degree zero. For every $a\leq-2$, they require injectivity of
$\Delta^a$ and surjectivity of $\Delta^{a-1}$, equivalently
bijectivity for all $a\leq-3$ and injectivity in degree $-2$.
This proves both directions of the vanishing criterion.

\textup{(\ref{frame:cone-support})}
Under these vanishing conditions, if $\Delta^{-1}$ is also
bijective, the short exact sequences in degrees zero and $-1$ give
$H^0\Hom_\Lambda^\bullet(P_Z^\bullet,Z)\cong\ker(\Delta^0)$ and
$H^{-1}\Hom_\Lambda^\bullet(P_Z^\bullet,Z)\cong\operatorname{coker}(\Delta^{-2})$.
The additional one-dimensionality assumptions give the asserted
complete self-Hom groups.
\end{proof}

\paragraph{Two normalized lifts.}

To apply the cone criterion, we construct $F$ and $v$ from two bimodule maps with the same nonzero stable evaluation. The following proposition gives an abstract form of the construction in
\cite[Section~7.1]{O}, replacing the explicit scalar conditions there by an invertibility condition on the difference of the induced actions.

\begin{Prop}\label{frame:factory}
Let $E$ be self-injective. Suppose that $X$ is an $E$-module and $Y,U_1,U_2$ are separately projective $E$-bimodules. For every $a\in\mathbb{Z}$, set
\[
 H^a=\underline{\Hom}_E(X,X[a]),\qquad
 W^a=\underline{\Hom}_E(X,(Y\otimes_EX)[a]).
\]
Assume the following conditions.
\begin{enumerate}
\item \emph{Support conditions.}\label{frame:factory-support}
$H^0=k1_X$, $H^{-3}=H^{-2}=0$, and
\[
 W^a\cong
 \begin{cases}
 k,&a=-3,0,\\
 0,&a\notin\{-3,0\}.
 \end{cases}
\]

\item \emph{Normalized lifts.}\label{frame:factory-lifts}
There are stable isomorphisms
$\phi_i:U_i\otimes_EX\xrightarrow{\sim}X$
and ordinary bimodule maps $g_i:U_i\to Y$.
Tensoring $g_i$ with $X$ gives the ordinary maps
\[
 g_i\otimes_E1_X:U_i\otimes_EX\longrightarrow Y\otimes_EX,
 \qquad u\otimes x\longmapsto g_i(u)\otimes x,\;\; u\in U_i, x\in X.
\]
We require that
\[
 \underline{g_1\otimes_E1_X}\circ\phi_1^{-1}
 =\underline{g_2\otimes_E1_X}\circ\phi_2^{-1}
 =w\ne0
 \quad\text{in }W^0.
\]

\item \emph{Invertible difference of induced actions.}
\label{frame:factory-actions}
For each $\xi\in H^a$, define
$$T_i^a:H^a\to H^a,\quad \xi\mapsto \phi_i[a](U_i\otimes_E\xi)\phi_i^{-1},$$
using the natural comparison with shifts. We require that
$T_1^a-T_2^a:H^a\to H^a$ is invertible whenever $a\ne0$ and $H^a\ne0$.
\end{enumerate}

Choose a surjection $\pi:Q\twoheadrightarrow Y$ with $Q$
projective over $E^e$, and define $F$ by
\[
 0\longrightarrow F\lraf{\iota_F}U_1\oplus U_2\oplus Q
 \xrightarrow{(g_1,-g_2,\pi)}Y\longrightarrow0.
\]
Define $r_i=\operatorname{pr}_{U_i}\circ\iota_F:F\to U_i$ for $i=1,2$,
where $\operatorname{pr}_{U_i}$ is the canonical projection
onto the summand $U_i$. Put
$$V^a=\underline{\Hom}_E(X,(F\otimes_EX)[a]).$$
Then:
\begin{enumerate}
\item\label{frame:factory-normalized-class}
There is a unique stable morphism
$v:X\to F\otimes_EX$ such that
\[
 \phi_i\circ\underline{r_i\otimes_E1_X}\circ v=1_X
 \qquad(i=1,2).
\]

\item\label{frame:factory-nonexceptional}
For this $v$, the associated maps
\begin{equation*}
\Delta^a: H^a\times H^a\longrightarrow V^a,
\qquad\Delta^a(g,h)=F(g)v-v[a]h \quad\mbox{for}\;\; g, h\in H^a
\end{equation*}
are isomorphisms for every $a\notin\{0,-2\}$.

\item\label{frame:factory-exceptions}
In the two exceptional degrees,
\[
\begin{aligned}
 &\Delta^0\text{ is surjective},
 &&\ker(\Delta^0)\cong k;\\
 &\Delta^{-2}\text{ is injective},
 &&\operatorname{coker}(\Delta^{-2})\cong k.
\end{aligned}
\]
\end{enumerate}
Consequently, the pair $(F,v)$ satisfies the equivalent vanishing
conditions and the additional hypotheses of
Theorem~\ref{frame:cone}\textup{(\ref{frame:cone-support})}.
Thus the cone construction has complete self-Hom supported
precisely in degrees $-1$ and $0$.
\end{Prop}

\begin{proof}
We use the evaluated defining sequence both to construct $v$ and to compute $\Delta^a$. First we check that evaluation gives a triangle in the stable category. Since $Q$ is $E^e$-projective and $Y$ is projective on each $E$-side, the defining sequence splits separately on the left and right. Thus $F$ is separately projective, and the right splitting preserves exactness after tensoring with $X$.
Moreover, ${_E}Q\otimes_E X$ is projective: evaluation of the free bimodule $E\otimes_kE$ is
$E\otimes_kX$, and this property passes to sums and summands.

\textup{(\ref{frame:factory-normalized-class})}
To construct the normalized class $v$, we will identify $V^0$
with the diagonal in $(H^0)^2$.
Using the stable isomorphisms $\phi_i$ and the vanishing of
$Q\otimes_E X$ in the stable category, the evaluated short
exact sequence gives a distinguished triangle
\[
 F\otimes_E X\xrightarrow{p}X\oplus X
    \xrightarrow{(w,-w)}Y\otimes_E X
    \longrightarrow(F\otimes_E X)[1].
\]
Here $p=(p_1,p_2)$ with $p_i=\phi_i\circ\underline{r_i\otimes_E1_X}:F\otimes_EX\to X$.
Applying stable Hom from $X$ gives the exact sequence of $k$-vector spaces:
\begin{equation*}
(H^{a-1})^2\lraf{\varepsilon^{a-1}}W^{a-1}\lra V^a\lraf{p^a}(H^a)^2\lraf{\varepsilon^a}W^a,
\qquad\varepsilon^a(g,h)=w[a]g-w[a]h.
\end{equation*}
The map $p^a$ is induced by postcomposition with $p[a]$; in particular, $p^0(f)=(p_1f,p_2f)$.
Identify $H^0=k$ by $c\mapsto c1_X$. Then
$\varepsilon^0(c,d)=(c-d)w$ is surjective, where $c,d\in k$.
Since $W^{-1}=0$, the map $p^0$ identifies $V^0$ with $\ker(\varepsilon^0)=k(1,1)$. This gives a unique $v\in V^0$ such that $p^0(v)=(1_X,1_X)$, as required.

\textup{(\ref{frame:factory-nonexceptional})}
We compare $\Delta^a$ with an endomorphism of $(H^a)^2$ by
postcomposing with $p^a$.
Note that
$$T_i^a(g)p_i=\phi_i[a](U_i\otimes_Eg)\phi_i^{-1}\phi_i(\underline{r_i\otimes_E1_X})
=\phi_i[a](U_i\otimes_Eg)(\underline{r_i\otimes_E1_X})
=\phi_i[a](\underline{r_i\otimes_E1_{X[a]}})F(g)=p_i[a]F(g).$$
The normalization $p_iv=1_X$ gives the two component calculations
\[
 p_i[a]F(g)v=T_i^a(g)p_iv=T_i^a(g),\qquad
 p_i[a]v[a]h=(p_iv)[a]h=h.
\]
Subtracting them yields $p^a\Delta^a(g,h)=(T_1^a(g)-h,\ T_2^a(g)-h)$, that is,
$$
\xymatrix{
p^a\Delta^a:\quad(H^a)^2\ar[rr]^-{\scriptstyle\begin{pmatrix}T_1^a&-1_{H^a}\\T_2^a&-1_{H^a}\end{pmatrix}}&&(H^a)^2.}
$$
Thus it suffices to prove that $p^a$ and $p^a\Delta^a$ are both invertible in every nonexceptional degree.

We first determine $p^a$ from the Hom exact sequence above. In degree $-3$, the vanishings $W^{-4}=H^{-3}=0$ give $V^{-3}=0$. For $a\notin\{0,-2,-3,1\}$, both $W^{a-1}$ and $W^a$ vanish, so $p^a$ is an isomorphism.
The same holds in degree $-3$, where both its source and target are zero. In degree one, surjectivity of $\varepsilon^0$ makes $W^0\to V^1$ zero; the vanishing of $W^1$ then makes $p^1$ an isomorphism. Hence $p^a$ is an isomorphism for every $a\notin\{0,-2\}$.

For each $a\notin\{0,-2\}$ with $H^a=0$, the isomorphism $p^a$ gives $V^a=0$, so $\Delta^a$ is an isomorphism. If $H^a\ne0$, subtracting the second component from the first makes $p^a\Delta^a$ block triangular, with diagonal entries $T_1^a-T_2^a$ and $-1_{H^a}$. Both are invertible by
assumption \textup{(\ref{frame:factory-actions})}. Together with the isomorphism $p^a$, this proves that
$\Delta^a$ is an isomorphism in every nonexceptional degree,
including $a=-1$.

\textup{(\ref{frame:factory-exceptions})}
It remains to check the two exceptional degrees.
In degree zero, $k$-linearity gives
$F(c1_X)=c1_{F\otimes_E X}$, so $\Delta^0(c,d)=(c-d)v$.
Thus $\Delta^0$ is surjective with kernel $k(1,1)$.
In degree $-2$, the vanishings $H^{-3}=H^{-2}=0$ identify
$V^{-2}$ with $W^{-3}\cong k$. The source $(H^{-2})^2$ is zero,
so $\Delta^{-2}$ is injective
with cokernel isomorphic to $k$.
These conclusions verify the vanishing conditions and the
additional hypotheses of
Theorem~\ref{frame:cone}\textup{(\ref{frame:cone-support})}.
Since $W^0\ne0$, the module $X$ is nonprojective, as required
for the cone construction.
\end{proof}

The following corollary combines the construction from two normalized
lifts with the cone criterion and the transfer theorem.

\begin{Koro}\label{frame:combined}
Assume the hypotheses of Proposition~\ref{frame:factory}, and let
$(F,v)$ be the pair constructed there. The associated triangular cone
construction gives a nonprojective $\Lambda$-module $Z$ with a complete
resolution $P_Z^\bullet$. The groups
$H^a\Hom_\Lambda^\bullet(P_Z^\bullet,Z)$ are isomorphic to $k$ for
$a=-1,0$ and vanish otherwise.
Moreover, $A=\Lambda\ltimes D\Lambda$ is symmetric, and
$M=A\otimes_\Lambda Z$ is a nonprojective self-orthogonal $A$-module.
\end{Koro}
\begin{proof}
Proposition~\ref{frame:factory} and parts
\textup{(\ref{frame:cone-properties}) and (\ref{frame:cone-support})}
of Theorem~\ref{frame:cone} give the assertions about $Z$ and its complete
resolution. The conclusion for $A$ and $M$ follows from
Theorem~\ref{frame:transfer} with $R=\Lambda$.
\end{proof}

\section{General Frobenius seeds and a fixed signed input}
\label{seed:section}

We prepare the two factors for the tensor product of algebras in
Section~\ref{ind:section}. Proposition~\ref{general:seed} gives their
common homological properties, and a quantum exterior algebra provides
the fixed involution.

\subsection{General Frobenius seeds}

The next proposition constructs the endomorphism algebra and its
distinguished simple module. The proof gives a projective resolution
of the simple module. It also resolves the left and right dual regular modules.

\begin{Prop} \label{general:seed}
Let $H$ be a finite-dimensional Frobenius $k$-algebra, and let
$M_0$ be an $H$-module with no nonzero projective direct summands.
Assume that
\begin{equation*}
 \tau_HM_0\simeq M_0,\qquad \underline{\End}_H(M_0)\simeq k,\qquad
 \Ext_H^i(M_0,M_0)=0\quad(i\ge2).
\end{equation*}
Set $N=H\oplus M_0$ and $C=\End_H(N)^{\mathrm{op}}$, and let
$e,f\in C$ be the idempotents corresponding to $H,M_0$,
respectively. Then:
\begin{enumerate}
\item\label{general:seed-simple}
The idempotent $f$ is primitive and $C/CeC\simeq k$. The quotient $S=C/CeC$ is a one-dimensional
simple left $C$-module. We write
$\chi_C:C\twoheadrightarrow C/CeC\simeq k$ for the quotient
character; thus $S=k_{\chi_C}$.

\item\label{general:seed-homology}
The extension groups satisfy
\begin{equation*}
 \Ext_C^i(S,S)=
 \begin{cases}k&i=0,\\0&i>0\end{cases},
 \qquad
 \Ext_C^i(S,C)=
 \begin{cases}DS&i=2,\\0&i\ne2\end{cases},
 \quad(i\ge0).
\end{equation*}
The isomorphisms in the second family are right $C$-linear.

\item\label{general:seed-dimension}
The homological dimensions are
\begin{equation*}
 \pd_C(DC)=\pd_{C^{\mathrm{op}}}(DC)
 =\id_C(C)=\id_{C^{\mathrm{op}}}(C)=2.
\end{equation*}

\item\label{general:seed-basic}
If $H$ is connected and split basic with $r$ simple modules, then
$C$ is connected and split basic with $r+1$ simple modules.

\end{enumerate}
\end{Prop}

For the proof and the subsequent seed constructions, put
$\Phi=\Hom_H(N,-)$, viewed as a functor to left $C$-modules
by precomposition.

\begin{proof}
\textup{(\ref{general:seed-simple})} To identify $C/CeC$, we first show that $M_0$ is
indecomposable. Indeed, each nonzero direct summand of $M_0$ has a nonzero stable identity, since it is
not projective. A nontrivial decomposition would therefore produce
nontrivial idempotents in $\underline{\End}_H(M_0)=k$.  It follows that
$fCf\simeq \End_H(M_0)^{\mathrm{op}}$ is local. Thus $f$ is primitive.
Since every projective module belongs to $\add(H)$, the ideal $CeC$
consists of endomorphisms factoring through $\add(H)$. Hence
$C/CeC\simeq\underline{\End}_H(M_0)^{\mathrm{op}}\simeq k$.
Thus $eS=0$ and $fS=S=k$.

\textup{(\ref{general:seed-homology})} To compute both families of extension groups, we construct
a projective resolution of $S$ from one of $M_0$.
Choose a minimal projective resolution
\[
 \cdots\longrightarrow P_2\longrightarrow P_1\lraf{d} P_0
 \lraf{p}M_0\longrightarrow 0.
\]
Yoneda identifies $\Phi(H)=Ce$, $\Phi(M_0)=Cf$ and $\Phi(N)=C$.
More generally, Yoneda gives the natural isomorphism
$\Hom_H(U,V)\simeq\Hom_C(\Phi(U),\Phi(V))$ for $U\in\add(N)$ and any module $V$.
In particular, $\Phi(P_j)\in\add(Ce)$ for every $j$.

Since $p$ is surjective with projective source, the image of
$\Phi(p)$ consists of the morphisms factoring through projectives.
Thus $\operatorname{coker}(\Phi(p))\simeq\underline{\Hom}_H(N,M_0)=S$.

For exactness at the remaining terms, it suffices to prove
$\Ext_H^1(N,\Omega_H^{j+1}M_0)=0$ for $j\ge1$.
Dimension shifting over the self-injective algebra $H$ gives
\[
 \underline{\Hom}_H(\Omega_H^{j+1}M_0,M_0)
 \simeq\Ext_H^{j+1}(M_0,M_0)=0.
\]
Auslander--Reiten duality (see the \hyperref[prelim:AR-duality]{discussion in Subsection~\ref*{prelim:stable}}) and
$\tau_HM_0\simeq M_0$ therefore give
\[
\Ext_H^1(N,\Omega_H^{j+1}M_0)\simeq  \Ext_H^1(M_0,\Omega_H^{j+1}M_0)
 \simeq D\underline{\Hom}_H(\Omega_H^{j+1}M_0,M_0)=0.
\]
Thus $\Phi$ preserves exactness of the short exact sequence
\[
 0\longrightarrow\Omega_H^{j+1}M_0\longrightarrow P_j
 \longrightarrow\Omega_H^jM_0\longrightarrow0
 \qquad(j\ge1).
\]
Left exactness identifies the kernel at $\Phi(P_0)$, and the preceding description of the cokernel identifies the augmentation.  It follows that
\begin{equation*}
 \cdots\longrightarrow\Phi(P_2)\longrightarrow\Phi(P_1)\longrightarrow\Phi(P_0)
 \longrightarrow Cf\longrightarrow S\longrightarrow0
\end{equation*}
is a projective resolution of $S$, with degree-zero term $Cf$ and
all negative cochain terms in $\add(Ce)$.

Since $eS=0$ and $fS=k$, applying $\Hom_C(-,S)$ leaves only the term
$k$ in degree zero. To compute the extensions into $C$, we work
with the Hom complex over $H$. Auslander--Reiten duality, together with
$\tau_HM_0\simeq M_0$ and $\underline{\End}_H(M_0)\simeq k$, gives
$\Ext_H^1(M_0,M_0)\simeq k$. Yoneda identifies the complex obtained by
applying $\Hom_C(-,C)$, up to the signs in our Hom convention, with
\begin{equation*}
 \Hom_H(M_0,N)\longrightarrow\Hom_H(P_0,N)
 \longrightarrow\Hom_H(P_1,N)\longrightarrow\cdots.
\end{equation*}
The degree-zero and degree-one cohomology groups vanish by exactness of
the projective presentation of $M_0$, and the cohomology in degree $i\ge2$ is
$\Ext_H^{i-1}(M_0,N)$.  The $H$-summand contributes no positive
extensions because $H$ is injective; the $M_0$-summand contributes
only in degree two by the assumed vanishing of higher self-extensions
and $\Ext_H^1(M_0,M_0)=k$.  Finally, the natural isomorphisms
\[
 \Ext_H^1(M_0,N)
 \simeq D\underline{\Hom}_H(N,\tau_HM_0)
 \simeq D\underline{\Hom}_H(N,M_0)=DS
\]
are right $C$-linear by naturality of Auslander--Reiten duality in
$N$. This proves~\textup{(\ref{general:seed-homology})}.

\textup{(\ref{general:seed-dimension})} We find the injective dimensions from the projective
dimensions of $DC$ on both sides. Put $\Psi=\Hom_H(-,N)$, viewed as a contravariant functor to
right $C$-modules by postcomposition. Yoneda identifies $\Psi(H)=eC$ and
$\Psi(M_0)=fC$, and $\Psi(N)=C$. Projective Hom duality
gives $D\Psi(Q)\simeq\Phi(\nu_HQ)$ for every projective module $Q$.
In particular, $D(eC)\simeq Ce$ as left modules because
$\nu_HH=DH\simeq H$, and dualization gives $D(Ce)\simeq eC$ as
right modules.  Thus ${_C}Ce$ and $eC_C$ are projective-injective.
Applying $\Psi$ to the minimal projective presentation of $M_0$ and then dualizing,
we obtain an exact sequence
\[
 D\Psi(P_1)\lraf{D\Psi(d)}D\Psi(P_0)
 \longrightarrow D\Psi(M_0)\longrightarrow0.
\]
Under projective Hom duality, its first map is $\Phi(\nu_Hd)$.
By the definition of $\tau_H=D\operatorname{Tr}_H$ in
Subsection~\ref{prelim:stable}, we have
$\ker(\nu_Hd)\simeq\tau_HM_0$.  Left exactness of $\Phi$ therefore
identifies the kernel of the first map with $\Phi(\tau_HM_0)\simeq Cf$.
Consequently there is an exact sequence
\begin{equation*}
 0\longrightarrow Cf\longrightarrow U_1\longrightarrow U_0
 \longrightarrow D(fC)\longrightarrow0,
 \qquad U_i=\Phi(\nu_HP_i)\in\add(_CCe).
\end{equation*}
Since $D(Ce)\simeq eC$, dualizing this sequence gives a projective
resolution on the right:
\begin{equation*}
 0\longrightarrow fC\longrightarrow DU_0\longrightarrow DU_1
 \longrightarrow D(Cf)\longrightarrow0,
 \qquad DU_i\in\add(eC_C).
\end{equation*}
The displayed resolutions give projective dimensions at most two.
Since $eS=0$ and $\Hom_C(Cf,S)\simeq k$, applying $\Hom_C(-,S)$
to the first resolution gives $\Ext_C^2(D(fC),S)\simeq k$.
Similarly, applying $\Hom_{C^{\mathrm{op}}}(-,DS)$ to the right
resolution gives $\Ext_{C^{\mathrm{op}}}^2(D(Cf),DS)\simeq k$.
Using $DC=D(eC)\oplus D(fC)$ on the left and
$DC=D(Ce)\oplus D(Cf)$ on the right, together with the projectivity
of $D(eC)\simeq Ce$ and $D(Ce)\simeq eC$, we obtain
\[
 \Ext_C^2(DC,S)\simeq k,\qquad
 \Ext_{C^{\mathrm{op}}}^2(DC,DS)\simeq k.
\]
Thus $\pd_C(DC)=2=\pd_{C^{\mathrm{op}}}(DC)$.
Finite-dimensional duality gives
$\id_C(C)=2=\id_{C^{\mathrm{op}}}(C)$,
proving~\textup{(\ref{general:seed-dimension})}.

\textup{(\ref{general:seed-basic})} Suppose that $H$ is connected and split basic with
$r$ simple modules.  The regular module is the direct sum of one
representative of each indecomposable projective.  These $r$ summands and $M_0$ are
pairwise nonisomorphic, and their endomorphism division rings are
all $k$.  The radical description for the endomorphism algebra of
their direct sum therefore gives
$C/\rad(C)\simeq k^{r+1}$.
A block decomposition of $C$ would partition the indecomposable
summands of $N$ into classes with no morphisms between them.
Connectedness of $H$ puts its projective summands in one class, and
$\Hom_H(H,M_0)=M_0\ne0$ places $M_0$ in the same class. Thus $C$
is connected.
\end{proof}

\begin{Rem}
By Proposition \ref{general:seed}\textup{(\ref{general:seed-dimension})} and its proof, $\operatorname{domdim}(C)=2=\id_C(C)$ and $\operatorname{domdim}(C^{\mathrm{op}})=2=\id_{C^{\mathrm{op}}}(C)$. Consequently $C$ is $1$-minimal Auslander--Gorenstein in the sense of \cite[Definition~1.1]{IS}, meaning that $\operatorname{domdim}(C)\ge2\ge\id_C(C)$. Since $H$ is Frobenius, $DH\simeq H$ as left modules.
Duality gives $C\simeq\End_{H^{\mathrm{op}}}(DN)$, where $DN$ is a generator over the
self-injective algebra $H^{\mathrm{op}}$, that is, $H^{\mathrm{op}}\in\add(DN)$.
Hence $C$ is a Morita algebra.
\end{Rem}

\subsection{The quantum exterior seed}
\label{seed:concrete-pair}

Fix an element $q\in k^\times$ of infinite multiplicative order, and set
\begin{equation*}
 H=k\langle x,y\rangle/(x^2,y^2,xy+qyx),\qquad z=yx.
\end{equation*}
For $a\in k^\times$, put $M(a)=H/H(x+ay)$ and $M_0=M(1)$.
Write $u_a,v_a$ for the classes of $1,x$ in $M(a)$, respectively;
the $H$-action is given by
\begin{equation*}
 xu_a=v_a,\qquad yu_a=-a^{-1}v_a,\qquad
 xv_a=yv_a=zu_a=zv_a=0.
\end{equation*}
These modules form the two-dimensional
family in \cite[Section~4.1]{Erd}, with parameter $-a^{-1}$
in that notation.
The following lemma records the seed conditions and the explicit
resolution needed for the signed comparison.

\begin{Lem}\label{seed:quantum-pair}
\leavevmode
\begin{enumerate}
\item
The algebra $H$ is a four-dimensional local Frobenius algebra, with
basis $1,x,y,z$ and Frobenius functional $\lambda:H\to k$ given by
the coefficient of $z$.

\item\label{seed:pair-modules}
For every $a\in k^\times$, $M(a)$ is a two-dimensional nonprojective
indecomposable module, and
\[
 \End_H(M(a))=k[\varepsilon_a]/(\varepsilon_a^2),\qquad
 \varepsilon_a(u_a)=v_a,\quad\varepsilon_a(v_a)=0.
\]
Moreover,
\[
 \tau_HM(a)\simeq M(a),\quad \underline{\End}_H(M(a))\simeq k,\quad
 \Ext_H^i(M(a),M(a))\simeq
 \begin{cases}k&i=1,\\0&i\ge2\end{cases}
 \quad(i\ge1).
\]
In particular, $(H,M_0)$ satisfies the hypotheses of
Proposition~\ref{general:seed}.

\item\label{seed:pair-resolution}
A minimal projective resolution $P(a)^\bullet$ of $M(a)$ has
$P_n(a)=P(a)^{-n}=H$ for $n\ge0$ and the augmented form
\begin{equation*}
 \begin{gathered}
 \cdots\longrightarrow P_2(a)\xrightarrow{\partial_2(a)}P_1(a)
 \xrightarrow{\partial_1(a)}P_0(a)\xrightarrow{p_a}M(a)
 \longrightarrow0,\\
 \partial_n(a)=\rho_{x+q^{n-1}ay},\quad p_a(1)=u_a.
 \end{gathered}
\end{equation*}
Here $\rho_b$ denotes right multiplication by $b$ and
$\partial_n(a)=d_{P(a)}^{-n}$ for $n\ge1$.
A generator of $\Ext_H^1(M(a),M(a))$ is represented by the map
$P_1(a)\to M(a)$ sending $1$ to $v_a$.
\end{enumerate}
\end{Lem}
\begin{proof}
The three quadratic relations are linearly independent, so the
degree-two component is one-dimensional with basis $z$.
Every word of length three vanishes, and hence $1,x,y,z$ form a basis.
The ideal $(x,y)$ is nilpotent and its quotient is $k$, so $H$ is
local. On this basis, the only nonzero pairings are
\[
 \lambda(1\cdot z)=\lambda(z\cdot1)=\lambda(y\cdot x)=1,
 \qquad \lambda(x\cdot y)=-q.
\]
Since $q\ne0$, the pairing is nondegenerate, proving (1).

The ideal $H(x+ay)$ is $k(x+ay)\oplus kz$. Thus $u_a,v_a$ form a
basis, with the stated $H$-actions. An endomorphism must commute with the actions of $x,y$, so it has the
form $c\,1_{M(a)}+d\varepsilon_a$ for $c, d\in k$. This endomorphism algebra is local,
so $M(a)$ is indecomposable.  Every projective $H$-module is free;
thus the dimensions $\dim_k(M(a))=2$ and $\dim_k(H)=4$ exclude
projectivity.

To compute the self-extensions of $M(a)$, we first verify its stated
projective resolution. For $b\in k^\times$, right multiplication $H\to H$ by
$x+by$ is given by
$$h=h_0+h_xx+h_yy+h_zz\mapsto h_0(x+by)+(-qb h_x+h_y)z,$$
where $h_0, h_x, h_y, h_z\in k$. Its image is $H(x+by)$ and its kernel is $H(x+qby)$.
These equalities prove exactness of the resolution
in~\textup{(\ref{seed:pair-resolution})}.
The kernel of $p_a$ lies in $\rad(H)$, and every differential has
image in $\rad(H)$, so the resolution is minimal.
We now apply $\Hom_H(-,M(a))$ to this resolution.
Evaluation at $1$ identifies $\Hom_H(P_n(a),M(a))$ with $M(a)$.
The differential from degree $n$ to degree $n+1$ is
$(-1)^{n+1}$ times the action of $x+q^n a y$ on $M(a)$, given by
$u_a\mapsto(1-q^n)v_a$ and $v_a\mapsto 0$.
It vanishes for $n=0$ and has kernel and image $kv_a$ for every
$n\ge1$, because $q$ has infinite order.  This proves the asserted
Ext groups and their degree-one representative.

To prove $\tau_HM(a)\simeq M(a)$, we compute $D\operatorname{Tr}_H(M(a))$.
Under evaluation at $1$, precomposition with $\rho_{x+ay}$ on
$\Hom_H(H,H)$ is left multiplication by $x+ay$.
Thus the minimal presentation in~\textup{(\ref{seed:pair-resolution})} gives
$\operatorname{Tr}_H(M(a))=H/(x+ay)H$.
Here $(x+ay)H=k(x+ay)\oplus kz$, and the right quotient has basis
$\bar1,\bar x$ with $\bar1x=\bar x$,
$\bar1y=-a^{-1}\bar x$ and $\bar x x=\bar x y=0$.
Let $\alpha,\beta$ be the $k$-dual basis of $\bar1,\bar x$ satisfying $\alpha(\bar1)=1=
\beta(\bar x)$ and $\alpha(\bar x)=0=\beta(\bar 1)$.
Then $x\beta=\alpha$, $y\beta=-a^{-1}\alpha$ and $x\alpha=y\alpha=0$.
Hence $u_a\mapsto\beta$, $v_a\mapsto\alpha$ identifies $M(a)$ with
$D\operatorname{Tr}_H(M(a))=\tau_HM(a)$.
Auslander--Reiten duality, together with $\tau_HM(a)\simeq M(a)$,
shows that $\underline{\End}_H(M(a))$ has dimension one.  Its identity is
nonzero because $M(a)$ is nonprojective, so it is isomorphic to $k$
as an algebra.  This completes \textup{(\ref{seed:pair-modules})} and \textup{(\ref{seed:pair-resolution})}.
\end{proof}

To simplify the notation for the endomorphism seed and its resolution, we write
$$P_n=P_n(1)=H,\;\; \partial_n=\partial_n(1)=\rho_{x+q^{n-1}y}\quad(n\geq 1),\;\; u=u_1,\;\;  v=v_1, \;\;  p=p_1:\; 1\mapsto u.$$
Let $N,C,e,f,S,\chi_C$ be the data obtained from $(H,M_0)$ by Proposition~\ref{general:seed}, and put
$\Phi=\Hom_H(N,-)$. Right multiplication identifies $eCe$ with $H$ as algebras, and
evaluation at $1$ identifies $eCf$ with $M_0$ as left $H$-modules.
Moreover, $fCf\simeq k[t]/(t^2)$, where $t=\varepsilon_1$. The proposition shows that $C$ is connected and split basic with primitive idempotents $e,f$ and $C/\rad(C)\simeq k\times k$.
Frobenius duality gives $\dim_k(\Hom_H(M_0,H))=2$, and therefore
$\dim_k(C)=10$.

Specializing the projective resolution constructed in the proof of
Proposition~\ref{general:seed}(\ref{general:seed-homology}) gives a minimal
projective resolution $P_S^\bullet$ of $S$:
\begin{equation*}
 \begin{gathered}
 \cdots\longrightarrow\Phi(P_2)\longrightarrow\Phi(P_1)\longrightarrow\Phi(P_0)
 \xrightarrow{\Phi(p)}Cf\longrightarrow S\longrightarrow0,\\
 P_S^0=Cf,\qquad P_S^{-n-1}=Ce\quad(n\ge0).
 \end{gathered}
\end{equation*}
By the Hom-complex identification in the proof of
Proposition~\ref{general:seed}\textup{(\ref{general:seed-homology})}, a generator of
$H^2\Hom_C^\bullet(P_S^\bullet,C)\simeq DS$ is represented by
$a:P_1\to M_0\subseteq N$ with $a(1)=v$, as specified in
Lemma~\ref{seed:quantum-pair}\textup{(\ref{seed:pair-resolution})}.

\subsection{The fixed signed auxiliary seed}
We equip the concrete seed with an involution whose comparison fixes the degree-two extension
class.

Define
$$\delta:\;\; H\lra H,\quad  x\mapsto -x,\;\; y\mapsto -y,\;\; z\mapsto z.$$
Then $\delta$ is an algebra automorphism satisfying $\delta^2=1$. Twisting the source gives an isomorphism of $H$-modules:
$$\kappa_{M_0}:\;\;{}_\delta M_0\lraf{\sim}M_0,\quad u\mapsto u, \;\; v\mapsto-v.$$
Together with $\delta:{}_\delta H\xrightarrow{\sim}H$,
it gives an isomorphism of $H$-modules:
$$\kappa_N=\delta\oplus\kappa_{M_0}:\;\;{}_\delta N\lraf{\sim}N,\qquad \kappa_N\circ{}_\delta\kappa_N=1_N.$$
Conjugation of twisted endomorphisms by $\kappa_N$ defines the algebra
automorphism
$$\gamma:C\longrightarrow C, \;\; c\mapsto \kappa_N\circ{}_\delta c\circ\kappa_N^{-1},$$
where $c$ is viewed in $\End_H(N)$ before passing to the opposite algebra.
The identity $\kappa_N\circ{}_\delta\kappa_N=1_N$ makes $\gamma$ involutive.

For twisting, we need the algebra involution together with a compatible
comparison on the resolution of $S$. The next lemma records this
comparison and its effect on the degree-two extension class.

\Needspace{4\baselineskip}
\begin{Lem}\label{seed:gamma}
\leavevmode
\begin{enumerate}
\item\label{seed:gamma-involution} The automorphism $\gamma$ is involutive, fixes $e,f$, and
preserves $\chi_C$. On $fCf=k[t]/(t^2)$, it satisfies $\gamma(t)=-t$.

\item\label{seed:gamma-comparison} There is a $C$-linear cochain isomorphism
$\kappa_S=(\kappa_S^j)_{j\leq 0}: {}_\gamma P_S^\bullet\xrightarrow{\sim}P_S^\bullet$
that induces $1_S$, using ${}_\gamma S=S$ through the character identification.


\item\label{seed:gamma-cohomology} Let $\gamma:{}_\gamma C\xrightarrow{\sim}C$ denote the
$C$-linear isomorphism given by the algebra automorphism on
underlying elements. Define a $k$-linear cochain automorphism
$\mathcal F_\gamma$ of
$\Hom_C^\bullet(P_S^\bullet,C)$ degreewise by
\[
 \mathcal F_\gamma^a(\xi)
   =\gamma\circ{}_\gamma\xi\circ(\kappa_S^{-a})^{-1},
 \qquad \xi\in\Hom_C(P_S^{-a},C).
\]
It satisfies
$\mathcal F_\gamma^a(\xi c)=\mathcal F_\gamma^a(\xi)\gamma(c)$
for $c\in C$.

Under the right $C$-module identification
$H^2\Hom_C^\bullet(P_S^\bullet,C)\simeq DS$, its induced map is
$H^2(\mathcal F_\gamma)=1_{DS}$.
\end{enumerate}
\end{Lem}
\begin{proof}
\textup{(\ref{seed:gamma-involution})} Involutivity was established in the preceding
construction. Since $\kappa_N$ preserves the two summands of $N$,
$\gamma$ fixes their projections and preserves $CeC$. Its induced
$k$-algebra automorphism of $C/CeC=k$ is the identity, so
$\chi_C\gamma=\chi_C$ and the character identifies ${}_\gamma S$
with $S$. In $fCf=\End_H(M_0)^{\mathrm{op}}$, $t=\varepsilon_1$
sends $u$ to $v$ and $v$ to zero; see Lemma \ref{seed:quantum-pair}\textup{(\ref{seed:pair-modules})}.
Conjugation by $\kappa_{M_0}$ therefore sends $t$ to $-t$.

\textup{(\ref{seed:gamma-comparison})} To construct $\kappa_S$, we first compare the augmented
$H$-projective resolution of $M_0$ with its $\delta$-twist. On the
resolution in Lemma~\ref{seed:quantum-pair}\textup{(\ref{seed:pair-resolution})}, put
$\kappa_n=(-1)^n\delta:{}_\delta P_n\xrightarrow{\sim}P_n$ for $n\ge0$.
Each $\kappa_n$ is $H$-linear, and
$p\kappa_0=\kappa_{M_0}\circ{}_\delta p$. For $n\ge1$,
\[
 \kappa_{n-1}(h(x+q^{n-1}y))
   =(-1)^n\delta(h)(x+q^{n-1}y)
   =\partial_n\kappa_n(h).
\]
Thus $\kappa_{n-1}\circ{}_\delta\partial_n
=\partial_n\kappa_n$, and these maps identify the twisted
augmented resolution with the original one.
They induce the following maps on $P_S^0=\Phi(M_0)$ and
$P_S^{-n-1}=\Phi(P_n)$ for $n\ge 0$:
\begin{equation*}
 \begin{aligned}
 \kappa_S^0(\varphi)
   &=\kappa_{M_0}\circ{}_\delta\varphi\circ\kappa_N^{-1},\quad\varphi\in  P_S^0;\\
 \kappa_S^{-n-1}(\varphi)
   &=\kappa_n\circ{}_\delta\varphi\circ\kappa_N^{-1},\quad\varphi\in  P_S^{-n-1}.
 \end{aligned}
\end{equation*}
On underlying vectors, the definition of $\gamma$ gives the identities
$$\kappa_S^j(c\cdot\varphi)=\kappa_S^j(\varphi \circ c)=\kappa_S^j(\varphi)\circ\gamma(c)=\gamma(c)\cdot\kappa_S^j(\varphi)$$
for $c\in C$ and $\varphi\in P_S^j$.
Since $\gamma^2=1$, $\kappa_S$ is $C$-linear from the
$\gamma$-twisted source. The comparison identities for $p$ and
$\partial_n$ give compatibility with the differentials, and each
component is invertible in $C\modcat$. The projection $N\to M_0$ is
fixed by these formulas for $\kappa_S$,
and its stable class generates $S=\underline{\Hom}_H(N,M_0)$. The induced
map on the augmentation is therefore $1_S$.

\textup{(\ref{seed:gamma-cohomology})} Twisting preserves the differentials, and the cochain
isomorphisms $\kappa_S$ and $\gamma:{}_\gamma C\to C$ give the stated
$k$-linear cochain automorphism. The right-action identity follows
from $\gamma$ being an algebra automorphism. To determine the induced
action on degree-two cohomology, we compute the image of an explicit generator.
Under the Yoneda identification $\Hom_C(\Phi(Q),C)\simeq\Hom_H(Q,N)$, the induced map
sends $b:Q\to N$ to $\kappa_N\circ{}_\delta b\circ\kappa_Q^{-1}$, where  $Q=M_0$ or $P_n$; the map $\kappa_Q:{}_\delta Q\to Q$ is $\kappa_{M_0}$ or $\kappa_n$, respectively. Since
$P_S^{-2}=\Phi(P_1)$, the degree-two class is represented by
$a:P_1\to M_0\subset N$, $a(1)=v$, by
Lemma~\ref{seed:quantum-pair}\,\textup{(\ref{seed:pair-resolution})}
and the Hom-complex identification in the proof of
Proposition~\ref{general:seed}\textup{(\ref{general:seed-homology})}.
Since $\kappa_1=-\delta$, the image of this representative satisfies
\[
 (\kappa_N\circ{}_\delta a\circ\kappa_1^{-1})(1)
       =-\kappa_{M_0}(v)=v.
\]
Thus the generator of the one-dimensional degree-two cohomology
is fixed. The right $C$-action on $DS$ factors through $\chi_C$, and
$\chi_C\gamma=\chi_C$, so the induced map is right $C$-linear
and equals $1_{DS}$. The two signs cancel in every characteristic;
in characteristic two they are both $1$.
\end{proof}

We distinguish the fixed auxiliary factor from the variable Frobenius
pair before forming their tensor product.

\par\medskip\refstepcounter{Theo}\noindent\textbf{Notation \theTheo.}\label{seed:auxiliary-notation}
Denote the concrete data by
\[
 \begin{aligned}
 (H_*,M_*,C_*,S_*,e_*,f_*)&=(H,M_0,C,S,e,f),\\
 (\gamma_*,\chi_*)&=(\gamma,\chi_C).
 \end{aligned}
\]
The starred objects remain fixed throughout the subsequent
construction.  Unstarred $H,M_0,C,S$ denote a variable Frobenius
pair and its associated homological seed.
\par\medskip

There is an algebra isomorphism from $C_*$ to the ten-dimensional
algebra $C$ of \cite[Section~3]{F} that identifies $S_*$ with the
simple module $s$ and intertwines $\gamma_*$ with the involutive
automorphism $\gamma$ defined there.

\section{Derived induction and its bimodule model}
\label{ind:section}

We need a finite bimodule whose evaluation can be computed from
derived induction. Proposition~\ref{ind:one-generator} gives the
Yoneda algebra, and Proposition~\ref{ind:Y-evaluation} supplies
the required bimodule.
Each seed contributes a derived Nakayama shift of two, so the tensor
product gives a shift of four and the induction triangle has a
connecting class of degree five. Stable duality and the shift by
$[-1]$ in the bimodule model then yield the stable Hom support
$\{-3,0\}$ required by Proposition~\ref{frame:factory}.

Let $(C_*,S_*,\gamma_*,\chi_*)$ be the fixed auxiliary seed of
Notation~\ref{seed:auxiliary-notation}, and let $(C,S)$ be the seed
associated with a Frobenius pair satisfying the hypotheses of
Proposition~\ref{general:seed}. By
Lemma~\ref{seed:gamma}\,\textup{(\ref{seed:gamma-cohomology})}, the fixed comparison for
$\gamma_*$ induces the identity on $\Ext_{C_*}^2(S_*,C_*)$.

Set
\[
 B=C_*\otimes_k C,\qquad X=S_*\otimes_k S,
 \qquad E=B\ltimes DB.
\]
Both seed modules are one-dimensional with characters $\chi_*:C_*\to k=S_*$ and $\chi_C: C\to k=S$; see Proposition \ref{general:seed}\textup{(\ref{general:seed-simple})}. They determine characters
$$\chi_B=\chi_*\otimes\chi_C:\;\;B\lra k\quad\mbox{and}\quad
\chi:\;\;E\lra k,\;\; (b,\varphi)\mapsto \chi_B(b).$$
We regard $X=k_{\chi_B}$ as an $E$-module through this quotient; in particular,
$DB$ annihilates $X$. By the trivial-extension construction in
Subsection~\ref{prelim:notation}, $E$ is symmetric and
$\dim_k(E)=20\dim_k(C)$.

\paragraph{The distinguished primitive idempotent.}
Let $e,f\in C$ be the idempotents in Proposition~\ref{general:seed}, and put
$p=f_*\otimes f\in B\subseteq E$. The local algebra $fCf$ has residue field $k$, while
$f_*C_*f_*\simeq k[t]/(t^2)$ also has residue field $k$.
Therefore $pBp=(f_*C_*f_*)\otimes_k(fCf)$ is local with residue field $k$.
Since $DB\subseteq\rad(E)$, the algebra $pEp$ is local as well.
It follows that $p$ is primitive in both $B$ and $E$. Consequently,
the character $\chi$ induces the projective cover $Ep\to X$, with
kernel $\rad(E)p$.
This argument applies whether or not the variable seed is split basic.

\subsection{Tensor induction and extension groups}

We first record the extension groups and one-sided homological
dimensions that control induction from $B$ to $E$.

\begin{Lem}\label{ind:tensor-input}
\leavevmode
\begin{enumerate}
\item\label{ind:tensor-input-ext} There are isomorphisms
\begin{equation*}
 \Ext_B^i(X,X)\simeq\begin{cases}k&i=0,\\0&i>0\end{cases},\qquad
 \Ext_B^i(X,B)\simeq\begin{cases}DX&i=4,\\0&i\ne4\end{cases}
 \qquad(i\ge0).
\end{equation*}
The second is an isomorphism of right $B$-modules.
\item\label{ind:tensor-input-dimensions} The one-sided dimensions satisfy
\begin{equation*}
 \pd_B(DB)=\pd_{B^{\mathrm{op}}}(DB)
 =\pd_B(E)=\pd_{B^{\mathrm{op}}}(E)=4,
 \qquad \id_B(B)=\id_{B^{\mathrm{op}}}(B)=4.
\end{equation*}
\end{enumerate}
\end{Lem}
\begin{proof}
\textup{(\ref{ind:tensor-input-ext})} Since $B$ and $X$ are tensor products, we compute the
required extension groups from those of the two seeds. The natural
K\"unneth isomorphism for Ext (see \cite[Theorem~XI.3.1]{CE}) is
\[
 \Ext_{A_1\otimes_k A_2}^n(M_1\otimes_k M_2,N_1\otimes_k N_2)
 \simeq\bigoplus_{i+j=n}
 \Ext_{A_1}^i(M_1,N_1)\otimes_k\Ext_{A_2}^j(M_2,N_2).
\]
The formula applies to the finite-dimensional algebras and modules
used here. Applying it to the two seeds gives the isomorphisms in
\textup{(\ref{ind:tensor-input-ext})}.
Naturality in the target modules preserves the right $B$-action
in the second identification, and
$DS_*\otimes_k DS\simeq D(S_*\otimes_k S)=DX$.

\textup{(\ref{ind:tensor-input-dimensions})} To obtain the exact dimensions, we need length-four
resolutions and nonvanishing in degree four. Tensoring the length-two
resolutions of $DC_*$ and $DC$ gives
length-four resolutions of $DB$ on both sides, so its two
projective dimensions are at most four. The same K\"unneth
isomorphism gives
\[
 \begin{aligned}
 \Ext_B^4(DB,X)&\simeq
  \Ext_{C_*}^2(DC_*,S_*)\otimes_k\Ext_C^2(DC,S)\simeq k,\\
 \Ext_{B^{\mathrm{op}}}^4(DB,DX)&\simeq
  \Ext_{C_*^{\mathrm{op}}}^2(DC_*,DS_*)\otimes_k
  \Ext_{C^{\mathrm{op}}}^2(DC,DS)\simeq k.
 \end{aligned}
\]
The degree-two groups are nonzero by the dual regular resolutions
in Proposition~\ref{general:seed}, applied to both seeds. Thus both
projective dimensions are exactly four. Since $E=B\oplus DB$ upon
restriction to $B$, the same equalities hold for $E$. Finite-dimensional
duality gives the asserted injective dimensions of $B$.
\end{proof}

We equip $\Ext_E^*(X,X)=\bigoplus_{i\ge0}\Ext_E^i(X,X)$ with its
Yoneda product. For homogeneous classes represented by
$\alpha:X\to X[r]$ and $\beta:X\to X[t]$, the product is
$\alpha\beta:=\alpha[t]\circ\beta:X\to X[r+t]$.
Put
$$J=E\otimes_B^{\mathbf L}X,\qquad L=DB\otimes_B^{\mathbf L}X.$$
These are objects of $\mathscr D^b(E)$ by
Lemma~\ref{ind:tensor-input}\textup{(\ref{ind:tensor-input-dimensions})}; the left $E$-action on $L$ factors through
$E\twoheadrightarrow B$.
We first identify $L$ with a shift of $X$ and use the resulting
triangle to compute the Yoneda algebra over $E$.

\begin{Prop}\label{ind:one-generator}
There is a triangle in $\mathscr D^b(E)$
\begin{equation*}
 X[4]\longrightarrow J\longrightarrow X
   \xrightarrow{\varpi}X[5].
\end{equation*}
Its connecting class generates the graded associative algebra
\begin{equation*}
 \Ext_E^*(X,X)\simeq k[\varpi].
\end{equation*}
Here $\varpi$ has degree five, and $k[\varpi]$ is the polynomial algebra
with this cohomological grading.
\end{Prop}
\begin{proof}
The triangle is obtained by tensoring the exact sequence of
$E$--$B$-bimodules
\[
 0\longrightarrow DB\longrightarrow E\longrightarrow B
 \longrightarrow0,
\]
where the left $E$-actions on $DB$ and $B$ factor through
$E\twoheadrightarrow B$. Applying $-\otimes_B^{\mathbf L}X$ and
using $B\otimes_B^{\mathbf L}X\simeq X$ gives the triangle in $\mathscr D^b(E)$:
$$L\lra J\lra X\lra L[1].$$ To obtain the asserted form, we must
identify $L$ with $X[4]$. Derived Nakayama and Ext--Tor duality give
\[
 L\simeq D\mathbf R\Hom_B(X,B),\qquad
 H^{-i}(L)\simeq D\Ext_B^i(X,B)
 \quad(i\ge0);
\]
see \cite[Lemma~9.72]{Rotman} for the cohomology identification.
By Lemma~\ref{ind:tensor-input}\textup{(\ref{ind:tensor-input-ext})}, the only nonzero cohomology is
$H^{-4}(L)\simeq DDX\simeq X$. All these identifications respect
the left $B$-action, and the $E$-actions factor through $B$.
Thus $L\simeq X[4]$ in $\mathscr D^b(E)$ via this cohomology
identification. This gives the stated induction triangle. Its last
morphism defines $\varpi$ using the convention for connecting morphisms
fixed in the preliminaries.

To compute the Yoneda algebra from this triangle, we first determine
$\Ext_{\mathscr D^b(E)}^i(J,X)$. Derived induction--restriction adjunction
\cite[Proposition~4.3.15]{KrauseHTR} and Lemma~\ref{ind:tensor-input}\textup{(\ref{ind:tensor-input-ext})} give, for every
$i\in\mathbb Z$,
\[
 Q^i=\Ext_{\mathscr D^b(E)}^i(J,X)
 \simeq\Ext_{\mathscr D^b(B)}^i(X,X)
 \simeq\begin{cases}k&i=0,\\0&i\ne0\end{cases}.
\]
No flatness assumption on $E_B$ is needed, since the tensor
product is derived. Applying $\Hom_{\mathscr D^b(E)}(-,X[i])$ to
the induction triangle yields the exact sequence
\[
 Q^{i-1}\longrightarrow\Ext_E^{i-5}(X,X)
 \xrightarrow{\alpha\mapsto\alpha[5]\circ\varpi}
 \Ext_E^i(X,X)\longrightarrow Q^i.
\]
Here negative ordinary Ext groups vanish because $X$ is concentrated
in degree zero. For $i=1$, both $\Ext_E^{-4}(X,X)$ and $Q^1$ vanish,
so $\Ext_E^1(X,X)=0$. For $i\ge2$, the two groups $Q^{i-1}$ and
$Q^i$ vanish, and the middle arrow is an isomorphism. This proves
vanishing in degrees $2,3,4$. In degree five, it sends $1_X$ to the
nonzero class $\varpi$. Iterating the same isomorphism shows that
$\varpi^m$ is a basis of $\Ext_E^{5m}(X,X)$ for every $m\ge0$, while
all other nonnegative degrees vanish. Thus the homomorphism
$k[T]\to\Ext_E^*(X,X)$, $T\mapsto\varpi$, is an isomorphism of
graded algebras.
\end{proof}

For the remainder of the construction, we put
$H^a=\underline{\Hom}_E(X,X[a])$ for $a\in\mathbb Z$.
We form the stable image $\overline J=\Phi_E(J)$ using the triangle functor defined in
Subsection~\ref{prelim:stable}. We also write $\varpi$ for the stable connecting
class obtained through the specified module and shift comparisons.
The cone criterion requires the support of stable morphism groups in
all integer degrees. We also compute the groups involving $\overline J$,
which will give the stable morphism spaces of the evaluated bimodule
model below.

\begin{Prop}\label{ind:Tate-and-middle}
\leavevmode
\begin{enumerate}
\item\label{ind:Tate-self} The module ${_E}X$ is nonprojective, and
\begin{equation*}
 H^{5m}=k\varpi^m,\qquad H^{-5m-1}\simeq D(k\varpi^m)\quad(m\ge0),
 \qquad H^a=0\quad\text{otherwise}.
\end{equation*}
\item\label{ind:Tate-middle} The stable image $\overline J$ satisfies
\begin{equation*}
 \underline{\Hom}_E(X,\overline J[a])\simeq
 \begin{cases}k&a=-4,-1,\\0&\text{otherwise}.
 \end{cases}
\end{equation*}
\end{enumerate}
\end{Prop}
\begin{proof}
\textup{(\ref{ind:Tate-self})} We first determine all the groups $H^a$. The nonzero
fifth self-extension proves that $X$ is nonprojective.
Since $\End_E(X)=k$, its identity has a nonzero stable class and
$H^0=k$. Positive stable Hom agrees with ordinary Ext, so
Proposition~\ref{ind:one-generator} determines $H^a$ for $a\ge0$.
Lemma~\ref{tool:duality} gives $H^{-i-1}\simeq DH^i$ for $i\ge0$,
which determines $H^a$ for $a<0$ and proves \textup{(\ref{ind:Tate-self})}.

\textup{(\ref{ind:Tate-middle})} The induction triangle expresses the required groups in
terms of the kernels and cokernels of multiplication by $\varpi$ on
$H^*$. We first determine these multiplication maps.
For $r\in\mathbb Z$, define $\mu_r:H^r\to H^{r+5}$ by
$\mu_r(\beta)=\varpi[r]\circ\beta$, where $\beta:X\to X[r]$ and
$\varpi[r]:X[r]\to X[r+5]$.
For $r\ge0$, these maps are isomorphisms by
Proposition~\ref{ind:one-generator}.

To determine the negative-degree maps, we use the composition formula
in symmetric stable duality. Let $\lambda=\lambda_X:H^{-1}\to k$
be the linear form of Lemma~\ref{tool:duality}. For $m\ge0$, write
$\vartheta_m=\vartheta^{-5m-1}_{X,X}$. Its composition formula
gives the duality isomorphism
$$
H^{5m}\xrightarrow{\sim}D H^{-5m-1},\quad \rho'\mapsto [\vartheta_m(\rho'): \;\beta'\mapsto \lambda\bigl(\rho'[-5m-1]\circ\beta'\bigr)]
$$
for $\rho': X\to X[5m]$ and $\beta': X\to X[-5m-1]$. We consider $r=-5m-1$ for $m\ge 1$, $\beta\in H^r$, and $\rho\in H^{5(m-1)}$. Shifted associativity gives
\[
 \vartheta_{m-1}(\rho)\bigl(\mu_r(\beta)\bigr)
 =\lambda\bigl(\rho[r+5]\circ\varpi[r]\circ\beta\bigr)
 =\vartheta_m\bigl(\rho[5]\circ\varpi\bigr)(\beta).
\]
Identifying $H^{5m}$ with $D H^{-5m-1}$, and $H^{5(m-1)}$ with $D H^{-5m+4}$ via these duality isomorphisms, the map $D(\mu_r)$ is the positive multiplication isomorphism
$$H^{5(m-1)}\lra H^{5m},\quad \rho\mapsto\rho[5]\circ\varpi.$$
Thus $\mu_r: H^{-5m-1}\to H^{-5m+4}$ is an isomorphism for $m\ge1$. Since $H^4=0$, the map $\mu_{-1}:H^{-1}\to H^4$ is zero.

To compute the required groups, we apply $\Phi_E$ to the triangle of
Proposition~\ref{ind:one-generator} and obtain a triangle
$X[4]\to \overline{J}\to X \xrightarrow{\varpi}X[5]$ in $E\text{-}\underline{\mathrm{mod}}$.
Its stable Hom sequence has the exact segment
\[
 H^{a-1}\xrightarrow{\mu_{a-1}}H^{a+4}
 \longrightarrow\underline{\Hom}_E(X,\overline J[a])
 \longrightarrow H^a\xrightarrow{\mu_a}H^{a+5},
\]
and hence a short exact sequence
\[
 0\longrightarrow\operatorname{coker}(\mu_{a-1})
 \longrightarrow\underline{\Hom}_E(X,\overline J[a])
 \longrightarrow\ker(\mu_a)\longrightarrow0.
\]
For $a\ge1$, both $\mu_a$ and $\mu_{a-1}$ are isomorphisms. For $a=0$,
the cokernel is zero by $H^4=0$, while $\mu_0:H^0\to H^5$
is an isomorphism. Hence
$\underline{\Hom}_E(X,\overline J[a])=0$ for all $a\ge0$.
For $a<0$, the cokernel is nonzero only at $a=-4$, where
$\mu_{-5}:0\to H^0=k$; the kernel is nonzero only at $a=-1$,
where $\mu_{-1}:H^{-1}=k\to H^4=0$.
Both exceptional terms are $k$, proving \textup{(\ref{ind:Tate-middle})}.
\end{proof}

\subsection{A finite-dimensional bimodule model}
\label{ind:finite-Y}

To apply Proposition \ref{frame:factory} in the framework of Proposition \ref{ind:Tate-and-middle}, we must represent $\overline J[-1]$ by the tensor product $Y\otimes_EX$  for a finite-dimensional, separately projective $E$-bimodule $Y$.

Let $\mathcal K=E\otimes_B^{\mathbf L}E\in\mathscr D^b(E^e)$.
Choose a bounded-above complex $K^\bullet$ of projective $E^e$-modules
representing $\mathcal K$, with $K^j=0$ for $j>0$.
Such a model is obtained by resolving the $(B,E)$-bimodule $E$ by
projective modules over $B\otimes_kE^{\mathrm{op}}$ and applying
$E\otimes_B(-)$. Its terms are projective $E^e$-modules, and the
underlying left $B$-projective resolution computes the derived
tensor product. Any model with the stated properties may be used.
Choose also a cochain map $m_K:K^\bullet\to E$ representing
derived multiplication $\mu:E\otimes_B^{\mathbf L}E\to E$.
The algebra extension $B^e\hookrightarrow E^e$ gives the adjoint pair $\big(E^e\otimes_{B^e}^{\mathbf L}-, \operatorname{Res}_{B^e}\big)$ of the derived induction and restriction functors, where $E^e\otimes_{B^e}^{\mathbf L}-\simeq E\otimesL_B-\otimesL_BE$. The morphism $\mu$ then corresponds to the inclusion $B\hookrightarrow\operatorname{Res}_{B^e}(E)$, or equivalently, $\mu$ is obtained by inducing this inclusion and then applying the counit. Once $E$ is fixed, the model $K^\bullet$ and the representative $m_K$ are chosen independently of the parameter $h$ introduced in Section~\ref{lift:construction-section}.

Set
\[
 \overline K=\Phi_{E^e}(K^\bullet),\qquad
 m=\Phi_{E^e}(m_K):\overline K\longrightarrow E.
\]
The next proposition represents $\overline K[-1]$ by a finite-dimensional, separately projective $E$-bimodule. Its evaluation at $X$ then supplies the stable morphism spaces needed for the lifting construction.

\begin{Prop}\label{ind:Y-evaluation}
\leavevmode
\begin{enumerate}
\item\label{ind:model-cohomology} The cohomology of $K^\bullet$ is concentrated in degrees $i$ with
$-4\leq i\leq 0$.

\item\label{ind:model-bimodule} There is a finite-dimensional, separately projective $E$-bimodule $Y$, with a specified isomorphism
$Y\simeq\overline K[-1]$ in
$E^e\text{-}\underline{\mathrm{mod}}$.

\item\label{ind:model-evaluation} Evaluation at $X$ gives a compatible isomorphism
$Y\otimes_E X\simeq\overline J[-1]$ in
$E\text{-}\underline{\mathrm{mod}}$. Consequently,
\begin{equation*}
 W^a=\underline{\Hom}_E(X,(Y\otimes_E X)[a])
  \simeq\begin{cases}k&a=-3,0,\\0&\text{otherwise}.
    \end{cases}
\end{equation*}
\end{enumerate}
The bimodule $Y$ and these identifications are chosen independently
of $h$.
\end{Prop}
\begin{proof}
\textup{(\ref{ind:model-cohomology})} We obtain the cohomology bound needed for the cutoff
from bounded one-sided projective models.
By Lemma~\ref{ind:tensor-input}\textup{(\ref{ind:tensor-input-dimensions})}, a projective resolution
$T^\bullet\to{}_BE$ may be chosen in degrees $i$ with $-4\leq i\leq 0$.
After forgetting the right $E$-action, $E\otimes_BT^\bullet$ is a
left $E$-projective model of $\mathcal K$ in these degrees.
Resolving $E_B$ gives the corresponding bounded right $E$-projective
model. Thus both one-sided restrictions of $K^\bullet$ are perfect,
and $H^j(K^\bullet)=0$ for $j<-4$; the vanishing for $j>0$ follows
from the chosen model.

\textup{(\ref{ind:model-bimodule})} Put $N_0=c^{-5}(K^\bullet)$, and let
$\pi_0:K^{-5}\twoheadrightarrow N_0$ be the quotient map.
By \textup{(\ref{ind:model-cohomology})}, the cokernel model for $\Phi_{E^e}$ gives
$\overline K\simeq N_0[5]$.
Thus four cosyzygies of $N_0$ will represent $\overline K[-1]$.
We show that $N_0$ is projective on each side by using a truncation
triangle. Let $K_{\ge-4}^\bullet$ be the brutal truncation of $K^\bullet$
to degrees $-4,\ldots,0$. In the quotient complex
$K^\bullet/K_{\ge-4}^\bullet$, the differential out of degree $-5$
is zero. Its cohomology there is therefore
$\operatorname{coker}(d_K^{-6})=N_0$, and its lower cohomology
vanishes by \textup{(\ref{ind:model-cohomology})}. The map induced by $\pi_0$ in degree $-5$
is thus a quasi-isomorphism to $N_0[5]$, giving a triangle in
$\mathscr D^b(E^e)$
\[
 K_{\ge-4}^\bullet\longrightarrow K^\bullet\longrightarrow N_0[5]
 \longrightarrow K_{\ge-4}^\bullet[1].
\]
On either $E$-side, $K_{\ge-4}^\bullet$ is bounded projective and $K^\bullet$
is perfect by \textup{(\ref{ind:model-cohomology})}. Hence $N_0$ has finite projective
dimension on each side. A finite projective resolution over the
self-injective algebra $E$ splits successively from its left end,
since projective modules are injective. Thus $N_0$ is projective,
and hence injective, on each side.
The differential induces a bimodule map
$\bar d:N_0\to K^{-4}$ satisfying $\bar d\pi_0=d_K^{-5}$.
Its kernel is $H^{-5}(K^\bullet)=0$, and its image is
$\operatorname{im}(d_K^{-5})$. Since $N_0$ is injective as a left and as a right $E$-module,
the embedding $\bar d$ splits on each side, with the two retractions chosen separately.

Choose successive $E^e$-injective envelopes to form the sequences
\[
 0\longrightarrow N_r\longrightarrow Q_r
   \longrightarrow N_{r+1}\longrightarrow0,
 \qquad r=0,1,2,3,
\]
and set $Y=N_4$.
Since $E$ and $E^e$ are symmetric and $N_0$ is projective on each $E$-side,
all four sequences split on each side, inductively. Thus $Y$ is
projective on each side, and their connecting morphisms specify
the stable isomorphism
$Y\simeq N_0[4]\simeq \overline K[-1]$.
These choices are independent of $h$.

\textup{(\ref{ind:model-evaluation})} To identify the evaluated model with $\overline J[-1]$,
we compare $K^\bullet\otimes_E X$ with $J$ and then tensor the
cutoff and cosyzygy sequences with $X$. The bounded-above complex $K^\bullet$ has right $E$-projective terms and
therefore computes the derived tensor product with $X$. Derived
associativity gives
\[
 K^\bullet\otimes_E X\simeq
 (E\otimes_B^{\mathbf L}E)\otimes_E^{\mathbf L}X
 \simeq E\otimes_B^{\mathbf L}X\simeq J.
\]
Under this identification, $\mu\otimes_E^{\mathbf L}X$ is the
induction counit $J\to X$ in Proposition~\ref{ind:one-generator}.
As in the proof of Proposition~\ref{frame:factory}, tensoring an
$E^e$-projective module with $X$ gives a projective left $E$-module.
Thus $K^\bullet\otimes_E X$ is a bounded-above projective model of $J$.

The induction triangle of Proposition~\ref{ind:one-generator} shows that $J$ has cohomology
concentrated in degrees $-4$ and $0$, so the cutoff $5$ is admissible.
Right exactness gives
$c^{-5}(K^\bullet\otimes_E X)\simeq N_0\otimes_E X$.
Under this identification, the quotient map is
$\pi_0\otimes_E1_X$ and the induced differential is
$\bar d\otimes_E1_X$.
The embedding $\bar d:N_0\hookrightarrow K^{-4}$ splits as a morphism
of right $E$-modules by \textup{(\ref{ind:model-bimodule})}, and hence remains injective
after tensoring with $X$.
The cokernel model for $\Phi_E$ therefore identifies
$\overline J\simeq(N_0\otimes_E X)[5]$.
The four cosyzygy sequences split on the right, so tensoring them
with $X$ preserves exactness and sends their middle terms to
projective $E$-modules. Their connecting isomorphisms give
$Y\otimes_E X\simeq(N_0\otimes_E X)[4]\simeq\overline J[-1]$.
The cutoff projections commute with tensoring.
The four connecting isomorphisms before and after evaluation come
from the same four right-split short exact sequences, with projective
middle terms. They are therefore compatible with the natural shift
comparisons and our fixed convention for connecting morphisms.
Applying Proposition~\ref{ind:Tate-and-middle}\textup{(\ref{ind:Tate-middle})} with shift
$a-1$ gives the asserted support of $W^a$.
\end{proof}

\section{Twisting and bimodule maps}
\label{lift:construction-section}

We compute how the two twists act (Proposition~\ref{ind:weights})
and then construct maps into the finite bimodule
(Proposition~\ref{lift:nonzero-lifts}). These are the remaining inputs
for the cone criterion.

We retain the notation of Section~\ref{ind:section}, including
the fixed bimodule $Y$ and the stable morphism spaces $H^a$ and $W^a$.

We obtain the twists by extending the fixed involution to $E$
and rescaling its square-zero ideal $DB$.
Put $\theta=\gamma_*\otimes1_C:B\to B$. Then $\theta$ is an involutive automorphism because $\gamma_*:C_*\to C_*$ is an involutive automorphism. Since $\gamma_*$ preserves $\chi_*$, we have
$\chi_B\theta=\chi_B$. Its extension
$$\widetilde\theta:\;\; E\lra E,\;\; (b,\varphi)\mapsto (\theta(b),\varphi\circ\theta^{-1})$$
is an involutive automorphism of $E$ by naturality of the
trivial-extension construction.
For $h\in k^\times$, define
\begin{equation*}
 \upsilon_h(b,\varphi)=(b,h\varphi),\qquad
 \sigma_h=\widetilde\theta\upsilon_h: \;\; (b,\varphi)\mapsto (\theta(b), h\varphi\circ\theta^{-1}), \qquad U_h=E_{\sigma_h}.
\end{equation*}
Both $\upsilon_h$ and $\sigma_h$ are automorphisms of $E$, and hence the bimodule $U_h$ is free on both sides. Since $\chi_B\theta=\chi_B$, the automorphisms $\upsilon_h$, $\widetilde\theta$ and $\sigma_h$ (and thus also their inverses) preserve $\chi$.


To compare the twists with derived induction, we use restriction
functors on both $B$-modules and $E$-modules.
Put $\mathcal F_\theta=\operatorname{Res}_{\theta^{-1}}$ on $B$-modules.
By the twist convention of Subsection~\ref{prelim:notation}, the
exact tensor functor $\mathcal T_h=U_h\otimes_E(-)$ is naturally
identified with $\operatorname{Res}_{\sigma_h^{-1}}$, where $\sigma_h^{-1}(b,\varphi)=(\theta^{-1}(b), h^{-1}\varphi\circ\theta)$.
These restriction functors extend to bounded derived categories.
Both functors preserve projective modules. Since $E$ is symmetric,
$\mathcal T_h$ also induces a triangle functor on the stable module
category, with natural shift comparisons. Fix the canonical identification
\begin{equation*}
 \iota_h:\;\;\mathcal T_hX=U_h\otimes_EX\xrightarrow{\sim}X,
 \qquad t\otimes x\longmapsto\sigma_h^{-1}(t)x=\chi(t)x,\qquad t\in E,\; x\in X.
\end{equation*}
Its inverse is $x\mapsto1\otimes x$.
Using the natural shift comparisons, it determines operators on $H^a$ for $a\in\mathbb{Z}$ by
\[
 T_h^a:\;\; H^a\lra H^a, \qquad \alpha\mapsto \iota_h[a]\circ\mathcal T_h(\alpha)\circ\iota_h^{-1},
 \qquad\alpha\in H^a.
\]
For the degree $-1$ action and the later lifting argument, we fix
the element $\zeta=(0,\chi_B)\in E$.
Trivial-extension multiplication and $\chi_B\theta=\chi_B$ give
$$\sigma_h(\zeta)=h\zeta \qquad\mbox{and}\qquad t_1\zeta t_2=\chi(t_1)\chi(t_2)\zeta,\qquad t_1, t_2\in E.$$

\subsection{Twisting weights from the induction triangle}

The following computation verifies the action condition in
Proposition~\ref{frame:factory} when the twisting parameters are
$1$ and $q$.

\begin{Prop}\label{ind:weights}
For every $m\ge0$ and $h\in k^\times$,
\begin{equation*}
 T_h^{5m}=h^{-m}1_{H^{5m}},\qquad
 T_h^{-5m-1}=h^{m+1}1_{H^{-5m-1}}.
\end{equation*}
These identities hold in every characteristic.
\end{Prop}
\begin{proof}
The powers of $\varpi$ determine the nonnegative-degree weights,
so we first compute $T_h^5(\varpi)$ by comparing the induction
triangle with its twist. For the negative degrees,
Proposition~\ref{ind:Tate-and-middle} reduces the calculation to
$T_h^{-1}$ via multiplication by powers of $\varpi$.
We use the restriction model
$\mathcal T_h=\operatorname{Res}_{\sigma_h^{-1}}$. Since
$\chi\sigma_h^{-1}=\chi$ and $\chi_B\theta^{-1}=\chi_B$, we have
$\mathcal T_hX=X$ and $\mathcal F_\theta X=X$; the identification
$\iota_h$ is then $1_X$. Restriction commutes with shifts of complexes.
In the stable category, all shifted identifications below use the
natural shift comparisons of the induced triangle functor.

\smallskip
\noindent\emph{Step 1. We construct an isomorphism from the twisted induction triangle to the original triangle.}
\par\nobreak\smallskip\nobreak
\noindent
For ${_E}Z{_B}=DB,E,B$, we use the coefficient maps
\[
 s_{DB,h}(\varphi)=h\varphi\circ\theta^{-1},\qquad
 s_{E,h}=\sigma_h,\qquad s_{B,h}=\theta,
\]
where $\varphi\in DB$. They satisfy
$$s_{Z,h}(tzb)=\sigma_h(t)s_{Z,h}(z)\theta(b),\qquad t\in E, z\in Z, b\in B.$$
These identities and the invertibility of the coefficient maps induce
natural $E$-linear isomorphisms for $B$-modules $V$:
$$\mathcal T_h(Z\otimes_BV)\lraf{\simeq} Z\otimes_B\mathcal F_\theta V,
\qquad z\otimes v\mapsto s_{Z,h}(z)\otimes v, \qquad z\in Z, v\in V;$$
their inverses are given by $z\otimes v\mapsto s_{Z,h}^{-1}(z)\otimes v$.
Since restriction is exact and preserves projectives, deriving
these coefficient comparisons gives
\[
 \rho_{Z,h,V}:\;\;\mathcal T_h(Z\otimes_B^{\mathbf L}\cpx{V})
 \xrightarrow{\sim}Z\otimes_B^{\mathbf L}\mathcal F_\theta \cpx{V},
 \qquad\cpx{V}\in\mathscr D^b(B).
\]
Let $\mathord{\text{\usefont{OML}{cmm}{m}{it}\symbol{124}}}:DB\to E$ and $\pi:E\to B$ be the inclusion and quotient, respectively.
The identities
\[
 \sigma_h\circ\mathord{\text{\usefont{OML}{cmm}{m}{it}\symbol{124}}}=\mathord{\text{\usefont{OML}{cmm}{m}{it}\symbol{124}}}\circ s_{DB,h},
 \qquad\pi\circ\sigma_h=\theta\circ\pi
\]
show that the coefficient comparisons respect the short exact
sequence $0\to DB\to E\to B\to 0$.
At $V=X$, the maps $\rho_{DB,h,X}$ and $\rho_{E,h,X}$ give isomorphisms
$c_{L,h}:\mathcal T_hL\xrightarrow{\sim}L$ and $c_{J,h}:\mathcal T_hJ\xrightarrow{\sim}J$. The quotient
comparison is $1_X$, since $\mathcal F_\theta X=X$ and $\chi_B\theta=\chi_B$.
Functoriality of derived tensor gives the following commutative diagram of
triangles, where the lower row is the original induction triangle in $\mathscr D^b(E)$
given in the proof of Proposition \ref{ind:one-generator}:
\[
 \begingroup\setlength{\arraycolsep}{2pt}
 \begin{array}{ccccccc}
 \mathcal T_hL&\xrightarrow{\mathcal T_h(\ell)}&\mathcal T_hJ
 &\xrightarrow{\mathcal T_h(\epsilon)}&X
 &\xrightarrow{\mathcal T_h(\delta)}&(\mathcal T_hL)[1]\\[2pt]
 {\scriptstyle c_{L,h}}\downarrow
 &&{\scriptstyle c_{J,h}}\downarrow
 &&{\scriptstyle1_X}\downarrow
 &&\downarrow{\scriptstyle c_{L,h}[1]}\\[2pt]
 L&\xrightarrow{\ell}&J&\xrightarrow{\epsilon}&X
 &\xrightarrow{\delta}&L[1].
 \end{array}
 \endgroup
\]

\smallskip
\noindent\emph{Step 2. We show that $c_{L,h}$ corresponds to $h1_{X[4]}$ under $L\simeq X[4]$.}
\par\nobreak\smallskip\nobreak
\noindent
We first check the case $h=1$ by duality.
Put $\mathcal V=\Ext_B^4(X,B)$, and define the induced action
\[
 \mathcal A_\theta(\alpha)
 =\theta[4]\circ\mathcal F_\theta(\alpha),
 \qquad\alpha:X\longrightarrow B[4].
\]
Here $\theta:\mathcal F_\theta B={_{\theta^{-1}}B}\to B$ is the coefficient
isomorphism, and the source is identified using
$\mathcal F_\theta X=X$. By the Ext K\"unneth
isomorphism in the proof of Lemma~\ref{ind:tensor-input},
\[
\begin{gathered}
\Ext^2_{C_*}(S_*, C_*)\otimes_k\Ext^2_C(S, C)\lraf{\simeq} \mathcal V,\\
f_1\otimes f_2\mapsto [X=S_*\otimes_k S\lraf{f_1\otimes_k f_2} (C_*[2])\otimes_k (C[2])=B[4]],
\end{gathered}
\]
where $f_1: S_*\to C_*[2]$ and $f_2: S\to C[2]$ are derived morphisms. By $\theta=\gamma_*\otimes 1_C$, the map $\mathcal A_\theta$ is the tensor product of the $\gamma_*$-action on $\Ext_{C_*}^2(S_*,C_*)$ and the identity on $\Ext_C^2(S,C)$. Recall that $\theta$ and $\gamma_*$ are involutive automorphisms.
Thus  Lemma~\ref{seed:gamma}\,\textup{(\ref{seed:gamma-cohomology})} gives $\mathcal A_\theta=1_{\mathcal V}$.

The natural tensor--Hom duality $DL\simeq\mathbf R\Hom_B(X,B)$ identifies $H^{-4}(L)$ with
$D\mathcal V$. We choose a bounded-above projective resolution
$P_X^\bullet\to X$ and a comparison
$\kappa_X:P_X^\bullet\to\mathcal F_\theta P_X^\bullet$ lifting
$1_X:X\to\mathcal F_\theta X=X$.
Both complexes resolve $X$, so $\kappa_X$ is a homotopy equivalence.
Write $\rho_X$ for the complex-level coefficient comparison
representing $\rho_{DB,1,X}$ from Step 1. Then $c_{L,1}^{-1}$ is
represented by
\[
 DB\otimes_BP_X^\bullet
 \xrightarrow{\ 1_{DB}\otimes\kappa_X\ }
 DB\otimes_B\mathcal F_\theta P_X^\bullet
 \xrightarrow{\ \rho_X^{-1}\ }
 \mathcal T_1(DB\otimes_BP_X^\bullet).
\]
Under tensor--Hom duality, identifying the underlying $k$-complexes
before and after restriction, the dual of $c_{L,1}^{-1}$ is represented by
\[
\begin{aligned}
  D(c_{L,1}^{-1}):\quad &
  \Hom_B^\bullet(P_X^\bullet,B)
 &\lraf{D(\rho_X^{-1})}\Hom_B^\bullet(\mathcal F_\theta P_X^\bullet,B)
 \longrightarrow\Hom_B^\bullet(P_X^\bullet,B),\\
&\quad\qquad f&\longmapsto\theta\circ\mathcal F_\theta(f)
 \longmapsto\theta\circ\mathcal F_\theta(f)\circ\kappa_X.
 \end{aligned}
\]
Here the first arrow comes from dualizing the inverse coefficient
map $\varphi\mapsto\varphi\circ\theta$ on $DB$, whose dual is $\theta$
on $B$; the second arrow is precomposition with $\kappa_X$ and gives
the specified source comparison. The induced map
$H^4\bigl(D(c_{L,1}^{-1})\bigr)$ is precisely $\mathcal A_\theta$,
so $H^4\bigl(D(c_{L,1}^{-1})\bigr)=1_{\mathcal V}$.
Duality gives $H^{-4}(c_{L,1})=1_{D\mathcal V}$.

The $E$-action on $L$ factors through $B$, and
$\sigma_h^{-1}|_B=\theta^{-1}$ is independent of $h$.
Thus $\mathcal T_hL$ and $\mathcal T_1L$ are naturally identified.
Moreover, $s_{DB,h}=h\,s_{DB,1}$, so $k$-linearity of derived
tensor gives $c_{L,h}=h\,c_{L,1}$ under this identification.
Hence $H^{-4}(c_{L,h})=h1_{H^{-4}(L)}$.
Let $\eta:L\xrightarrow{\sim}X[4]$ denote the cohomology isomorphism
used in the proof of Proposition~\ref{ind:one-generator}. Since $L$ has cohomology only
in degree $-4$, naturality of the truncation isomorphism gives
\[
 \begin{array}{ccc}
 \mathcal T_hL&\xrightarrow{\ \mathcal T_h(\eta)\ }&X[4]\\[2pt]
 {\scriptstyle c_{L,h}}\downarrow
 &&\downarrow{\scriptstyle h1_{X[4]}}\\[2pt]
 L&\xrightarrow{\ \eta\ }&X[4].
 \end{array}
\]

\smallskip
\noindent\emph{Step 3. We prove that $T_h^5(\varpi)=h^{-1}\varpi$ and deduce the weights in nonnegative degrees.}
\par\nobreak\smallskip\nobreak
\noindent
We write the induction triangle as
$$X[4]\lraf{u}J\lraf{v}X\lraf{\varpi}X[5].$$
The comparisons in Steps 1 and 2 give the following commutative
diagram of triangles:
\[
 \begingroup\setlength{\arraycolsep}{3pt}
 \begin{array}{ccccccc}
 X[4]&\xrightarrow{\mathcal T_h(u)}&\mathcal T_hJ
 &\xrightarrow{\mathcal T_h(v)}&X
 &\xrightarrow{\mathcal T_h(\varpi)}&X[5]\\[2pt]
 {\scriptstyle h 1_{X[4]}}\downarrow
 &&{\scriptstyle c_{J,h}}\downarrow
 &&{\scriptstyle1_X}\downarrow
 &&\downarrow{\scriptstyle h 1_{X[5]}}\\[2pt]
 X[4]&\xrightarrow{u}&J&\xrightarrow{v}&X
 &\xrightarrow{\varpi}&X[5].
 \end{array}
 \endgroup
\]
The last vertical map is the shift of the first, so it also has
the factor $h$. Since restriction is exact and preserves projectives,
the cokernel model in Subsection~\ref{prelim:stable} gives a natural
isomorphism $\Phi_E\mathcal T_h\simeq\mathcal T_h\Phi_E$, compatible
with the shift comparisons induced by cosyzygy sequences.
Applying $\Phi_E$ to the diagram, its rightmost square therefore gives
$h\,T_h^5(\varpi)=\varpi$, and hence $T_h^5(\varpi)=h^{-1}\varpi$.
The triangle functor $\mathcal T_h$ preserves shifted composition. Hence
$T_h^{5m}(\varpi^m)=h^{-m}\varpi^m$ for $m\ge0$.
Since $H^{5m}=k\varpi^m$, this proves the nonnegative-degree formula.

\smallskip
\noindent\emph{Step 4. We prove that $T_h^{-1}=h1_{H^{-1}}$ and deduce the remaining negative-degree weights.}
\par\nobreak\smallskip\nobreak
\noindent
We use the character surjection onto $X$. Put $K_X=\ker(\chi)$, and consider
the short exact sequence
\[
 0\longrightarrow K_X\longrightarrow E\lraf{\chi}X
 \longrightarrow0.
\]
Let $\delta_X:X\to K_X[1]$ be its stable connecting morphism.
Set $b_X=\delta_X[-1]$. Since $E$ is projective, the morphism $b_X:X[-1]\to K_X$
is an isomorphism in the stable category.
The map $\sigma_h:\mathcal T_hE\to E$ restricts to
$\kappa_h=\sigma_h|_{K_X}:\mathcal T_hK_X\to K_X$ and induces
$1_X$ on the quotient, because $\chi\sigma_h=\chi$.
Applying inverse rotation to the induced stable triangles and using
the natural shift comparisons, we obtain the commutative diagram
\[
 \begingroup\setlength{\arraycolsep}{2pt}
 \begin{array}{ccccccc}
 \mathcal T_h(X[-1])
 &\xrightarrow{-\mathcal T_h(b_X)}&\mathcal T_hK_X
 &\longrightarrow&\mathcal T_hE
 &\xrightarrow{\mathcal T_h(\chi)}&X\\[2pt]
 {\scriptstyle\xi_h}\downarrow
 &&{\scriptstyle\kappa_h}\downarrow
 &&{\scriptstyle\sigma_h}\downarrow
 &&\downarrow{\scriptstyle1_X}\\[2pt]
 X[-1]&\xrightarrow{-b_X}&K_X
 &\longrightarrow&E&\xrightarrow{\chi}&X.
 \end{array}
 \endgroup
\]
Here $\xi_h:\mathcal T_h(X[-1])\xrightarrow{\sim}X[-1]$
is induced by the shift comparison and $\mathcal T_hX=X$;
the minus signs come from inverse rotation.
The leftmost square gives
$\kappa_h\circ\mathcal T_h(b_X)=b_X\circ\xi_h$.

Since $\chi(\zeta)=0$, the element $\zeta$ lies in $K_X$.
The identity $t\zeta=\chi(t)\zeta$ for any $t\in E$ shows that
$i_X:X\to K_X$, $1\mapsto\zeta$, is nonzero and $E$-linear.
Symmetry and coinduction give
$\Hom_E(X,E)\simeq\Hom_E(X,DE)\simeq DX\simeq k$,
and hence $\Hom_E(X,K_X)=ki_X$.
Since $b_X$ induces an isomorphism
\[
 H^{-1}\xrightarrow{\sim}\underline{\Hom}_E(X,K_X),
 \qquad\beta\longmapsto b_X\circ\beta,
\]
and $H^{-1}\simeq k$ by Proposition~\ref{ind:Tate-and-middle},
the stable class $\underline{i_X}$ is nonzero and spans the target.
Since $\kappa_h(\zeta)=\sigma_h(\zeta)=h\zeta$, we have
$\kappa_h\circ\mathcal T_h(i_X)=h\,i_X$.

Let $\beta_0\in H^{-1}$ be the unique class with
$b_X\circ\beta_0=\underline{i_X}$. It generates $H^{-1}$.
The identity above and functoriality give
\[
 \begin{aligned}
 b_X\circ T_h^{-1}(\beta_0)
 &=b_X\circ\xi_h\circ\mathcal T_h(\beta_0)
  =\kappa_h\circ\mathcal T_h(b_X)\circ\mathcal T_h(\beta_0)=\kappa_h\circ\mathcal T_h(b_X\circ\beta_0)\\
 &=\underline{\kappa_h\circ\mathcal T_h(i_X)}=h\underline{i_X}=h\,b_X\circ\beta_0.
 \end{aligned}
\]
Since $b_X$ is invertible, $T_h^{-1}(\beta_0)=h\beta_0$, and
therefore $T_h^{-1}=h1_{H^{-1}}$.

For $m\ge0$, Proposition~\ref{ind:Tate-and-middle} gives the
multiplication isomorphism
\[
 \mu^{(m)}:H^{-5m-1}\xrightarrow{\sim}H^{-1},\qquad
 \beta\longmapsto\varpi^m[-5m-1]\circ\beta,
\]
where $\beta\in H^{-5m-1}$. The identities
$T_h^{-1}=h1_{H^{-1}}$ and
$T_h^{5m}(\varpi^m)=h^{-m}\varpi^m$, together with preservation
of shifted composition, give
\[
 \begin{aligned}
 h\mu^{(m)}(\beta)
 &=T_h^{-1}\bigl(\mu^{(m)}(\beta)\bigr)
  =T_h^{-1}\bigl(\varpi^m[-5m-1]\circ\beta\bigr)
  =T_h^{5m}(\varpi^m)[-5m-1]\circ T_h^{-5m-1}(\beta)\\
 &=h^{-m}\varpi^m[-5m-1]\circ T_h^{-5m-1}(\beta)
  =h^{-m}\mu^{(m)}\bigl(T_h^{-5m-1}(\beta)\bigr).
 \end{aligned}
\]
Since $\mu^{(m)}$ is injective,
$T_h^{-5m-1}(\beta)=h^{m+1}\beta$, as required.
\end{proof}

\subsection{An obstruction from a symmetric corner algebra}

The obstruction to perfect factorization used below comes from a
symmetric corner algebra of the fixed auxiliary seed. The variable seed
requires neither an automorphism nor such a corner algebra. Set
\[
 R=f_*C_*f_*=k[t]/(t^2),\qquad
 \omega=\gamma_*|_R,\qquad b=f_*\otimes1_C.
\]
The fixed seed satisfies $\gamma_*(f_*)=f_*$ and
$\chi_*(f_*)=1$. Hence
$$\theta(b)=b,\qquad bBb=R\otimes_k C, \qquad
\theta(x\otimes z)=\omega(x)\otimes z$$ for $x\in R$ and $z\in C$.
Write $\chi_R=\chi_*|_R$; see Subsection~\ref{seed:concrete-pair} and
Lemma~\ref{seed:gamma}\textup{(\ref{seed:gamma-involution})}. Then $\omega(t)=-t$,
$\chi_R(t)=0$, and $\chi_B|_{bBb}=\chi_R\otimes\chi_C$, due to $\chi_B=\chi_*\otimes\chi_C$.
The standard $B$-bimodule structure on $DB$ is
\[
 (b_1\lambda b_2)(b_3)=\lambda(b_2b_3b_1)
 \qquad(b_1, b_2, b_3\in B,\ \lambda\in DB).
\]
Let $(DB)_\theta$ have the same left action and right action $\lambda\cdot b_2=\lambda\theta(b_2)$.
With $B$ carrying its usual $B$-bimodule structure, we define the character map
\[
 \eta_\theta:\;\;B\longrightarrow(DB)_\theta,\qquad
 \eta_\theta(b')=\chi_B(b')\chi_B\quad(b'\in B).
\]
Likewise, $R_\omega$ and $(DR)_\omega$ denote the corresponding
$R$-bimodules with their usual left actions and right actions
induced by $\omega$.

For a $B$-bimodule $L$, we set $\mathcal R(L)=bLb$ and regard it as an
$R$-bimodule through the canonical inclusion $R\to bBb$, $r\mapsto r\otimes1_C$.
This induces an additive functor
$\mathcal R\colon B^e\modcat\to R^e\modcat$.
The following lemma supplies the properties needed to apply the functor $\mathcal R$ to a factorization through a perfect complex.

\begin{Lem}\label{lift:corner-preparation}
The map $\eta_\theta$ is a $B$-bimodule homomorphism.
The functor $\mathcal R$ is exact and preserves projective modules.
Consequently, its termwise extension $\mathcal R:\mathscr D^b(B^e)\to\mathscr D^b(R^e)$
preserves perfect objects.
\end{Lem}

\begin{proof}
We regard $k$ as a $B$-bimodule with both actions through
$\chi_B$. Since $\chi_B\theta=\chi_B$, the functional $\chi_B$
spans a copy of this bimodule in $(DB)_\theta$. Hence
$\eta_\theta$ is the composite of bimodule maps
$B\xrightarrow{\chi_B}k\hookrightarrow(DB)_\theta$.

The corner functor $B^e\modcat\to(bBb)^e\modcat$, $L\mapsto bLb$, and restriction along $R\hookrightarrow bBb$ are exact,
so $\mathcal R$ is exact. To prove preservation of projectives,
it suffices to show that the image of a rank-one free
$B^e$-module is $R^e$-projective. To describe its two tensor factors, we first show that
${}_R(f_*C_*)\cong R^2$ and $(C_*f_*)_R\cong R^2$.
For the right-hand decomposition, we use $C_*f_*=e_*C_*f_*\oplus R$ and
$e_*C_*f_*\cong M_*$ as right $R$-modules. The action of $t$ on
$M_*$ sends $u$ to $v$ and $v$ to zero, so $u$ generates a free
rank-one right $R$-module. For the left-hand decomposition, Frobenius duality gives
$f_*C_*e_*=\Hom_{H_*}(M_*, H_*)\simeq \Hom_{H_*}(M_*, DH_*)\simeq DM_*$ as left $R$-modules.
Since $M_*\cong R_R$ and $R$ is symmetric, it follows that
$DM_*\simeq {}_R R$; adding the summand $f_*C_*f_*=R$ gives the
left-hand decomposition.

Since $bB=f_*C_*\otimes_kC$ and $Bb=C_*f_*\otimes_kC$, the left $R$-module $bB$ and the right $R$-module
$Bb$ are free of rank $2\dim_k(C)$. Therefore the image of a rank-one free $B^e$-module is
\[
 \mathcal R(B\otimes_k B)\simeq bB\otimes_k Bb=f_*C_*\otimes_kC\otimes_kC_*f_*\otimes_kC
 \simeq (R^e)^{\,4(\dim_k(C)^2)}
 \quad\text{as $R$-bimodules}.
\]
Passing to finite sums and direct summands proves that $\mathcal R$ preserves projective modules. Applying it termwise to bounded complexes of projectives proves the last assertion.
\end{proof}

The following proposition gives the character-map obstruction
needed for our construction; compare \cite[Lemma~6.5]{O} and
\cite[Proposition~5.3]{F}. We prove it by restriction to the fixed
symmetric corner algebra $R$.

\begin{Prop}\label{lift:trace-obstruction}
The character map $\eta_\theta$ does not factor
through a perfect object of $\mathscr D^b(B^e)$.
\end{Prop}

\begin{proof}
We reduce the assertion to showing that the $R$-bimodule map
$\psi:R\to R_\omega$, sending $r\in R$ to $rt$, has a nonzero stable class.
Indeed, we obtain $\psi$ from $\eta_\theta$ by applying $\mathcal R$
and composing with bimodule maps. These operations preserve
factorizations through perfect objects.

We first identify the target after restriction. Restriction of functionals gives a $(bBb)$-bimodule isomorphism $b(DB)b\cong D(bBb)$. Its inverse sends
$\lambda\in D(bBb)$ to the functional $b'\mapsto\lambda(bb'b)$ on $B$, where $b'\in B$.
Since $\theta(b)=b$ and its restriction to $bBb=R\otimes_k C$ is
$\omega\otimes1_C$, this identifies
$\mathcal R((DB)_\theta)$ with $D(R\otimes_k C)$, with its usual
dual left $R$-action and its dual right action precomposed with
$\omega$. Indeed, for $\lambda\in D(R\otimes_k C)$,
$a,d, r\in R$ and $z\in C$, the twisted right action gives
\[
 (a\cdot\lambda\cdot d)(r\otimes z)
 =\lambda\bigl((\omega(d)\otimes1_C)(r\otimes z)(a\otimes1_C)\bigr)
 =\lambda\bigl((\omega(d)ra)\otimes z\bigr).
\]
Let $i:R\to\mathcal R(B)=R\otimes_k C$ be the bimodule map $r\mapsto r\otimes1_C$. Dualizing $i$ and twisting the right actions by $\omega$ gives an $R$-bimodule map
$\rho:\mathcal R((DB)_\theta)\to(DR)_\omega$. This can be seen as follows.

Under the preceding identification with $D(R\otimes_k C)$,
we have $\rho(\lambda)(r)=\lambda(r\otimes1_C)$ for
$\lambda\in\mathcal R((DB)_\theta)$ and $r\in R$.
For $a,d,r\in R$, the action formula gives
\[
 \rho(a\cdot\lambda\cdot d)(r)
 =\lambda((\omega(d)ra)\otimes1_C)
 =(a\cdot\rho(\lambda)\cdot d)(r).
\]
This shows $\rho(a\cdot\lambda\cdot d)=a\cdot\rho(\lambda)\cdot d$, and thus $\rho$ is an $R$-bimodule map.

The identities
$\chi_B|_{bBb}=\chi_R\otimes\chi_C$ and $\chi_C(1_C)=1$ give
\[
 \bigl(\rho\,\mathcal R(\eta_\theta)\,i(r)\bigr)(x)
 =\eta_\theta(r\otimes1_C)(x\otimes1_C)
 =\chi_R(r)\chi_R(x),
\]
where $r,x\in R$. Hence the composite $\rho\,\mathcal R(\eta_\theta)\,i: R\to (DR)_\omega$ sends $r$ to $\chi_R(r)\chi_R$.

Note that $R$ is a commutative and symmetric algebra with a symmetrizing functional
$\delta_R:R\to k$ given by $\delta_R(k_1+k_2 t)=k_2$ for $k_1, k_2\in k$.
In particular, the map $\phi: R\to DR$, $a\mapsto[r\mapsto\delta_R(ra)]$, is an isomorphism of $R$-bimodules. It sends $1,t$ to $\delta_R, \chi_R$, respectively, since $\delta_R(rt)=\chi_R(r)$.
Further, $\delta_R$ and $\chi_R$ form a basis of $DR$. Twisting both right actions by $\omega$ gives
an isomorphism $R_\omega\xrightarrow{\sim}(DR)_\omega$, which we continue to denote by $\phi$.
Composing $\rho\,\mathcal R(\eta_\theta)\,i$ with the inverse of
$\phi$ gives the $R$-bimodule homomorphism
\[
 \psi=\phi^{-1}\rho\,\mathcal R(\eta_\theta)\,i:\;\; R\longrightarrow R_\omega,\qquad
 r\mapsto \chi_R(r)t=rt\quad(r\in R),
\]
where the last equality follows by writing $r=k_1+k_2 t$, with $k_1, k_2\in k$.
By Lemma~\ref{lift:corner-preparation}, a perfect factorization
of $\eta_\theta$ would give a perfect factorization of $\psi$ in
$\mathscr D^b(R^e)$. It remains to rule out such a factorization.

We do this by showing that $\psi$ has a nonzero stable class. Set
\[
 R^e=k[\ell,r]/(\ell^2,r^2),\qquad
 u=\ell-r,\quad v=\ell+r,
\]
where $\ell$ and $r$ denote $t\otimes1$ and $1\otimes t$ in $R\otimes_k R^{\rm op}$, respectively.
The element $u$ annihilates the source $R$, so the image of $1$
under any $R^e$-linear map $R\to R^e$ must belong to
$\operatorname{ann}_{R^e}(u)$.
To compute this annihilator, we write
$w=a_0+a_1\ell+a_2r+a_3\ell r$, with $a_0,a_1,a_2,a_3\in k$.
Then $uw=a_0(\ell-r)+(a_2-a_1)\ell r$, and hence
$\operatorname{ann}_{R^e}(u)=k(\ell+r)\oplus k\ell r=R^ev$.
On the target $R_\omega$, the elements $\ell$ and $r$ act as
multiplication by $t$ and $\omega(t)=-t$, respectively. Thus $v$
annihilates $R_\omega$.
Consequently, every pair of $R^e$-linear maps
$\alpha:R\to R^e$ and $\beta:R^e\to R_\omega$ satisfies
$\alpha(1)\in R^ev$ and $\beta(R^ev)=0$, so $\beta\alpha=0$.
The same conclusion holds for a factorization through a finite sum
of copies of $R^e$, and hence through any finitely generated
projective $R^e$-module.
Since $\psi(1)=t\ne0$, the stable class of $\psi$ is nonzero.
This argument also applies in characteristic two, when $u=v$.

Since $R^e$ is symmetric, the triangle functor $\Phi_{R^e}$
sends $\psi$ to its nonzero stable class and annihilates perfect objects. Therefore $\psi$ cannot factor through a perfect object of $\mathscr D^b(R^e)$.
The reduction above proves the assertion for $\eta_\theta$.
\end{proof}


\subsection{Bimodule maps with nonzero stable evaluation}
\label{lift:nonzero-section}

We now use the obstruction just proved to construct ordinary bimodule maps
whose evaluations at $X$ have nonzero stable classes. The twisted
sources $U_h$ and their identifications $\iota_h$ are those fixed
at the beginning of this section.
We use the bimodule $Y$ and the stable multiplication map
$m=\Phi_{E^e}(m_K):\overline K\to E$ from
Subsection~\ref{ind:finite-Y}, where $m_K$ represents derived
multiplication $\mu$.

Evaluation on separately projective bimodules is well defined
on stable morphisms: $E^e$-projectives evaluate to $E$-projectives,
as in the proof of Proposition~\ref{frame:factory}.
Syzygy and cosyzygy sequences of such bimodules with
$E^e$-projective middle terms split on each $E$-side, so their kernels and cokernels remain
separately projective. Tensoring these right-split sequences with
$X$, as in the proof of Proposition~\ref{ind:Y-evaluation}\textup{(\ref{ind:model-evaluation})},
gives the natural shift comparisons. Thus evaluation preserves
stable isomorphisms and commutes with shifts. In particular,
$\overline K$ and all its shifts admit such representatives.

\Needspace{14\baselineskip}
The following stable form of tensor--Hom adjunction follows the argument of
\cite[Lemma~6.6]{O}; see also \cite[(21)]{F}. We use it to detect
nonvanishing after evaluation at the character module.

\begin{Lem}\label{lift:stable-adjunction}
Let $E$ be a finite-dimensional symmetric $k$-algebra, let
$\chi:E\to k$ be a character, and put $X=k_\chi$.
Regard $k$ as an $E$-bimodule via
$a\cdot c\cdot b=\chi(a)c\chi(b)$ for $a,b\in E$ and $c\in k$.
For every separately projective $E$-bimodule $L$, the assignment
$\psi\mapsto\widehat\psi$, where
\[
 \widehat\psi(\ell\otimes x)=\psi(\ell)x
 \qquad(\ell\in L,\ x\in X),
\]
induces a $k$-linear isomorphism
\[
 \underline{\Hom}_{E^e}(L,k)
 \xrightarrow{\sim}\underline{\Hom}_E(L\otimes_E X,X),
 \qquad \underline\psi\longmapsto\underline{\widehat\psi},
\]
natural in $L$.
\end{Lem}

\begin{proof}
For any $E$-bimodule $L$, the ordinary tensor--Hom adjunction gives
\[
 \Hom_E(L\otimes_E X,X)
 \simeq\Hom_{E^e}(L,\Hom_k(X,X))
 \simeq\Hom_{E^e}(L,k).
\]
Here the $E$-bimodule action on $\Hom_k(X,X)$ is
$(afb)(x)=a\,f(bx)$ for $a,b\in E$, $f\in\Hom_k(X,X)$ and $x\in X$.
Since $X=k_\chi$, the map $c1_X\mapsto c$ for $c\in k$
identifies this bimodule with $k$ carrying the stated character actions.
Under these identifications, $\psi$ corresponds to $\widehat\psi$.
We show that the two subspaces of maps factoring through projective
modules correspond under this adjunction.

If $\psi$ factors through an $E^e$-projective module $Q$, then
$\widehat\psi$ factors through $Q\otimes_E X$.
This is an $E$-projective module: $Q$ is a direct summand of a finite
sum of copies of $E\otimes_k E$ with the outer bimodule actions,
and $(E\otimes_k E)\otimes_E X\simeq E\otimes_k X$ is left $E$-free.

Conversely, suppose that $\widehat\psi=vu$ through an
$E$-projective module $P$, with $u:L\otimes_E X\to P$ and
$v:P\to X$. Since $E^e$ is symmetric, choose an embedding
$j:L\hookrightarrow V$ with $V$ projective-injective over $E^e$.
The right $E$-module $L$ is projective and hence injective, so
$j$ splits on the right. Thus $j\otimes_E1_X$ is a monomorphism.
Since $E$ is self-injective, $P$ is injective, and $u$ extends to
$\widetilde u:V\otimes_E X\to P$ satisfying
$\widetilde u(j\otimes_E1_X)=u$.
Under the ordinary adjunction, $v\widetilde u$ corresponds to an
$E$-bimodule map $\widetilde\psi:V\to k$.
Writing $1$ for the standard vector of $X=k_\chi$, we have
\[
 \widetilde\psi(j(\ell))1
 =(v\widetilde u)(j(\ell)\otimes1)
 =\widehat\psi(\ell\otimes1)=\psi(\ell)1
 \qquad(\ell\in L).
\]
Therefore $\widetilde\psi j=\psi$, so $\psi$ factors through the
$E^e$-projective module $V$.
The ordinary adjunction consequently descends to the asserted
isomorphism of stable Hom spaces. Its naturality in $L$ descends
from the naturality of the ordinary adjunction.
\end{proof}

For the fixed character $\chi$ above, we use the character bimodule
$k$ from Lemma~\ref{lift:stable-adjunction}.
Then $\chi:E\to k$ is an $E$-bimodule homomorphism.
When composing with stable morphisms, we regard ordinary bimodule
maps as stable classes. Put
$\varphi=\chi m:\overline K\to k$.
By Proposition~\ref{ind:Y-evaluation}\textup{(\ref{ind:model-evaluation})},
$W^0=\underline{\Hom}_E(X,Y\otimes_E X)\simeq k$.

\begin{Prop}\label{lift:nonzero-lifts}
For every $h\in k^\times$ there is an $E$-bimodule homomorphism
$g_h:U_h\to Y$ for which the stable class of
\[
 X\xrightarrow{\iota_h^{-1}}U_h\otimes_E X
   \xrightarrow{g_h\otimes1}Y\otimes_E X
\]
is nonzero in $W^0$. The target $Y$ is independent of $h$.
\end{Prop}

\begin{proof}
Since $Y\simeq\overline K[-1]$ by Proposition \ref{ind:Y-evaluation}\textup{(\ref{ind:model-bimodule})}, we seek a stable morphism
$\beta:U_h[1]\to\overline K$ whose evaluation at $X$ (namely, $\beta\otimes_EX$) is nonzero; shifting it will give the desired map to $Y$. Stable duality will produce $\beta$ with $\varphi\beta\ne0$. We will then use this condition to prove nonvanishing after evaluation.


We first define a $k$-linear map $b_0:k\to U_h=E_{\sigma_h}$ sending $1$ to $\zeta=(0, \chi_B)$. Recall from Section \ref{lift:construction-section} that $t_1\zeta t_2=\chi(t_1)\chi(t_2)\zeta$ for any $t_1, t_2\in E$. Since $\chi\sigma_h=\chi$, the map $b_0$ is a homomorphism of $E$-bimodules.
We show that $b_0\varphi\ne0$ in $\underline{\Hom}_{E^e}(\overline K,U_h)$.
To use Proposition~\ref{lift:trace-obstruction}, we show that vanishing of this class would force a perfect factorization of $\eta_\theta$.

Suppose that the stable class $b_0\varphi$ is zero. The derived morphism
$\alpha=b_0\chi m_K:K^\bullet\to U_h[0]$ has image $b_0\varphi$ under $\Phi_{E^e}$, by its compatibility with ordinary module maps. By Proposition~\ref{ind:Y-evaluation}\textup{(\ref{ind:model-cohomology})}, the bounded-above projective model $K^\bullet$ has bounded cohomology. Since $E^e$ is self-injective, Lemma~\ref{lift:zero-perfect} implies that $\alpha$ factors through a bounded complex of projective $E^e$-modules.
Under $K^\bullet\simeq\mathcal K=E\otimes_B^{\mathbf L}E$, this is a perfect factorization of
$b_0\chi\mu$. Here $\mu:\mathcal K\to E$ is the derived multiplication map and is
represented by $m_K:K^\bullet\to E$.

To obtain a factorization of $\eta_\theta$, we use the following derived induction and restriction induced by the canonical inclusion $B^e\to E^e$:
\[
 \mathsf F=E^e\otimes_{B^e}^{\mathbf L}(-)
  \simeq E\otimes_B^{\mathbf L}(-)\otimes_B^{\mathbf L}E,
 \qquad
 \mathsf R=\operatorname{Res}_{B^e}.
\]
We first verify that $\mathsf R$ preserves perfect objects. In fact, a rank-one free $E^e$-module restricts to $E\otimes_kE$ with outer $B$-actions. Tensoring left and right $B$-projective resolutions of $E$, each of length four by Lemma~\ref{ind:tensor-input}\textup{(\ref{ind:tensor-input-dimensions})}, gives a $B^e$-projective resolution of length at most eight. The same bound holds for finite sums and direct summands. Thus the restriction of any bounded complex of projective $E^e$-modules to $B^e$ is perfect. The same argument on the other side gives finite right $B^e$-projective dimension of $E^e$. Hence $(\mathsf F, \mathsf R)$ is an adjoint pair of functors between $\mathscr D^b(B^e)$ and $\mathscr D^b(E^e)$.

We verify the canonical identification $\mathsf F(B)\simeq \mathcal K$ on a common projective model. Choose a bounded-above $B^e$-projective resolution $\epsilon_B:P_B^\bullet\to B$. Its augmented complex splits after restriction to right $B$-modules: $B_B$ is projective, and the successive kernels are projective by induction. Thus $P_B^\bullet\otimes_BE\to E$ is exact. Its terms are direct summands of finite sums of $B\otimes_kE$, so they are projective as $B$--$E$-bimodules. Consequently, $E\otimes_BP_B^\bullet\otimes_BE$ computes both $\mathsf F(B)$ and $\mathcal K$.

Let $u_B:B\to\mathsf R\mathsf F(B)$ be the unit at $B$ and let $\varepsilon_E:\mathsf F\mathsf R(E)\to E$ be the counit at $E$. Identifying $\mathcal K$ with $\mathsf F(B)$, the morphism
$\mu:\mathcal K\to E$ is adjoint to the inclusion
$j_B:B\hookrightarrow\mathsf R(E)$. The unit--counit identities give
$\mu=\varepsilon_E\mathsf F(j_B)$ and $\mathsf R(\mu)u_B=j_B$.
These are derived identities, independent of the representative $m_K$.
Since $\sigma_h|_B=\theta$, restriction gives $\mathsf R(U_h)=B_\theta\oplus(DB)_\theta$.
Let $\operatorname{pr}_{DB}$ be the projection onto the second
summand. It follows from the identity $\mathsf R(\mu)u_B=j_B$ that
\[
 \operatorname{pr}_{DB}\,
 \mathsf R(b_0\chi\mu)u_B= \operatorname{pr}_{DB}\,
 \mathsf R(b_0\chi) \mathsf R(\mu)u_B
 =\operatorname{pr}_{DB}\,\mathsf R(b_0\chi)j_B
 =\eta_\theta.
\]
The last equality is represented by ordinary $B$-bimodule maps and sends
$b_1\in B$ to $\chi_B(b_1)\chi_B$.
Note that the morphism $b_0\chi\mu$ factors through a bounded complex of projective $E^e$-modules
and $\mathsf R$ preserves perfect objects. Consequently, there is a perfect factorization of
$\mathsf R(b_0\chi\mu)$. Composing this factorization with $u_B$ and $\operatorname{pr}_{DB}$ gives a perfect
factorization of $\eta_\theta$. This contradicts Proposition~\ref{lift:trace-obstruction} and proves
$b_0\varphi\ne0$.

We can now use this nonvanishing to construct $\beta$.
Applying the degree-zero case of Lemma~\ref{tool:duality},
we obtain a natural isomorphism of additive functors:
\[
 D\underline{\Hom}_{E^e}(U_h[1], -)
       \simeq\underline{\Hom}_{E^e}(-,U_h):\;\;
       \bigl(E^e\text{-}\underline{\mathrm{mod}}\bigr)^{\mathrm{op}}\lra k\modcat.
\]
Under this isomorphism, the dual of
\[
 \varphi_*:\;\;\underline{\Hom}_{E^e}(U_h[1],\overline K)
        \longrightarrow\underline{\Hom}_{E^e}(U_h[1],k),
 \qquad \beta\longmapsto\varphi\beta
\]
is the map
\[
 \underline{\Hom}_{E^e}(k,U_h)
       \longrightarrow\underline{\Hom}_{E^e}(\overline K,U_h),
 \qquad \alpha\longmapsto \alpha\varphi.
\]
The dual map is nonzero at $b_0$ because
$b_0\varphi\ne0$; finite-dimensionality therefore
implies that $\varphi_*$ is nonzero. So we can choose a stable morphism
$\beta:U_h[1]\to \overline K$ with $\varphi\beta\ne0$.

It remains to prove that $\beta$ stays nonzero after evaluation at $X$.
For this, we take a cosyzygy sequence
$0\to U_h\to Q_h\to L_h\to0$ in $E^e\modcat$ with $Q_h$ projective-injective. Then $L_h$ is separately
projective, and its connecting morphism gives
$L_h\otimes_EX\simeq(U_h\otimes_EX)[1]$ in
$E\text{-}\underline{\mathrm{mod}}$.
Further, we choose a separately projective representative $\widetilde K$ of $\overline K$ and ordinary
$E$-bimodule maps $\beta_0:L_h\to \widetilde K$ and $m_0:\widetilde K\to E$ representing
$\beta$ and $m$, respectively. Set $\psi=\chi\circ m_0\circ\beta_0:L_h\to k$.
Its stable class is $\varphi\beta\ne0$.

By Lemma~\ref{lift:stable-adjunction}, the adjoint
$\widehat\psi:L_h\otimes_E X\to X$ has a nonzero stable class.
Since the $E$-action on $X$ is given by $\chi$, it satisfies
\[
 \widehat\psi(\ell\otimes x)=\psi(\ell)x
 =\big(m_0\beta_0(\ell)\big)\cdot x
 \qquad(\ell\in L_h,\ x\in X).
\]
Thus $\widehat\psi$ is the composite
\[
 L_h\otimes_E X
 \xrightarrow{\beta_0\otimes_E1_X}\widetilde K\otimes_E X
 \xrightarrow{m_0\otimes_E1_X}E\otimes_E X
 \xrightarrow{\sim}X,
\]
where the last map is given by multiplication.
It follows that $\beta_0\otimes_E1_X$ also has a nonzero stable class.

Finally, we compose $\beta[-1]$ with the fixed isomorphism $\overline K[-1]\simeq Y$.
The resulting stable morphism $U_h\to Y$ has an ordinary $E$-bimodule representative $g_h$, since stable Hom is ordinary Hom modulo projective factorizations. By the natural evaluation and shift
comparisons, the stable class of $g_h\otimes_E1_X$ is the composite
\[
 U_h\otimes_E X\simeq(L_h\otimes_E X)[-1]
 \xrightarrow{\ \underline{\beta_0\otimes_E1_X}\,[-1]\ }
 (\widetilde K\otimes_E X)[-1]\simeq (\overline K\otimes_E X)[-1]\simeq Y\otimes_E X.
\]
The middle map is a shift of the nonzero stable class of
$\beta_0\otimes_E1_X$, so this composite is nonzero.
Precomposing with the fixed identification $\iota_h^{-1}$ therefore gives a nonzero element of $W^0$.
\end{proof}

Fix $0\ne w\in W^0$.  For each $h\in k^\times$, we choose $g_h$ as
in Proposition~\ref{lift:nonzero-lifts}. After evaluation at $X$ and
precomposition with $\iota_h^{-1}$, we have
\[
 \underline{g_h\otimes_E1_X}\circ\iota_h^{-1}=c_hw
\]
for a unique $c_h\in k^\times$, since $\dim_k(W^0)=1$.
Replacing $g_h$ by $c_h^{-1}g_h$ normalizes this class to $w$.
In particular, taking $h=1,q$ gives ordinary maps
$g_1:U_1\to Y$ and $g_q:U_q\to Y$ with the same nonzero stable evaluation. All choices and rescalings
take place over $k$.

\section{The symmetric algebra and its self-orthogonal module}
\label{app:section}

We now apply the abstract criteria to the data from
Sections~\ref{seed:section}--\ref{lift:construction-section}. This
proves Theorem~\ref{general:construction}, including the quotient
identifications that recover the original input. We then prove that
the resulting module has complexity three. Finally, we construct
two-term dg models for derived recollements associated with
self-orthogonal generators and determine an explicit four-dimensional
formal model for our examples.

\subsection{The cone module and its complete resolution}
Retain the fixed parameter $q$ and the twisted $E$-bimodules $U_1,U_q$.

Take the normalized ordinary $E$-bimodule maps $g_1,g_q$ from
Subsection~\ref{lift:nonzero-section}. For $h=1,q$, let
$\iota_h:U_h\otimes_E X\to X$ be the fixed source
identification. Their stable evaluations
$\underline{g_h\otimes_E1_X}\circ\iota_h^{-1}: X\to Y\otimes_EX$ agree and are nonzero.

\paragraph{The kernel bimodule.}
We use the two normalized lifts to construct the bimodule in
Proposition~\ref{frame:factory}. Choose a projective cover
$\pi_Q:Q\twoheadrightarrow Y$ in $E^e\modcat$, and define
\begin{equation}\label{app:F}
 F=\ker\bigl((g_1,-g_q,\pi_Q):U_1\oplus U_q\oplus Q\to Y\bigr).
\end{equation}
Since $Y$ is a separately projective $E$-bimodule, the
defining sequence splits after forgetting either one of the two
actions. Consequently $F$ is projective on each side.

We verify the hypotheses of Proposition~\ref{frame:factory}.
The algebra $E$ is symmetric, $X$ is nonprojective, and
Proposition~\ref{ind:Tate-and-middle}\textup{(\ref{ind:Tate-self})} gives $H^0=k$ and
$H^{-3}=H^{-2}=0$. Proposition~\ref{ind:Y-evaluation}\textup{(\ref{ind:model-evaluation})}
gives the entire support of $W$.
This verifies assumption~\textup{(\ref{frame:factory-support})}. The maps $g_1,g_q$
supply the two normalized lifts in assumption~\textup{(\ref{frame:factory-lifts})}.
For assumption~\textup{(\ref{frame:factory-actions})},
Proposition~\ref{ind:weights} gives the actions in the following table:
\begin{center}
\begin{tabular}{@{}llll@{}}
\hline
Degree $a$ & $H^a$ & $T_1^a$ & $T_q^a$\\
\hline
$5m$, $m\ge1$ & $k$ & $1$ & $q^{-m}$\\
$-5m-1$, $m\ge0$ & $k$ & $1$ & $q^{m+1}$\\
$0$ & $k$ & $1$ & $1$\\
\hline
\end{tabular}
\end{center}
The difference $T_1^a-T_q^a$ is invertible in every
nonzero degree in the support, since $q$ has infinite order.

\paragraph{The cone module.}

To form the cone, we need a stable map $X\to F\otimes_E X$
normalized by the two projections. Let $r_1:F\to U_1$ and
$r_q:F\to U_q$ be the restrictions of the direct-sum projections in the defining
sequence of $F$. By Proposition~\ref{frame:factory}\textup{(\ref{frame:factory-normalized-class})},
there is a unique stable morphism $v_0:X\to F\otimes_E X$ such that
\[
 \iota_h\circ\underline{r_h\otimes_E1_X}\circ v_0=1_X
 \qquad(h=1,q).
\]
Choose an ordinary representative
$\widetilde v:X\to F\otimes_EX$. The remaining input for the
cone is an embedding of $X$ into a projective-injective $E$-module;
we will use $Ep$. Recall that $X=k_\chi$ and
$p=f_*\otimes f\in B\subseteq E$. We use the element $\zeta=(0, \chi_B)$
fixed in Section~\ref{lift:construction-section}.

Since $\chi_B(p)=1$, we have $\zeta p=\zeta\ne0$. The character
identity from Section~\ref{lift:construction-section} therefore gives
a nonzero $E$-module homomorphism $j_X:X\to Ep$, $1\mapsto\zeta$.
Since $X$ is simple, $j_X$ is injective.

Using this embedding, define
\begin{equation}\label{app:Z}
 \Lambda=\begin{pmatrix}E&0\\F&E\end{pmatrix},\qquad
 Z=\left(X,
 \frac{F\otimes_EX\oplus Ep}
 {\{(\widetilde v(x),-j_X(x)):x\in X\}},\alpha\right),
\end{equation}
where $\alpha$ is induced by inclusion of the first summand:
$\alpha(y)=[(y,0)]$ for $y\in F\otimes_E X$.
For the embedding $i:X\hookrightarrow J_X$ in
Subsection~\ref{frame:cone-section}, we take $i=j_X$ and $J_X=Ep$.
By Theorem~\ref{frame:cone}\textup{(\ref{frame:cone-properties})},
the signed cone of the two column complexes gives a complete
resolution $P_Z^\bullet$. Thus $Z$ is a nonprojective
Gorenstein-projective $\Lambda$-module. Proposition~\ref{frame:factory}\textup{(\ref{frame:factory-nonexceptional})--(\ref{frame:factory-exceptions})}
verifies the hypotheses of
Theorem~\ref{frame:cone}\textup{(\ref{frame:cone-support})}, which gives
\[
 H^a\Hom_\Lambda^\bullet(P_Z^\bullet,Z)\simeq
 \begin{cases}k,&a=-1,0,\\0,&a\notin\{-1,0\}.
 \end{cases}
\]

We also record the self-injective dimensions of $\Lambda$.

\begin{Rem}\label{app:gorenstein}
The algebra $\Lambda$ is $1$-Gorenstein, with
$\id_\Lambda(\Lambda)=\id_{\Lambda^{\mathrm{op}}}(\Lambda)=1$.
Indeed, since $E$ is symmetric and $F$ is separately projective,
\cite[Proposition~3.7(2)]{Wang} gives
$\id_\Lambda(\Lambda)\le1$ and
$\id_{\Lambda^{\mathrm{op}}}(\Lambda)\le1$.
It remains to show that neither the left nor the right regular module is injective.

First note that $Y$ is not projective over $E^e$.
Indeed, otherwise $Y\otimes_E X$ would be a projective left $E$-module,
forcing $W^0=\underline{\Hom}_E(X,Y\otimes_E X)=0$,
contrary to $W^0\simeq k$.
Thus the projective cover $Q\twoheadrightarrow Y$ has nonzero kernel,
so $\dim_k(Q)>\dim_k(Y)$. The defining sequence of $F$ now gives
\[
 \dim_k(F)=2\dim_k(E)+\dim_k(Q)-\dim_k(Y)>2\dim_k(E)>0.
\]
The nonsplit exact sequence
$0\to(0,F,0)\to(E,F,1_F)\to(E,0,0)\to0$
with projective left term shows that ${}_\Lambda\Lambda$ is not injective.
Applying the same argument to
$\Lambda^{\rm op}$, with the two diagonal idempotents interchanged,
rules out injectivity of the right regular module.
The upper bounds therefore give the asserted equalities.
\end{Rem}

\subsection{The general construction theorem}

The required cone module is now available. We can therefore complete
the transfer and recover the original seed from the resulting pair.

\begin{Theo}
\label{general:construction}
Let $k$ be a field containing an element $q\in k^\times$ of infinite
multiplicative order. Let $H$ be a finite-dimensional Frobenius
$k$-algebra and let
$M_0$ be an indecomposable nonprojective left
$H$-module such that
\begin{equation*}
 \tau_HM_0\simeq M_0,\qquad
 \underline{\End}_H(M_0)\simeq k,\qquad
 \Ext_H^i(M_0,M_0)=0\quad(i\ge2).
\end{equation*}
Set $C=\End_H(H\oplus M_0)^{\mathrm{op}}$ and $S=C/CeC$, where $e,f\in C$ are the idempotents corresponding to the direct summands $H, M_0$, respectively.
Fix the signed auxiliary seed $(C_*,S_*,\gamma_*)$ of Section~\ref{seed:section},
constructed with this parameter $q$, and
form
\[
 B=C_*\otimes_k C,\qquad E=B\ltimes DB,\qquad
 X=S_*\otimes_k S,\qquad \theta=\gamma_*\otimes1_C.
\]
Choose the kernel bimodule $F$ and the module $Z$ by
\eqref{app:F} and \eqref{app:Z}, using the twists $U_1,U_q$
determined by $\theta$ as in Section~\ref{lift:construction-section}, and put
\begin{equation*}
 \Lambda=\begin{pmatrix}E&0\\F&E\end{pmatrix},\qquad
 A=\Lambda\ltimes D\Lambda,\qquad M=A\otimes_\Lambda Z.
\end{equation*}
Then the following statements hold.
\begin{enumerate}
\item\label{general:construction-module} The algebra $A$ is finite-dimensional and symmetric. The
$A$-module $M$ is finite-dimensional and nonprojective, has complexity
three, and satisfies
\[
 \Ext_A^i(M,M)=0\quad(i>0),\qquad
 \underline{\End}_A(M)\simeq k[\varepsilon]/(\varepsilon^2).
\]
\item\label{general:construction-recovery} There is a specified algebra epimorphism $\pi:A\twoheadrightarrow C$
such that $C\otimes_A M\simeq S$.
Writing $I=\ker(\pi)$, the input and output are related by
\begin{equation*}
 A/I\simeq C,\qquad M/IM\simeq S,\qquad
 e(A/I)e\simeq H,\qquad e(A/I)f\simeq M_0.
\end{equation*}
The last isomorphism is an isomorphism of left $H$-modules, using the
preceding identification of the corner algebra.
\item\label{general:construction-basic} If $H$ is connected and split basic with $r$ isomorphism classes
of simple modules, then $A$ is connected and split basic with
$4(r+1)$ isomorphism classes of simple modules.
\end{enumerate}
\end{Theo}
\begin{proof}
\textup{(\ref{general:construction-module})} The preceding subsection verifies the hypotheses of
Corollary~\ref{frame:combined}, which gives the symmetry of $A$
and the nonprojectivity and self-orthogonality of $M$.
Both are finite-dimensional by construction.
To compute the stable endomorphism algebra of $M$, we use the
same corollary, which gives
$H^a\Hom_\Lambda^\bullet(P_Z^\bullet,Z)\simeq k$ for $a=-1,0$.
Theorem~\ref{frame:transfer}\textup{(2)}, applied with $R=\Lambda$, gives
$\underline{\End}_A(M)\simeq k[\varepsilon]/(\varepsilon^2)$.
The complexity assertion is proved in Proposition~\ref{general:complexity}
below.

\textup{(\ref{general:construction-recovery})} To recover $S$ from $M$, we project to the first
diagonal corner algebra of $\Lambda$, identified with $E$, and apply
the character of the fixed factor. Let $\chi_*:C_*\to k$ be the
character defining $S_*$, and let $\rho_0:\Lambda\to E$ be this
projection.
Define $\pi$ as the composite of algebra epimorphisms
\begin{equation*}
 A\twoheadrightarrow\Lambda
 \xrightarrow{\rho_0} E
 \twoheadrightarrow B=C_*\otimes_k C
 \xrightarrow{\chi_*\otimes1_C} k\otimes_kC\simeq C.
\end{equation*}
The first and third maps are the natural trivial-extension
projections.  All four maps are unital and surjective.
Write $\varepsilon=\operatorname{diag}(1_E,0)\in\Lambda$.
As a right $\Lambda$-module via $\rho_0$, $E$ is $\varepsilon\Lambda$,
so $E\otimes_\Lambda Z\simeq\varepsilon Z=X$ as $E$-modules.
Since the $E$-module structure of $X$ is obtained by restriction along $E\twoheadrightarrow B$,
$C\otimes_E X\simeq C\otimes_B X$ as $C$-modules.
Associativity of ordinary tensor products consequently gives
$$
C\otimes_A M\simeq C\otimes_\Lambda Z\simeq C\otimes_E(E\otimes_\Lambda Z)
\simeq C\otimes_E X\simeq C\otimes_B X.
$$
Set $J=\ker(\chi_*)$. Then $C\simeq (C_*\otimes_k C)/(J\otimes_k C)$ as algebras. It follows from $S_*=k_{\chi_*}$ that
$$
C\otimes_B X\simeq X/((J\otimes_kC)X)\simeq (S_*/(JS_*))\otimes_kS\simeq k\otimes_kS\simeq S.
$$
Thus $C\otimes_A M\simeq S$. This proves the asserted quotient-module identification and the two
quotient isomorphisms. The seed construction identifies
$eCe\simeq\End_H(H)^{\mathrm{op}}\simeq H$ by right multiplication.
Moreover,  $eCf\simeq\Hom_H(H,M_0)\simeq M_0$.
These identifications are compatible with $A/I\simeq C$ and prove
the remaining isomorphisms in \textup{(\ref{general:construction-recovery})}. The idempotents $e,f\in C$ lift to
$\Lambda\subseteq A$ by tensoring with $1_{C_*}$ and embedding into the
first diagonal corner algebra of $\Lambda$.

\textup{(\ref{general:construction-basic})} We track connectedness and the semisimple quotients
through $C$, $B$, $E$, $\Lambda$ and $A$.
Assume that $H$ is connected and split basic with $r$ simple modules.
Proposition~\ref{general:seed} shows that $C$ is connected
and split basic with $r+1$ simple modules.

The fixed algebra $C_*$ is connected and split basic with two simple
modules, so $B$ is split basic with $2(r+1)$ simple modules.
To show that $B$ is connected, it suffices to prove that its centre
is local.

We first show a general result: for a connected and split basic $k$-algebra $R$, its centre is local with residue field $k$. Indeed, connectedness excludes nontrivial idempotents in the centre, and a finite-dimensional commutative algebra without such idempotents is local. Thus the centre $\mathrm Z(R)$ of $R$ is local. Since $\mathrm Z(R)\cap\rad(R)=\rad(\mathrm Z(R))$, the residue field $\kappa$ of the centre embeds in $R/\rad(R)\simeq k^s$ for some $s\geq 1$. Each simple $R$-module $V$ is one-dimensional. The central action on $V$ factors through $\kappa$ and induces a unital $k$-algebra homomorphism $\lambda_V:\kappa\to\End_R(V)\simeq k$. This map is injective, so $\kappa\simeq k$.

By the above result, the centres of $C_*$ and $C$ are local algebras with residue field $k$. Moreover,
$\mathrm Z(B)=\mathrm Z(C_*)\otimes_k\mathrm Z(C)$.
The sum of the two radical tensor ideals is nilpotent and has
quotient $k$. Thus $\mathrm Z(B)$ is local with residue field $k$ and $B$ is connected.


We next pass from $B$ to $E$. For any algebra $\Gamma$, the square-zero
ideal $D\Gamma$ in $\Gamma\ltimes D\Gamma$ lies in the radical and
$(\Gamma\ltimes D\Gamma)/\rad(\Gamma\ltimes D\Gamma)\simeq \Gamma/\rad(\Gamma)$ as algebras.
Trivial extension also preserves connectedness. A central idempotent projects to a central idempotent of $\Gamma$. If $\Gamma$ is connected, that image is zero or one, so either the idempotent or
its complement lies in the square-zero ideal $D\Gamma$ and must be zero.
In particular, $E$ is connected and split basic with $2(r+1)$ simple modules.

For $\Lambda$, the condition $F\ne0$ from Remark~\ref{app:gorenstein} forces the two diagonal entries of any central idempotent to agree. Indeed, a central idempotent commutes with the two diagonal idempotents, so it has the form $\operatorname{diag}(c_0,c_1)$, with $c_0,c_1$ central idempotents of $E$. Since $E$ is connected, each $c_i$ is zero or one. Centrality further requires $c_1u=uc_0$ for every $u\in F$.
The condition $F\ne0$ therefore forces $c_0=c_1$, leaving only
the zero and identity idempotents. Thus $\Lambda$ is connected.
It remains to count the simple modules. We have
\[
 \rad(\Lambda)=\begin{pmatrix}\rad(E)&0\\F&\rad(E)\end{pmatrix},\qquad
 \Lambda/\rad(\Lambda)\simeq(E/\rad(E))\times(E/\rad(E)).
\]
Hence $\Lambda$ is split basic with $4(r+1)$ simple modules.  Passing to
$A=\Lambda\ltimes D\Lambda$ preserves both assertions.
\end{proof}

We record a few remarks on the construction and its consequences.

\begin{Rem}
The choice of the finite bimodule model $Y$ may affect the dimensions
of the resulting algebra and module. For fixed $E$, $Y$ and $Q$,
the dimension of $F$, and hence that of $A$, is independent of the
normalized lifts. The recovery relations in
Theorem~\ref{general:construction}\textup{(\ref{general:construction-recovery})} hold for every such choice; no
minimality is asserted. Since $\dim_k(C_*)=10$, the dimensions at the
last stages are
\[
 \dim_k(E)=20\dim_k(C),\qquad
 \dim_k(\Lambda)=2\dim_k(E)+\dim_k(F),\qquad
 \dim_k(A)=4\dim_k(E)+2\dim_k(F).
\]
Under the hypotheses of Theorem~\ref{general:construction}\textup{(\ref{general:construction-basic})},
connectedness and the number of simple modules are independent of
these choices.
\end{Rem}

The projective bimodule model for $\mathcal K$ must be unbounded to the left.
To see this, we show that the derived bimodule $\mathcal K$ is not perfect
over $E^e$. Otherwise its stable image $\overline K$ would vanish, so
$Y\simeq \overline K[-1]$ would be an $E^e$-projective module. Then
$Y\otimes_E X$ would be $E$-projective, contradicting $W^0=k$.
On each side, the derived bimodule has a bounded projective model.

It remains to identify the ground fields admitting the required
parameter. In characteristic zero, $q=2$ has infinite multiplicative order.
In positive characteristic, a field not algebraic over its prime field
contains a transcendental element, which also has infinite order.
In a field algebraic over a finite prime field,
every nonzero element lies in a finite subfield and has finite order.
The hypotheses here therefore impose no condition on the
characteristic but require the specified field parameter.

The resulting self-orthogonal module also yields counterexamples to
related homological conjectures.

\begin{Rem}
\label{audit:section}
The pair $(A,M)$ in Theorem~\ref{general:construction} also violates
the \emph{Auslander--Reiten conjecture}: since $A$ is self-injective,
$\Ext_A^i(M,M\oplus A)=0$ for every $i>0$, although $M$ is
nonprojective. As every $A$-module is Gorenstein-projective, this
also gives a counterexample to the \emph{Gorenstein-projective conjecture}.

Moreover, put $G=A\oplus M$ and $\Gamma=\End_A(G)^{\mathrm{op}}$.
M\"uller's theorem on dominant dimension \cite[Lemma~3]{Mueller} gives
$\operatorname{domdim}(\Gamma)=\infty$.
To see that $\Gamma$ is not self-injective, we use symmetry of $A$ to choose
an embedding $M\hookrightarrow P$ with $P\in\add(A)$ and apply
$\mathcal H=\Hom_A(G,-)$. If ${}_\Gamma\Gamma$ were injective, the projective module
$\mathcal H(M)$ would be injective, so the monomorphism
$\mathcal H(M)\hookrightarrow\mathcal H(P)$ would split.
Full faithfulness of $\mathcal H$ on $\add(G)$ would then make $M\hookrightarrow P$ split,
contradicting nonprojectivity of $M$. Thus $\Gamma$ violates the \emph{Nakayama conjecture}.

Let $W$ be the direct sum of representatives of the indecomposable
projective-injective left $\Gamma$-modules. Since
$\operatorname{domdim}(\Gamma)=\infty$, the minimal injective coresolution
of $\Gamma$ has all terms in $\add(W)$ and remains exact under
$\Hom_\Gamma(-,W)$. Thus $W$ is Wakamatsu tilting with $\pd_\Gamma(W)=0$.
If it were tilting, then $\add(W)=\add(\Gamma)$, contradicting the
non-self-injectivity of $\Gamma$.
Hence $W$ is not tilting, giving a counterexample to the
\emph{Wakamatsu tilting conjecture}.

For an algebra $R$, the \emph{little} and \emph{big finitistic dimensions}
$\operatorname{findim}(R)$ and $\operatorname{Findim}(R)$ are the suprema of the
finite projective dimensions of finitely generated and arbitrary left
$R$-modules, respectively. The \emph{little and big finitistic dimension
conjectures} assert, respectively, that these dimensions are finite for every
algebra. By the \emph{finitistic dimension conjecture}, we mean the little version.

The terms of the minimal injective coresolution $0\to\Gamma\to I^0\to I^1\to\cdots$ are projective.
Put $C_0={}_\Gamma\Gamma$ and define the successive cosyzygies by
$0\to C_j\to I^j\to C_{j+1}\to 0$ for $j\ge0$. Then $\pd_\Gamma(C_n)\le n$ for $n\ge1$.
The sequence $0\to\Gamma\to I^0\to C_1\to0$ does not split,
since $\Gamma$ is not self-injective. Its class is therefore nonzero in
$\Ext_\Gamma^1(C_1,\Gamma)$.
Since each $I^j$ is projective, dimension shifting gives
$0\ne\Ext_\Gamma^1(C_1,\Gamma)\simeq\Ext_\Gamma^n(C_n,\Gamma)$
for $n\ge1$. Hence $\pd_\Gamma(C_n)=n$ for every $n\ge1$, and
$\operatorname{Findim}(\Gamma)\ge\operatorname{findim}(\Gamma)=\infty$.
This gives counterexamples to the \emph{little and big finitistic dimension
conjectures}. See also \cite{EM} and \cite[Section~8]{F} for further connections among these homological conjectures.
\end{Rem}

\Needspace{12\baselineskip}
\subsection{Complexity of the constructed module}
\label{general:complexity-section}

The notion of complexity originates in Alperin's work on modular
representations of finite groups \cite{Alperin}. We use its usual
formulation for finite-dimensional algebras.

Let $R$ be such an algebra,
let $U$ be a finite-dimensional $R$-module, and let $P_U^\bullet\to U$
be a minimal projective resolution concentrated in nonpositive degrees.
The \emph{complexity} $\operatorname{cx}_R(U)$ is the least integer
$d\ge0$ for which there is a constant $c>0$ such that
$\dim_kP_U^{-n}\le c(n+1)^{d-1}$ for all sufficiently large $n$;
we set $\operatorname{cx}_R(U)=\infty$ if no such integer exists.
Thus $\operatorname{cx}_R(U)=0$ means that $U$ has finite projective
dimension, while $\operatorname{cx}_R(U)\le1$ means that the dimensions
of the terms in its minimal projective resolution are bounded.
For self-injective $R$, complexity zero is equivalent to projectivity.

\begin{Prop}\label{general:complexity}
For the algebras $B,E,A$ and modules $X,M$ defined in
Theorem~\ref{general:construction}, we have
$\operatorname{cx}_B(X)=2$ and
\mbox{$\operatorname{cx}_E(X)=\operatorname{cx}_A(M)=3$}.
\end{Prop}
\begin{proof}
We first estimate the dimensions of the terms in minimal projective
resolutions of $X$ over $B$ and $E$. We then construct a projective
$A$-resolution of $M$ to obtain $\operatorname{cx}_A(M)\le3$.
For the reverse inequality, we recover a projective $E$-resolution
of $X$ from a minimal projective $A$-resolution of $M$, without
increasing the dimensions of its terms.

\smallskip\noindent
\emph{Step 1. We prove that $\operatorname{cx}_B(X)=2$ by establishing linear growth of the projective terms.}
\par\nobreak\smallskip\nobreak
\noindent
We first show that the terms in the minimal projective resolutions
of the two seed modules are nonzero and have bounded dimensions.
Retain $N=H\oplus M_0$, and choose the minimal projective resolution
$\cdots\to P_1\to P_0\to M_0\to0$ used in the proof of
Proposition~\ref{general:seed}. Since
$\tau_HM_0\simeq\Omega_H^2\nu_HM_0\simeq M_0$ and $\nu_H$
commutes with syzygies, we have
$\Omega_H^{n+2}M_0\simeq\nu_H^{-1}\Omega_H^nM_0$ for $n\ge0$.
The Nakayama functor preserves minimal projective covers, so
$P_{n+2}\simeq\nu_H^{-1}P_n$. As $H$ is Frobenius, $\nu_H$ is
an automorphism twist and preserves vector-space dimensions.
Thus the dimensions of the $P_n$ are bounded. Every $P_n$ is nonzero:
otherwise $M_0$ would be projective by the self-injectivity of $H$.

By the proof of Proposition~\ref{general:seed}\textup{(\ref{general:seed-homology})},
the projective $C$-resolution of $S$ has degree-zero term $Cf$ and degree-$(-n-1)$ term
$\Hom_H(N,P_n)$ for $n\ge0$. This resolution is minimal. Indeed, the functor $\Hom_H(N,-):\add(N)\to C\text{-}\mathrm{proj}$ is an equivalence and preserves radical morphisms. The differentials
between the $P_n$ are radical, and the components of $P_0\to M_0$
are also radical since $M_0$ is indecomposable and nonprojective.
Moreover, $\Hom_H(H,P_n)\simeq P_n$ is a direct summand of
$\Hom_H(N,P_n)$ as a vector space, whereas
$\dim_k\Hom_H(N,P_n)\le(\dim_kN)(\dim_kP_n)$.
Thus the terms of this minimal resolution have dimensions bounded
above and below by positive constants. The same argument applies
to the fixed seed $S_*$ over $C_*$.

We now tensor the two seed resolutions over $k$ to resolve $X$ over
$B=C_*\otimes_kC$. To check minimality, note that $\rad(B)=\rad(C_*)\otimes_kC+C_*\otimes_k\rad(C)$.
Indeed, the ideal on the right is nilpotent.
Since $C_*$ is split basic with two simple modules, the quotient of $B$ by this ideal is
$k^2\otimes_k(C/\rad(C))\simeq(C/\rad(C))^2$, which is semisimple.
Each differential in the tensor total complex therefore has image
in the radical of its target, so this complex is a minimal
projective $B$-resolution $P_X^\bullet\to X$.
Its degree-$(-n)$ term is the direct sum of $n+1$ tensor products
of seed terms. Its dimension is consequently bounded above and
below by positive constant multiples of $n+1$. This proves
$\operatorname{cx}_B(X)=2$.

To compare the resolutions over $B$ and $E$, we will measure their
terms using extension groups with a semisimple target. Set
$T=E/\rad(E)\simeq B/\rad(B)$ and define
$b_n=\dim_k\Ext_B^n(X,T)$ for $n\ge0$.
Then $$b_n=\dim_k\Hom_B(P_X^{-n},T)=\dim_k\Hom_T(P_X^{-n}/\rad(B)P_X^{-n},T)=\dim_k(P_X^{-n}/\rad(B)P_X^{-n}),$$
where the last equality uses the symmetry of the semisimple algebra $T$. By Nakayama's lemma, lifts of a $k$-basis of this quotient generate $P_X^{-n}$, so $b_n\le\dim_kP_X^{-n}\le(\dim_kB)b_n$.
The linear growth of $P_X^{-n}$ therefore gives constants
$c_1,c_2>0$ such that $c_1(n+1)\le b_n\le c_2(n+1)$.

\smallskip\noindent
\emph{Step 2. We prove that $\operatorname{cx}_E(X)=3$ by establishing quadratic growth of the projective terms.}
\par\nobreak\smallskip\nobreak
\noindent
We relate $\Ext_E^n(X,T)$ to the groups $\Ext_B^n(X,T)$ by using the triangle $L\to J\to X\to L[1]$
from Proposition~\ref{ind:one-generator}, where
$L\simeq X[4]$ and $J=E\otimes_B^{\mathbf L}X$.
On the projective model $P_X^\bullet$, its first morphism is
represented by the inclusion
$DB\otimes_BP_X^\bullet\to E\otimes_BP_X^\bullet$.
The latter complex is bounded above with projective $E$-terms,
so every derived morphism $J\to T[n]$ is represented by a cochain map
$E\otimes_BP_X^\bullet\to T[n]$. Every such cochain map vanishes on
$DB\otimes_BP_X^\bullet$ by $E$-linearity, because $DB$ annihilates $T$.
Hence the induced map
$\Hom_{\mathscr D^b(E)}(J,T[n])\to
\Hom_{\mathscr D^b(E)}(L,T[n])$ is zero for every integer $n$.
Applying $\Hom_{\mathscr D^b(E)}(-,T[n])$ to the triangle, using
these zero maps in degrees $n-1$ and $n$, and applying
induction--restriction adjunction gives
\[
 0\longrightarrow\Ext_E^{n-5}(X,T)
 \longrightarrow\Ext_E^n(X,T)
 \longrightarrow\Ext_B^n(X,T)\longrightarrow0
 \qquad(n\ge0).
\]
Here negative ordinary Ext groups are zero, as in our conventions.
Write $e_n=\dim_k\Ext_E^n(X,T)$ for $n\ge0$, and set $e_n=0$
for $n<0$. The short exact sequence gives $e_n=e_{n-5}+b_n$.
Let $r$ be the largest nonnegative integer satisfying $5r\le n$.
Iterating this recurrence and using $e_{n-5(r+1)}=0$, we obtain
\[
e_n=\sum_{j=0}^{r}b_{n-5j}\qquad(n\ge0).
\]
By Step 1, for each $0\le j\le r$ we have
$c_1(n-5j+1)\le b_{n-5j}\le c_2(n-5j+1)$.
Note that
\[
\sum_{j=0}^{r}(n-5j+1)=(r+1)(n+1-5r/2).
\]
Since $5r\le n<5r+5$, we have
$(n+1)/5\le r+1\le n+1$ and
$(n+1)/2\le n+1-5r/2\le n+1$. Hence
\[
\frac{(n+1)^2}{10}\le\sum_{j=0}^{r}(n-5j+1)\le(n+1)^2.
\]
Combining these inequalities gives
$(c_1/10)(n+1)^2\le e_n\le c_2(n+1)^2$ for $n\ge0$.
Let $P^\bullet\to X$ be a minimal projective $E$-resolution.
As in Step 1, we have  $e_n\le\dim_kP^{-n}\le(\dim_kE)e_n$ for $n\ge0$.
These bounds prove $\operatorname{cx}_E(X)=3$.

\smallskip\noindent
\emph{Step 3. We transfer the quadratic growth bounds to $M$ and prove that $\operatorname{cx}_A(M)=3$.}
\par\nobreak\smallskip\nobreak
\noindent
For the upper bound, we extend $P^\bullet$ to a complete $E$-resolution
by splicing it with an injective coresolution beginning with the
prescribed embedding $j_X:X\hookrightarrow Ep$.
Writing $\epsilon:P^0\twoheadrightarrow X$ for the augmentation,
we have $P^1=Ep$ and $d_P^0=j_X\epsilon$, as required in
Subsection~\ref{frame:cone-section}. The cokernel description in
Theorem~\ref{frame:cone}\textup{(\ref{frame:cone-properties})}
shows that the resulting cone still represents the specified module $Z$.
Its projective terms satisfy
$P_Z^{-n}=D_0(P^{-n})\oplus D_1(P^{1-n})$ for $n\ge1$.
The column functors $D_0,D_1$ increase dimensions by at most fixed
constant factors. Thus Step 2 gives a constant $c_3>0$, independent
of $n$, such that $\dim_kP_Z^{-n}\le c_3(n+1)^2$ for $n\ge1$.
By the proof of Theorem~\ref{frame:transfer},
$A\otimes_\Lambda P_Z^\bullet$ is a complete projective resolution
of $M$. Since
$\dim_k(A\otimes_\Lambda P_Z^{-n})\le(\dim_kA)\dim_kP_Z^{-n}$,
there is a constant $c_4>0$ such that
$\dim_k(A\otimes_\Lambda P_Z^{-n})\le c_4(n+1)^2$ for $n\ge1$.
A minimal projective resolution is a direct summand, degree by
degree, of this resolution, so $\operatorname{cx}_A(M)\le3$.

For the lower bound, let $P_M^\bullet\to M$ be a minimal projective
$A$-resolution. We will turn it into a projective $E$-resolution of
$X$ without increasing the dimensions of its terms.
Regard $\Lambda$ as a right $A$-module through
$A\twoheadrightarrow\Lambda$. Tensoring the nonpositive part of
$A\otimes_\Lambda P_Z^\bullet$ with $\Lambda$ over $A$ recovers
the nonpositive part of $P_Z^\bullet$. Hence
$\Lambda\otimes_AM\simeq Z$ and
$\operatorname{Tor}_i^A(\Lambda,M)=0$ for $i>0$.
These vanishings ensure that
$\Lambda\otimes_AP_M^\bullet\to Z$ remains exact and is therefore
a projective $\Lambda$-resolution.

We define $\epsilon_0=\operatorname{diag}(1_E,0)\in\Lambda$.
Since $\epsilon_0\Lambda\simeq E$ as left $E$-modules, the exact
functor $\epsilon_0(-)$ sends projective $\Lambda$-modules to
projective $E$-modules. As $\epsilon_0Z=X$, applying this functor
to $\Lambda\otimes_AP_M^\bullet\to Z$ gives a projective $E$-resolution
$\epsilon_0(\Lambda\otimes_AP_M^\bullet)\to X$.
For any $A$-module $U$, the quotient
$\Lambda\otimes_AU\simeq U/(D\Lambda)U$ and its component
$\epsilon_0(\Lambda\otimes_AU)$ have dimension at most $\dim_kU$.
In particular,
$\dim_k\epsilon_0(\Lambda\otimes_AP_M^{-n})\le\dim_kP_M^{-n}$ for $n\ge0$.
Since the nonpositive part of $P^\bullet$ is minimal, its terms
are direct summands of the corresponding terms of this resolution.
Thus
$\dim_kP^{-n}\le\dim_kP_M^{-n}$ for $n\ge0$, and consequently
$3=\operatorname{cx}_E(X)\le\operatorname{cx}_A(M)$.
Together with the upper bound, this proves $\operatorname{cx}_A(M)=3$.
\end{proof}

\begin{Rem}\label{general:bounded-selforthogonal}
The modules in our counterexamples to (TC2) have finite complexity
by Proposition~\ref{general:complexity}. We now explain why a
counterexample cannot have complexity at most one.

More generally, for a finite-dimensional self-injective algebra $R$, every
finite-dimensional self-orthogonal module $U$ with
$\operatorname{cx}_R(U)\le1$ is projective. This is a consequence
of Dugas's periodicity criterion \cite[Theorem~1]{Dugas}.
To apply this criterion, we first extend scalars to an algebraic closure:
self-injectivity,
self-orthogonality and the existence of a projective resolution with
bounded terms are preserved, and projectivity descends along a field
extension. We may also pass to a Morita equivalent basic algebra,
which preserves these properties.
If $U$ had a nonprojective indecomposable summand $V$, then
$\operatorname{cx}_R(V)=1$ and $\Ext_R^1(V,V)=0$.
The latter condition implies that $V$ has an open orbit in its module
variety, by the standard tangent-space criterion.
Dugas's theorem makes $V$ periodic, say $\Omega_R^rV\simeq V$
for some $r>0$. This contradicts
$0=\Ext_R^r(V,V)\simeq\underline{\Hom}_R(\Omega_R^rV,V)
\simeq\underline{\End}_R(V)\ne0.$
\end{Rem}

\subsection{Two-term dg models and derived recollements}
\label{general:dg-models}

For a symmetric algebra $A$ and a self-orthogonal generator
$G=A\oplus M$, we compare the derived categories of $A$ and
$\End_A(G)^{\mathrm{op}}$. A recollement describes the difference
between them by the derived category of a two-term $2$-symmetric dg
algebra. For our counterexamples, this dg algebra is formal and
admits an explicit four-dimensional model.

For a dg algebra $T$, let $\mathscr D(T)$ denote the unbounded
derived category of all left dg $T$-modules, without any
finite-dimensionality restriction. A dg algebra morphism
$R\to T$ is a \emph{homological epimorphism} if the multiplication
$T\otimes_R^{\mathbf L}T\to T$ is a quasi-isomorphism, or equivalently
if derived restriction is fully faithful.

The following proposition gives an explicit two-term form of
derived quotient results; compare Braun--Chuang--Lazarev
\cite[Remark~9.5]{BCL} and Booth \cite[Corollary~3.2.4]{BoothDQ}.
We include a self-contained proof for completeness.
For the definition of recollements of triangulated categories,
we refer to \cite[Section~1.4.3]{BBD}.

\begin{Prop}\label{dg:two-term}
Let $R$ be a finite-dimensional $k$-algebra and $e\in R$ an idempotent.
Put $\Delta_e=Re\otimes_{eRe}eR$ and let $\mu_e:\Delta_e\to R$ be
multiplication.
\begin{enumerate}
\item There is a two-term dg algebra
$\mathcal B_e=(\Delta_e\xrightarrow{\mu_e}R)$ in degrees $-1,0$,
with the natural $R$-bimodule actions and square-zero degree-$(-1)$ part.
Write $\lambda_e:R\to\mathcal B_e$ for the inclusion.
\item The map $\lambda_e$ is a homological
epimorphism if and only if
$\operatorname{Tor}^{eRe}_i(Re,eR)=0$ for all $i>0$.
Under these conditions, there is a recollement of derived categories
\[
\begin{tikzcd}[column sep=5em]
 \mathscr D(\mathcal B_e) \arrow[r]
 & \mathscr D(R)
   \arrow[l, shift right=3ex]
   \arrow[l, shift left=3ex]
   \arrow[r]
 & \mathscr D(eRe)
   \arrow[l, shift right=3ex]
   \arrow[l, shift left=3ex]
\end{tikzcd}
\]
where the left embedding is restriction along $\lambda_e$ and
the right quotient functor is $e(-)$.
\end{enumerate}
\end{Prop}
\begin{proof}
\textup{(1)} To verify the dg structure, we use the balancing identity
$\mu_e(\delta)\delta'=\delta\mu_e(\delta')$ for
every $\delta,\delta'\in\Delta_e$. This is the Leibniz identity in
degree $(-1,-1)$; the other cases follow from bimodule linearity.
We will also use that $\mathcal B_e e$ and $e\mathcal B_e$
are contractible, since their differentials identify with the
identity maps of $Re$ and $eR$.

\textup{(2)} To prove the homological epimorphism criterion, we apply
$\mathcal B_e\otimes_R^{\mathbf L}-$ to the defining triangle
\[
 \Delta_e\lraf{\mu_e}R\lraf{\lambda_e}\mathcal B_e
 \longrightarrow\Delta_e[1].
\]
The resulting map
$\mathcal B_e\to\mathcal B_e\otimes_R^{\mathbf L}\mathcal B_e$
composes with multiplication to the identity. Thus $\lambda_e$
is a homological epimorphism precisely when
$\mathcal B_e\otimes_R^{\mathbf L}\Delta_e=0$.
To test this vanishing, put
$\mathcal L_e=Re\otimes_{eRe}^{\mathbf L}eR$ and let
$u:\mathcal L_e\to\Delta_e$ be the natural comparison.
Contractibility gives
$\mathcal B_e\otimes_R^{\mathbf L}\mathcal L_e
\simeq(\mathcal B_e e)\otimes_{eRe}^{\mathbf L}eR=0$.
If the higher Tor groups vanish, then $u$ is a quasi-isomorphism,
so $\mathcal B_e\otimes_R^{\mathbf L}\Delta_e=0$ and
$\lambda_e$ is a homological epimorphism.

Conversely, suppose that $\lambda_e$ is a homological epimorphism,
so $\mathcal B_e\otimes_R^{\mathbf L}\Delta_e=0$.
To prove that $u$ is a quasi-isomorphism, let $\mathcal C_e$ be its cone.
Then
$\mathcal B_e\otimes_R^{\mathbf L}\mathcal C_e=0$.
Since $H^0(u)$ is an isomorphism and $eu$ identifies with the
identity of $eR$, we have $\mathcal C_e\in\mathscr D^{\le-2}(R)$
and $e\mathcal C_e=0$.
It follows that $0=H^n(e\mathcal C_e)=eH^n(\mathcal C_e)$ for every $n$.
If $\mathcal C_e\ne0$, choose the largest $n$ with
$H^n(\mathcal C_e)\ne0$. As $\mathcal B_e$ is nonpositive,
truncation at $n$ gives
\[
 H^n(\mathcal B_e\otimes_R^{\mathbf L}\mathcal C_e)
 \simeq H^0(\mathcal B_e)\otimes_R H^n(\mathcal C_e)
 \simeq (R/ReR)\otimes_R H^n(\mathcal C_e)
 \simeq H^n(\mathcal C_e)\ne0.
\]
This contradiction gives $\mathcal C_e=0$, so $u$ is a
quasi-isomorphism and the higher Tor groups vanish.

To obtain the recollement, we identify the image of restriction.
Under these equivalent conditions, the defining triangle is
$\mathcal L_e\to R\to\mathcal B_e\to\mathcal L_e[1]$.
Tensoring it with $X\in\mathscr D(R)$ satisfying $eX=0$ gives
$X\simeq\mathcal B_e\otimes_R^{\mathbf L}X$.
Together with the contractibility of $e\mathcal B_e$, this identifies
the image of restriction with $\operatorname{Ker}(e(-))$.
The functor $e(-)$ has fully faithful adjoints
$Re\otimes_{eRe}^{\mathbf L}-$ and
$\mathbf R\Hom_{eRe}(eR,-)$; restriction has both adjoints as well.
Hence the recollement follows.
\end{proof}

We next apply this construction to generators over symmetric
algebras, where a dg bimodule duality gives additional structure.

For a dg algebra $T$, we write $DT$ for its graded $k$-dual and
$T^e=T\otimes_kT^{\mathrm{op}}$, using the graded opposite and
Koszul signs. We call $T$ \emph{formal} if it is quasi-isomorphic,
as a dg algebra, to its graded cohomology algebra
$H^*(T):=\bigoplus_{n\in\mathbb Z}H^n(T)$, equipped with zero differential.
Following \cite[Definition~2.2]{Jin}, for a positive
integer $d$, a nonpositive dg algebra $T$ with finite-dimensional total
cohomology is \emph{$d$-symmetric} if $\add(T)=\add(DT[d-1])$ in $\mathscr D(T^e)$;
the analogous equality in $\mathscr D(T)$ defines $d$-self-injectivity.

\begin{Prop}\label{dg:symmetric-application}
Let $A$ be a symmetric algebra and $M$ an $A$-module.
Put $G=A\oplus M$ and $\Gamma=\End_A(G)^{\mathrm{op}}$.
Let $e\in\Gamma$ be the idempotent endomorphism of $G$
that is the identity on $A$ and zero on $M$, and write
$\overline\Gamma=\Gamma/\Gamma e\Gamma
\simeq\underline{\End}_A(M)^{\mathrm{op}}$.
Let $\mathcal B_e$ be the dg algebra of
Proposition~\ref{dg:two-term} with $R=\Gamma$.
\begin{enumerate}
\item The canonical map $\Gamma\to\mathcal B_e$ is a homological
epimorphism if and only if $M$ is self-orthogonal. In this case,
there is a recollement of derived categories
\[
\begin{tikzcd}[column sep=5em]
 \mathscr D(\mathcal B_e) \arrow[r]
 & \mathscr D(\Gamma)
   \arrow[l, shift right=3ex]
   \arrow[l, shift left=3ex]
   \arrow[r]
 & \mathscr D(A)
   \arrow[l, shift right=3ex]
   \arrow[l, shift left=3ex]
\end{tikzcd}
\]
\item Without assuming self-orthogonality, we have
$H^0(\mathcal B_e)\simeq\overline\Gamma$ and
$H^{-1}(\mathcal B_e)\simeq D\overline\Gamma$.
Moreover, $\mathcal B_e\simeq D\mathcal B_e[1]$ in
$\mathscr D((\mathcal B_e)^e)$, so $\mathcal B_e$ is $2$-symmetric
and hence $2$-self-injective.
\end{enumerate}
\end{Prop}
\begin{proof}
\textup{(1)} To apply Proposition~\ref{dg:two-term}, we write
$G^*=\Hom_A(G,A)$, viewed as a right $A$-module.
Then $e\Gamma e\simeq A$ as algebras,
$\Gamma e\simeq G^*$ as right $A$-modules, and
$e\Gamma\simeq G$ as left $A$-modules.
Since $A$ is symmetric, $DG^*\simeq G$ as left $A$-modules. Thus
\[
 D\operatorname{Tor}^{e\Gamma e}_i(\Gamma e,e\Gamma)
 \simeq\Ext_A^i(G,DG^*)
 \simeq\Ext_A^i(G,G)
 \simeq\Ext_A^i(M,M)
 \qquad(i>0).
\]
The last isomorphism uses that $A$ is projective and injective.
Proposition~\ref{dg:two-term} gives the criterion and recollement.

\textup{(2)} For the dg bimodule duality, recall that $\Gamma$ is gendo-symmetric
and $\Gamma e$ is faithful projective-injective. The double centralizer property and
\cite[Lemma~5.1(2)]{CFX} give a $\Gamma$-bimodule isomorphism
$\gamma_e:\Delta_e\xrightarrow{\sim}D\Gamma$.
To define the desired dg bimodule isomorphism, we first construct
a graded $k$-linear map $\tau_e:\mathcal B_e\to k[1]$.
For $\delta\in\Delta_e=\mathcal B_e^{-1}$, put
$\tau_e(\delta)=\gamma_e(\delta)(1_\Gamma)$, and let $\tau_e$
vanish on $\Gamma=\mathcal B_e^0$.
For $a\in\Gamma$ and $\delta\in\Delta_e$, bimodule linearity gives
$\tau_e(a\delta)=\gamma_e(\delta)(a)=\tau_e(\delta a)$.

We then define $\Theta_e:\mathcal B_e\to D\mathcal B_e[1]$ by
\[
 \Theta_e(x)(y)=\tau_e(xy)
 \qquad\bigl(x\in\mathcal B_e^i,\ y\in\mathcal B_e^{-i-1},\ i=-1,0\bigr).
\]
The degree-$(-1)$ component of $\Theta_e$ is $\gamma_e$, while
its degree-zero component is $D\gamma_e$, under the identification
$\Gamma\simeq DD\Gamma$. Both components are isomorphisms.

Associativity and the $\Gamma$-bimodule linearity of $\gamma_e$
show that $\Theta_e$ respects both actions.
It remains to check compatibility with the differentials.
For $\delta,\delta'\in\Delta_e$, this reduces to
$\tau_e(\mu_e(\delta)\delta'-\delta\mu_e(\delta'))=0$,
by the balancing identity in the proof of Proposition~\ref{dg:two-term}.
Thus $\Theta_e$ is an isomorphism of dg bimodules.
Taking cohomology gives
$H^{-1}(\mathcal B_e)\simeq DH^0(\mathcal B_e)=D\overline\Gamma$
as bimodules.
\end{proof}

\Needspace{10\baselineskip}
For our counterexamples, we can make the dg model explicit by
determining its quasi-isomorphism type.

\begin{Koro}\label{dg:four-dimensional}
Let $(A,M)$ be the pair of Theorem~\ref{general:construction}, and
let $\Gamma,e,\mathcal B_e$ be as above. Then $\mathcal B_e$ is
quasi-isomorphic to the zero-differential dg algebra
\[
 \mathcal F=
 k\langle\varepsilon,\eta\rangle/
 (\varepsilon^2,\eta^2,\varepsilon\eta-\eta\varepsilon),
 \qquad |\varepsilon|=0,\quad |\eta|=-1.
\]
In particular, $\mathcal B_e$ is formal, and there is a recollement of derived categories
\[
\begin{tikzcd}[column sep=5em]
 \mathscr D(\mathcal F) \arrow[r]
 & \mathscr D(\Gamma)
   \arrow[l, shift right=3ex]
   \arrow[l, shift left=3ex]
   \arrow[r]
 & \mathscr D(A)
   \arrow[l, shift right=3ex]
   \arrow[l, shift left=3ex]
\end{tikzcd}
\]
\end{Koro}
\begin{proof}
To construct a unital dg quasi-isomorphism $\mathcal F\to\mathcal B_e$,
we first seek a unital section of the quotient map
$\pi_e:\Gamma\twoheadrightarrow\overline\Gamma$.
By Theorem~\ref{general:construction},
$\overline\Gamma\simeq k[\varepsilon]/(\varepsilon^2)$, so such a section
will follow from a strictly square-zero lift of $\varepsilon$.

Recall that $A=\Lambda\ltimes D\Lambda$ and $M=A\otimes_\Lambda Z$.
The proof of Theorem~\ref{frame:transfer}\textup{(2)}, with $R=\Lambda$,
gives $\underline{\Hom}_\Lambda(Z,D\Lambda\otimes_\Lambda Z)\simeq k$.
We choose a $\Lambda$-linear map $\beta:Z\to D\Lambda\otimes_\Lambda Z$
representing a nonzero stable class. Under the decomposition
$\operatorname{Res}_\Lambda M=Z\oplus(D\Lambda\otimes_\Lambda Z)$,
we define $\widetilde\varepsilon:M\to M$ by
$\widetilde\varepsilon(z,v)=(0,\beta(z))$ for
$z\in Z$ and $v\in D\Lambda\otimes_\Lambda Z$.
The same proof shows that $\widetilde\varepsilon$ is $A$-linear
and that its stable class generates the radical of
$\underline{\End}_A(M)$.
Its image lies in the second summand, which it kills. Thus
$\widetilde\varepsilon^2=0$.
Extending $\widetilde\varepsilon$ by zero on the summand $A$ gives
$\widehat\varepsilon=0_A\oplus\widetilde\varepsilon\in\Gamma$.
We now identify $\overline\Gamma$ with $k[\varepsilon]/(\varepsilon^2)$
by sending $\pi_e(\widehat\varepsilon)$ to $\varepsilon$.
Then the $k$-linear map $\sigma_e:\overline\Gamma\to\Gamma$, defined on the basis
$1,\varepsilon$ of $\overline\Gamma$ by
$\sigma_e(1)=1_\Gamma$ and $\sigma_e(\varepsilon)=\widehat\varepsilon$,
is a unital algebra section of $\pi_e$.

To extend $\sigma_e$ to a dg map, we use the symmetrizing form
$\ell\in D\overline\Gamma$, $\ell(a+b\varepsilon)=b$ for $a,b\in k$,
to identify $\mathcal F\simeq\overline\Gamma\ltimes D\overline\Gamma[1]$
via $\varepsilon\mapsto(\varepsilon,0)$ and $\eta\mapsto(0,\ell)$,
with $\ell$ in degree $-1$.
The natural $\overline\Gamma$-bimodule
$K_e=\ker(\mu_e)=H^{-1}(\mathcal B_e)$ admits a $\overline\Gamma$-bimodule isomorphism
$\kappa_e:D\overline\Gamma\xrightarrow{\sim}K_e$
by Proposition~\ref{dg:symmetric-application}\textup{(2)}.
Define $\Phi_e:\overline\Gamma\ltimes D\overline\Gamma[1]\to\mathcal B_e$
by $\Phi_e^0=\sigma_e$ and $\Phi_e^{-1}=\kappa_e$,
viewing $K_e$ as a submodule of $\Delta_e$.
The section property of $\sigma_e$ and bimodule linearity of $\kappa_e$
ensure compatibility with mixed products.
Both negative parts have square zero and $\mu_e(K_e)=0$,
so $\Phi_e$ is a unital dg map.
It induces the identity on $H^0$ and $\kappa_e$ on $H^{-1}$,
hence is a quasi-isomorphism.
Proposition~\ref{dg:symmetric-application}\textup{(1)} gives the recollement.\qedhere
\end{proof}

\begin{Rem}
The algebra $\Gamma$ in Corollary~\ref{dg:four-dimensional} is not compactly Gorenstein in the sense of
\cite[Definition~1.3]{CXSO}, thereby refuting the conjecture in \cite[Section~1.2]{CXSO}.
Following the notation in \cite[Section~4.2]{CXSO}, we take a $G$-filtered
$A$-module $X_\infty$ stably isomorphic to $H_*^0(G)$, where $H_*^0$ denotes the homological functor associated with the torsion pair in $A\text{-}\underline{\mathrm{Mod}}$ determined by ${_A}G$. Let $Y_\infty:=\Hom_A(G,X_\infty)$. The construction in \cite[proof of Corollary~1.4]{CXSO} shows that $Y_\infty$ is a countably generated (but not finitely generated) nonprojective Gorenstein-projective
$\Gamma$-module. It admits an ascending filtration $Y_\infty=\bigcup_{n\ge0}Y_n$ by finitely generated submodules of finite projective dimension, and it is compact in the stable category of all Gorenstein-projective $\Gamma$-modules. Since $A$ is symmetric and $\overline\Gamma\simeq k[\varepsilon]/(\varepsilon^2)$ is Frobenius, it follows from \cite[Corollary~4.11 and Lemma~3.17(2)]{CXSO} that $\Omega_\Gamma^2Y_\infty\simeq Y_\infty$ in this stable category.
\end{Rem}

\appendix
\section{Cyclic coverings of the quantum exterior seed}
\label{family:section}

We check the seed conditions for the cyclic coverings using graded
modules and syzygies. Corollary~\ref{family:outputs} gives the
corresponding symmetric algebras.

Retain $q\in k^\times$ of infinite multiplicative order and the algebra
$H=k\langle x,y\rangle/(x^2,y^2,xy+qyx)$, and put $z=yx$.
For $n\ge2$, let $Q_n$ have vertices
$i\in\mathbb Z/n\mathbb Z$ and arrows $x_i,y_i:i\to i+1$.
Write $e_i$ for the trivial path at vertex $i$.
Products of paths are read with the rightmost arrow traversed first.
Define
\begin{equation*}
 H_n=kQ_n/\bigl(x_{i+1}x_i,\ y_{i+1}y_i,
        \ x_{i+1}y_i+q\,y_{i+1}x_i\mid i\in\mathbb Z/n\mathbb Z\bigr),
 \qquad z_i=y_{i+1}x_i.
\end{equation*}
For $q=1$, the quiver and relations coincide with those of the
circular double-arrow algebra studied over an algebraically closed
field in \cite[Section~1 and Proposition~2.2]{Itaba}. Our
infinite-order assumption excludes this case.

To transfer the extension calculations from $H$ to its cyclic
covering, we give $H$ the $\mathbb Z/n\mathbb Z$-grading determined by
$\deg(x)=\deg(y)=1$ and $\deg(z)=2$.  Then there is an exact equivalence
\begin{equation}\label{family:graded-equivalence}
 H_n\text{-}\mathrm{mod}\simeq
 (\mathbb Z/n\mathbb Z)\text{-}\mathrm{grmod}\,H.
\end{equation}
Here $\mathrm{grmod}\,H$ denotes the category of graded left
$H$-modules with degree-zero morphisms, with the finite-dimensional
convention fixed in Subsection~\ref{prelim:notation}.
A graded module $V=\bigoplus_iV_i$ becomes an $H_n$-module by letting
$e_i$ project onto $V_i$ and $x_i,y_i$ act as the restrictions of
$x,y$ to $V_i$.  Conversely, an $H_n$-module acquires the grading
$V_i=e_iV$ and the $H$-action $x=\sum_i x_i$, $y=\sum_i y_i$.
The relations make these functors mutually inverse, and exactness
is tested on the graded or vertex components. This gives the left-module version
of the graded-module and smash-product equivalence in
\cite[Theorems~1.3 and~2.1]{NR}. The cited results are formulated
for right modules; applying them to $H^{\mathrm{op}}$ gives the
left-module convention used here. The direct construction requires
no restriction on $\operatorname{char}(k)$ relative to $n$.

To compare extension groups, let $U$ denote the forgetful functor and
use the shift convention $W\langle d\rangle_i=W_{i+d}$ for graded
modules $W$. The decomposition by internal degree gives natural
isomorphisms $\Hom_H(UV,UW)\simeq
\bigoplus_{d\in\mathbb Z/n\mathbb Z}\Hom_{\mathrm{grmod}\,H}(V,W\langle d\rangle)$
for graded modules $V,W$. For $f:UV\to UW$, its component indexed
by $d$ is given by
$f_d|_{V_i}=\operatorname{pr}_{W_{i+d}}f|_{V_i}$, where $\operatorname{pr}_{W_{i+d}}$ is the canonical projection onto the summand $W_{i+d}$.
Homogeneity of the action makes each $f_d$ $H$-linear.
Since $U$ sends shifts of $H$ to free $H$-modules and preserves
direct summands, it preserves projectives.
Applying a graded projective resolution therefore gives
\begin{equation}\label{extension-isomorphism}
 \Ext_H^j(UV,UW)\simeq
 \bigoplus_{d\in\mathbb Z/n\mathbb Z}
 \Ext_{\mathrm{grmod}\,H}^j(V,W\langle d\rangle)
 \qquad(j\ge0).
\end{equation}
The Hom differentials preserve internal degree, and cohomology
commutes with the finite direct sum.  In particular, every graded self-extension group
is a direct summand of the corresponding ordinary self-extension
group. We now combine the graded description with syzygy and Nakayama calculations
to verify the seed conditions.

\begin{Prop}\label{family:seed}
The algebra $H_n$ is connected, split basic and Frobenius, with
dimension $4n$ and $n$ simple modules.  For $a\in k^\times$ and
$r\in\mathbb Z/n\mathbb Z$, we define the $H_n$-module $M_n(a,r)$ to have basis
$u,v$, supported at vertices $r,r+1$, respectively, with
$x_ru=v$ and $y_ru=-a^{-1}v$; all other arrow actions are zero.
Then
\begin{equation*}
 \Omega_{H_n}M_n(a,r)\simeq M_n(qa,r+1),\quad
 \nu_{H_n}M_n(a,r)\simeq M_n(q^{-2}a,r-2),\quad
 \tau_{H_n}M_n(a,r)\simeq M_n(a,r).
\end{equation*}
These modules are nonprojective and satisfy
\begin{equation*}
 \begin{gathered}
 \End_{H_n}(M_n(a,r))=\underline{\End}_{H_n}(M_n(a,r))=k,\quad
 \Ext_{H_n}^1(M_n(a,r),M_n(a,r))=k,\\
 \Ext_{H_n}^j(M_n(a,r),M_n(a,r))=0\quad(j\ge2).
 \end{gathered}
\end{equation*}
In particular, $(H_n,L_n)$, where $L_n=M_n(1,0)$, satisfies the
hypotheses of Proposition~\ref{general:seed}.  Its associated algebra
$C_n=\End_{H_n}(H_n\oplus L_n)^{\mathrm{op}}$ is connected and
split basic, with dimension $4n+5$ and $n+1$ simple modules.
\end{Prop}
\begin{proof}
At each starting vertex, the quadratic relations leave the single
length-two basis path $z_i$.  Every path of length three vanishes:
an alternating path reduces to one with two consecutive $x$-arrows
or two consecutive $y$-arrows.  Hence $\{e_i,x_i,y_i,z_i\}_{i\in\mathbb Z/n\mathbb Z}$ is
a basis of $H_n$, the radical of $H_n$ has cube zero, and
$H_n/\rad(H_n)\simeq k^n$.  Since the quiver is connected, $H_n$ is
connected and split basic.  This argument also applies when $n=2$;
in that case $z_i$ is a loop, but remains linearly independent of
$e_i$ because the relations are homogeneous in path length.

To show that $H_n$ is Frobenius, we require an explicit Frobenius form.
Define $\lambda_n:H_n\to k$ by $\lambda_n(z_i)=1$ and let it
vanish on the remaining basis vectors.
For $b=\sum_i(\alpha_ie_i+\beta_ix_i+\gamma_iy_i+\delta_iz_i)\in H_n$
with $\alpha_i, \beta_i, \gamma_i, \delta_i\in k$, one has
$\lambda_n(bz_i)=\alpha_{i+2}$, $\lambda_n(by_i)=-q\beta_{i+1}$,
$\lambda_n(bx_i)=\gamma_{i+1}$ and $\lambda_n(be_i)=\delta_i$.
These identities show that the associative pairing is nondegenerate,
so $H_n$ is Frobenius. Its Nakayama automorphism is given by
$\eta_n(e_i)=e_{i-2}$, $\eta_n(x_i)=-q x_{i-2}$,
$\eta_n(y_i)=-q^{-1}y_{i-2}$ and $\eta_n(z_i)=z_{i-2}$.
These formulas preserve the relations.  For the defining
identity $\lambda_n(ab)=\lambda_n(b\eta_n(a))$, the only nonzero
arrow pairings to verify are
$\lambda_n(x_iy_{i-1})=-q=\lambda_n(y_{i-1}\eta_n(x_i))$ and
$\lambda_n(y_ix_{i-1})=1=\lambda_n(x_{i-1}\eta_n(y_i))$.
The vertex and socle pairings follow from the shift $i\mapsto i-2$;
every other basis pairing vanishes on both sides.

To compute the Auslander--Reiten translation, we now determine the
syzygies and the Nakayama twist. The projective cover
$H_ne_r\to M_n(a,r)$, sending $e_r$ to $u$, has kernel $k(x_r+a y_r)\oplus kz_r$.  The assignment
$u'\mapsto x_r+a y_r$, $v'\mapsto-qa z_r$, where $\{u', v'\}$ is a basis of $M_n(qa,r+1)$, defines an isomorphism from $M_n(qa,r+1)$ onto this kernel.  This proves the syzygy formula.

For the Nakayama functor, we use $\nu_{H_n}V\simeq{}_{\eta_n^{-1}}V$. In particular, $\nu_{H_n}(H_ne_i)\simeq H_ne_{i-2}$. In the twist of $M_n(a,r)$, the top is supported at $r-2$, and the
actions of $x_{r-2},y_{r-2}$ on $u$ are $-q^{-1}v,qa^{-1}v$,
respectively. With $v''=-q^{-1}v$, these actions become $x_{r-2}u=v''$ and $y_{r-2}u=-(q^{-2}a)^{-1}v''$.
Thus the twist is isomorphic to $M_n(q^{-2}a,r-2)$. The formula $\tau\simeq\Omega^2\nu$ from
Subsection~\ref{prelim:stable}, together with the syzygy and twist formulas, now gives $\tau_{H_n}M_n(a,r)\simeq M_n(a,r)$: both the parameter change and the vertex shift cancel.

For the endomorphism and extension calculations, we use
\eqref{family:graded-equivalence} to regard $M_n(a,r)$ as $M(a)$ with
$u_a$ in degree $r$ and $v_a$ in degree $r+1$.
By Lemma~\ref{seed:quantum-pair}\textup{(\ref{seed:pair-modules})},
$\End_H(M(a))=k1\oplus k\varepsilon_a$, with $1$ and $\varepsilon_a$ in degrees zero
and one, respectively. Since $n\ge2$, the degree-zero endomorphisms are precisely the
scalar multiples of the identity. Thus $\End_{H_n}(M_n(a,r))=k$.
If $M_n(a,r)$ were projective, its
underlying module $M(a)$ would be projective, which is impossible.
Its stable identity is therefore nonzero, and its stable
endomorphism algebra is also $k$.

Lemma~\ref{seed:quantum-pair}\textup{(\ref{seed:pair-modules})} gives $\Ext_H^1(M(a),M(a))=k$ and $\Ext_H^j(M(a),M(a))=0$ for $j\ge2$; see also \cite[Section~4.1, Corollary~4.3]{Erd} for these ext-finite modules. The direct-sum decomposition \eqref{extension-isomorphism}
implies $\Ext_{H_n}^j(M_n(a,r),M_n(a,r))=0$ for $j\ge2$.

In degree one, we must check that the ordinary generator of $\Ext_H^1(M(a), M(a))$ in Lemma~\ref{seed:quantum-pair}\textup{(\ref{seed:pair-resolution})} has internal
degree zero, so that it gives an unshifted graded extension.
We use the resolution in Lemma~\ref{seed:quantum-pair}\textup{(\ref{seed:pair-resolution})}
with the grading in which the canonical generator $1$ of the $j$-th free module $P_j(a)$ of rank $1$ has degree $r+j$. Its differential is right multiplication by $x+q^{j-1}ay$ for $j\ge1$, and is homogeneous of degree zero in this grading. The representative of the ordinary first extension sends the generator of $P_1(a)$ (of degree $r+1$) to $v_a$, which also has degree $r+1$. Thus its internal degree is zero, so the entire one-dimensional group belongs to the unshifted graded summand of $\Ext_H^1(M(a),M(a))$. Hence
$\Ext_{H_n}^1(M_n(a,r),M_n(a,r))=k$.

Finally, Frobenius duality gives
$\dim_k(\Hom_{H_n}(H_n,L_n))=\dim_k(\Hom_{H_n}(L_n,H_n))=2$, and $\dim_k(\End_{H_n}(L_n))=1$. Consequently $\dim_k(C_n)=4n+5$; connectedness and the number of simple modules follow from
Proposition~\ref{general:seed}.
\end{proof}

\Needspace{5\baselineskip}
Applying the main theorem to these seeds gives the following family.

\begin{Koro}\label{family:outputs}
For every $n\ge2$, the pair $(H_n,L_n)$ determines, through
Theorem~\ref{general:construction}, a connected finite-dimensional
symmetric algebra $A_n$ and a nonprojective $A_n$-module $T_n$
such that
$\Ext_{A_n}^i(T_n,T_n)=0$ for all $i>0$.
The algebra $A_n$ has $4(n+1)$ simple modules.
\end{Koro}
\begin{proof}
We apply Theorem~\ref{general:construction} to the seed $(H_n,L_n)$ from Proposition~\ref{family:seed} by taking $r=n$. Part~\textup{(\ref{general:construction-basic})} then gives $4(n+1)$ simple modules over $A_n$.
\end{proof}

\section*{Declaration of AI assistance}
Building on the constructions of OpenAI~\cite{O} and Tang~\cite{F},
this work develops a general framework for constructing counterexamples
to Tachikawa's second conjecture. The author directed the research, contributed homological and
categorical ideas, developed new arguments and alternative approaches
to some proofs, and critically reviewed and revised the manuscript.
GPT-6.1 (OpenAI) provided substantial assistance in exploring
proof strategies, drafting and revising mathematical
arguments, checking calculations and notation, and improving the
organization and exposition. The author takes full responsibility for the mathematical content
and the final manuscript.

% Bibliographic convention: for Chinese authors, retain the initials of every given-name syllable.
\begin{thebibliography}{99}
\setlength{\itemsep}{1pt}

\bibitem{Alperin}{{\sc J.~L.~Alperin},
Periodicity in groups, \emph{Illinois J. Math.} \textbf{21} (1977), no.~4, 776--783.
\href{https://doi.org/10.1215/ijm/1256048927}
{doi:10.1215/ijm/1256048927}.}

\bibitem{ARS}{{\sc M.~Auslander, I.~Reiten} and {\sc S.~Smal{\o}},
\emph{Representation theory of Artin algebras}, Cambridge Studies in Advanced Mathematics, vol.~36.
Cambridge University Press, Cambridge, 1995. }

\bibitem{BBD}{{\sc A.~A.~Beilinson, J.~Bernstein} and {\sc P.~Deligne},
Faisceaux pervers, in \emph{Analyse et topologie sur les espaces singuliers, I},
\emph{Ast\'erisque} \textbf{100} (1982), 5--171.
\url{https://www.numdam.org/item/AST_1982__100__1_0/}.}

\bibitem{BoothDQ}{{\sc M.~Booth},
Singularity categories via the derived quotient,
\emph{Adv. Math.} \textbf{381} (2021), 107631.
\href{https://doi.org/10.1016/j.aim.2021.107631}
{doi:10.1016/j.aim.2021.107631}.
\url{https://arxiv.org/abs/2003.05439v3}.}

\bibitem{BCL}{{\sc C.~Braun, J.~Chuang} and {\sc A.~Lazarev},
Derived localisation of algebras and modules,
\emph{Adv. Math.} \textbf{328} (2018), 555--622.
\href{https://doi.org/10.1016/j.aim.2018.02.004}
{doi:10.1016/j.aim.2018.02.004}.
\url{https://arxiv.org/abs/1505.01146v3}.}

\bibitem{CE}{{\sc H.~Cartan} and {\sc S.~Eilenberg},
\emph{Homological Algebra},
Princeton Mathematical Series, vol.~19,
Princeton University Press, Princeton, NJ, 1956.}

\bibitem{CFX}{{\sc H.~X.~Chen, M.~Fang} and {\sc C.~C.~Xi},
Tachikawa's second conjecture, derived recollements,
and gendo-symmetric algebras,
\emph{Compos. Math.} \textbf{160} (2024), 2704--2737.
\href{https://doi.org/10.1112/S0010437X24007395}
{doi:10.1112/S0010437X24007395}.}

\bibitem{CXSO}{{\sc H.~X.~Chen} and {\sc C.~C.~Xi},
Homological theory of self-orthogonal modules,
\emph{Trans. Amer. Math. Soc.} \textbf{378} (2025), 7287--7335.
\href{https://doi.org/10.1090/tran/9426}{doi:10.1090/tran/9426}.}

\bibitem{CXVG}{{\sc H.~X.~Chen} and {\sc C.~C.~Xi},
Virtually Gorenstein algebras of infinite dominant dimension,
\emph{J. Pure Appl. Algebra} \textbf{230} (2026), no.~4, 108224.
\href{https://doi.org/10.1016/j.jpaa.2026.108224}
{doi:10.1016/j.jpaa.2026.108224}.}

\bibitem{Dugas}{{\sc A.~S.~Dugas},
On periodicity in bounded projective resolutions,
preprint (2012), arXiv:1203.2408.
\url{https://arxiv.org/abs/1203.2408}.}

\bibitem{EM}{{\sc H.~Enomoto} and {\sc R.~Marczinzik},
On the homological conjectures for Artin algebras,
\mbox{arXiv:2609.19172v2} (2026).
\url{https://arxiv.org/abs/2609.19172v2}.}

\bibitem{Erd}{{\sc K.~Erdmann},
Ext-finite modules for weakly symmetric algebras with radical cube zero,
\emph{J. Aust. Math. Soc.} \textbf{102} (2017), 108--121.
\href{https://doi.org/10.1017/S1446788716000331}{doi:10.1017/S1446788716000331}.
\url{https://arxiv.org/abs/1511.01418}.}

\bibitem{Itaba}{{\sc A.~Itaba},
On Hochschild cohomology of a self-injective special biserial
algebra obtained by a circular quiver with double arrows,
\emph{Comm. Algebra} \textbf{44} (2016), no.~1, 404--415.
\href{https://doi.org/10.1080/00927872.2014.982815}{doi:10.1080/00927872.2014.982815}.
\url{https://arxiv.org/abs/1404.2032}.}

\bibitem{IS}{{\sc O.~Iyama} and {\sc \O.~Solberg},
Auslander--Gorenstein algebras and precluster tilting,
\emph{Adv. Math.} \textbf{326} (2018), 200--240.
\href{https://doi.org/10.1016/j.aim.2017.11.025}{doi:10.1016/j.aim.2017.11.025}.
\url{https://arxiv.org/abs/1608.04179v2}.}

\bibitem{Jin}{{\sc H.~B.~Jin},
Cohen--Macaulay differential graded modules and negative Calabi--Yau configurations,
\emph{Adv. Math.} \textbf{374} (2020), 107338.
\href{https://doi.org/10.1016/j.aim.2020.107338}
{doi:10.1016/j.aim.2020.107338}.}

\bibitem{KrauseSD}{{\sc H.~Krause},
The stable derived category of a noetherian scheme,
\emph{Compos. Math.} \textbf{141} (2005), no.~5, 1128--1162.
\href{https://doi.org/10.1112/S0010437X05001375}
{doi:10.1112/S0010437X05001375}.}

\bibitem{KrauseHTR}{{\sc H.~Krause},
\emph{Homological Theory of Representations},
Cambridge Studies in Advanced Mathematics, vol.~195,
Cambridge University Press, Cambridge, 2022.
\href{https://doi.org/10.1017/9781108979108}
{doi:10.1017/9781108979108}.}

\bibitem{LiVG}{{\sc W.~Q.~Li},
Virtually Gorenstein Artin algebras are weakly Gorenstein,
arXiv:2608.10049 (2026).}

\bibitem{Mueller}{{\sc B.~J.~M\"uller},
The classification of algebras by dominant dimension,
\emph{Canad. J. Math.} \textbf{20} (1968), 398--409.
\href{https://doi.org/10.4153/CJM-1968-037-9}
{doi:10.4153/CJM-1968-037-9}.}

\bibitem{NR}{{\sc C.~N\u{a}st\u{a}sescu} and {\sc N.~Rodin\`o},
Group graded rings and smash products,
\emph{Rend. Sem. Mat. Univ. Padova} \textbf{74} (1985), 129--137.
\url{https://www.numdam.org/item/RSMUP_1985__74__129_0/}.}

\bibitem{O}{{\sc OpenAI},
A counterexample to Tachikawa's second conjecture,
preprint, September~23, 2026.
\href{https://github.com/openai/math/blob/adc7f1241b42e322a6451854ab7e4b4c146bf78a/preprints/A-counterexample-to-Tachikawas-second-conjecture-September-23-2026/paper.pdf}
{Fixed repository version \texttt{adc7f1241b42}}.}

\bibitem{Rotman}{{\sc J.~J.~Rotman},
\emph{An Introduction to Homological Algebra},
second edition, Universitext,
Springer, New York, 2009.
\href{https://doi.org/10.1007/978-0-387-68324-9}
{doi:10.1007/978-0-387-68324-9}.}

\bibitem{F}{{\sc Q.~Y.~Tang},
A note on counterexamples to homological conjectures,
Hexagon preprint hexagon:2610.00143v4, 2026.
\url{https://hexagonmath.org/2610.00143v4}.}

\bibitem{Wang}{{\sc R.~Wang},
Gorenstein triangular matrix rings and category algebras,
\emph{J. Pure Appl. Algebra} \textbf{220} (2016), 666--682.
\href{https://doi.org/10.1016/j.jpaa.2015.07.011}
{doi:10.1016/j.jpaa.2015.07.011}.}

\bibitem{Xia}{{\sc W.~H.~Xia},
The Auslander--Reiten conjecture for quantum complete intersections,
\mbox{arXiv:2609.24007v1} (2026).
\url{https://arxiv.org/abs/2609.24007v1}.}

\bibitem{ZZ}{{\sc X.~J.~Zhang} and {\sc P.~Y.~Zhou},
Igusa--Todorov algebras and the Auslander--Reiten conjecture,
\mbox{arXiv:2609.14453v2} (2026).
\url{https://arxiv.org/abs/2609.14453v2}.}

\end{thebibliography}

\bigskip
\noindent\textsc{Hongxing Chen}\\
School of Mathematical Sciences \& Academy for Multidisciplinary Studies,\\
Capital Normal University, 100048 Beijing, P.~R.~China\\
Email: \href{mailto:chenhx@cnu.edu.cn}{\texttt{chenhx@cnu.edu.cn}}

\end{document}







